\documentclass[10pt,a4paper]{article}

% ----------------- Paquetes -------------------%
\usepackage[T1]{fontenc}
\usepackage[utf8]{inputenc}
\usepackage[a4paper,margin=2.5cm]{geometry}
\usepackage{mathptmx}             
\usepackage{changepage} 
\usepackage{setspace}
\usepackage[spanish]{babel}
\usepackage[labelfont=bf,labelsep=space]{caption}
\usepackage{amssymb}
\usepackage{amsmath}
\usepackage{ebgaramond}
\usepackage{placeins}
\usepackage{etoolbox}
\usepackage{setspace}
\usepackage{parskip} 
\usepackage{tcolorbox}
\usepackage{needspace}
\usepackage{xparse}
\usepackage{ocgx2}
\usepackage{hyperref}
\usepackage{hypcap}
\usepackage{soul}\sethlcolor{yellow}% resaltado amarillo: \hl{texto} (Panel TCEE, ⌃H)


%%%%%%%%%%%%%%%%%%%%%%%%%%%%%%%%%%%%%%%%%%%%%%%%%%%%%%%%%%%%%%%%
% --- Fija primero tus valores globales "buenos" ---
\setstretch{1.2}                 % Interlineado global
\setlength{\parskip}{0.7\baselineskip}
\setlength{\parindent}{0pt}
\newlength{\GlobalParskip}
\setlength{\GlobalParskip}{\parskip}

\newcounter{eqnum}[subsection]
\renewcommand{\theeqnum}{\arabic{eqnum}}
\newcounter{alphaitem}
\newcounter{numitem}

%% ------------------- Recuadros----------------------%%
\newcommand{\puntosclave}[1]{%
  \begin{tcolorbox}[
    colframe=black,
    title={Puntos clave},
    before upper={\setlength{\parindent}{0.5cm}\setlength{\parskip}{0.5\baselineskip}}
  ]
  #1
  \end{tcolorbox}
}

\newcommand{\modificaciones}[1]{
  \begin{tcolorbox}[
    colback=white,
    colframe=red,
    title={Modificaciones a realizar},
    before upper={\setlength{\parindent}{0.5cm}\setlength{\parskip}{0.5\baselineskip}}
  ]
  #1
  \end{tcolorbox}
}

%% ------------------- Preguntas TEST----------------------%%
% Opciones: radio (single-choice)
\NewDocumentCommand{\OptRadio}{m m m}{%
  \noindent\CheckBox[radio,name=#1,value=#2,width=1.2ex,height=1.2ex]{}%
  \;\textbf{#2.}~#3\par
}

% Opciones: checkbox (multi-choice)
\NewDocumentCommand{\OptCheck}{m m}{%
  \noindent\CheckBox[name=#1,width=1.2ex,height=1.2ex,bordercolor={0 0 0}]{}%
  \;#2\par
}

% Pregunta single-choice (radio)
% Args: 1=id, 2=enunciado, 3=opciones, 4=solución+justificación, 5=imagen (o {}), 6=origen
\NewDocumentCommand{\PreguntaTestSingle}{m m m m m m}{%
  \par\medskip
  \noindent\textbf{#1.}~#2\par\smallskip

    % Imagen (si no es {})
    \IfBlankTF{#5}{}{%
      \begin{center}
        \includegraphics[width=0.75\linewidth]{#5}
      \end{center}
    }

    \begin{Form}
      #3
    \end{Form}

    \par\smallskip
    \switchocg{sol-#1}{\fbox{Mostrar / ocultar solución}}%
    \hfill{\footnotesize #6}\par\smallskip

    \begin{ocg}{Solución}{sol-#1}{0}
      \noindent #4
    \end{ocg}
  \par\medskip
}

% Pregunta multi-choice (checkboxes)
\NewDocumentCommand{\PreguntaTestMulti}{m m m m m m}{%
  \par\medskip
  \noindent\textbf{#1.}~#2\par\smallskip

    % Imagen (si no es {})
    \IfBlankTF{#5}{}{%
      \begin{center}
        \includegraphics[width=0.75\linewidth]{#5}
      \end{center}
    }
    \begin{Form}
      #3
    \end{Form}

    \par\smallskip
    \switchocg{sol-#1}{\fbox{Mostrar / ocultar solución}}%
    \hfill{\footnotesize #6}\par\smallskip

    \begin{ocg}{Solución}{sol-#1}{0}
      \noindent #4
    \end{ocg}
  \par\medskip
}

%% ------------------- CITA ----------------------%%
% \begin{cita}[Autor][Año][Obra] texto \end{cita}: cita sangrada y, debajo a la derecha, «AUTOR (Año), Obra». Los tres datos son opcionales.
% En el Panel TCEE: escribir \cita e Intro (el cursor va al texto; Tab pasa a Autor, Año y Obra).
\NewDocumentEnvironment{cita}{ O{} O{} O{} +b }{%
  \begin{adjustwidth}{2cm}{0pt}
  \raggedright
  #4\par
  \raggedleft
  \if\relax\detokenize{#1#2#3}\relax
    % sin firma
  \else
    #1%
    \if\relax\detokenize{#2}\relax\else\if\relax\detokenize{#1}\relax\else\space\fi(#2)\fi
    \if\relax\detokenize{#3}\relax\else\if\relax\detokenize{#1#2}\relax\else, \fi\textit{#3}\fi
  \fi
  \end{adjustwidth}
}{}



%% ------------------- Listas----------------------%%
    % --- Lista alfabética ---
    \newenvironment{la}
      {\begin{list}{\alph{alphaitem})}{%
          \usecounter{alphaitem}%
          \setlength{\leftmargin}{1cm}%
          \setlength{\parskip}{3pt}% 
          \setlength{\itemsep}{0pt}% ← espacio entre ítems
          \setlength{\parsep}{0pt}%  ← párrafos dentro de un ítem
      }}
      {\end{list}}
      
    % --- Lista numérica ---
    \newenvironment{lnum}
      {\begin{list}{\arabic{numitem}.}{%
          \usecounter{numitem}%
          \setlength{\leftmargin}{1cm}%
          \setlength{\parskip}{3pt}% 
          \setlength{\itemsep}{0pt}% ← espacio entre ítems
          \setlength{\parsep}{0pt}%  ← párrafos dentro de un ítem
      }}
      {\end{list}}
      
% ------------------- Imágenes----------------------%
    \graphicspath{{Img/}}% Rutas por defecto 
    % --- Imagen adaptativa ---
    % Registros de dimensión definidos una sola vez
    \newdimen\imgcap
    \newdimen\imgmin
    \newdimen\imgmax
    \newdimen\imgremain
    \newdimen\imgh
    \newdimen\imgneed
    
    \newcommand{\imagenfit}[3]{%
      \addvspace{10pt}% separación antes
      \par\begingroup
        % Parámetros de composición (asignaciones, no \newdimen)
        \imgcap=1\baselineskip            % alto estimado de título+fuente
        \imgmin=0.15\textheight
        \imgmax=0.15\textheight
        \imgremain=\pagegoal
        \ifdim\pagegoal=\maxdimen \imgremain=0pt\fi
        \advance\imgremain by -\pagetotal
        \advance\imgremain by -\imgcap
        % Decisión de altura destino
        \ifdim\imgremain<\imgmin
          \newpage
          \imgh=0.20\textheight
        \else
          \imgh=\imgremain
          \ifdim\imgh>\imgmax \imgh=\imgmax \fi
        \fi
        % Reservar espacio para evitar cortes del bloque completo
        \imgneed=\imgh \advance\imgneed by \imgcap
        \Needspace{\imgneed}
        % Cuerpo
        \noindent\begin{minipage}{\textwidth}
          \centering
          \captionof{figure}{\textemdash\ #2}\par\vspace{0.3em}% título arriba
          \includegraphics[width=\linewidth,height=\imgh,keepaspectratio]{\detokenize{#1}}\par
          \vspace{0.3em}\small\textit{Fuente: }#3% fuente abajo
        \end{minipage}%
      \endgroup
      \par\addvspace{10pt}%
    }

%------------ Bloque de RECUADRO ESPECIAL -------------%
    \newenvironment{specialblock}[1]{%
      \begin{quote} % crea sangría a izquierda y derecha
      \small % texto más pequeño
      \noindent\textbf{#1}\\[-0.3em] % título en negrita
      \rule{\linewidth}{0.4pt}\\[0.6em]}{
      \end{quote}}
      
%-------------------------- Portada -----------------------%
\newcommand{\oldtitle}[1]{\def\OldTitle{#1}}
\newcommand{\changerationale}[1]{\def\ChangeWhy{#1}}

\newcommand{\changebox}{%
\begin{tcolorbox}[title={Historial de cambios del tema}]
\textbf{Título anterior:} \OldTitle\par
\textbf{Motivo del cambio:} \ChangeWhy
\end{tcolorbox}
}

%-------------------------- Author Font -----------------------%
    \newcommand{\authorfont}[1]{{\fontfamily{EBGaramond-TLF}\selectfont #1}}

%-------------------------- Anexos -----------------------%
    \makeatletter
    \newcounter{anexo}
    \renewcommand{\theanexo}{\Roman{anexo}}
    \newcommand{\anexo}[1]{%
      \clearpage
      \phantomsection
      \refstepcounter{anexo}%
      \noindent\textbf{Anexo \theanexo.- }#1\par\medskip
      \addcontentsline{stc}{section}{Anexo \theanexo.- #1}%
    }
    \makeatother

    %-------------------------- Lista de figuras (personal) -----------------------%
\makeatletter
\newcommand{\MyListOfFigures}{%
  \clearpage
  \phantomsection
  {\centering\bfseries\Large Índice de figuras\par}%
  \medskip
  % Entrada en nuestro índice de contenido (stc)
  \addcontentsline{stc}{section}{Índice de figuras}%
  % Recuperación del listado (sin cabecera propia)
  \begingroup
    \parindent=0pt\parskip=2pt
    \@starttoc{lof}%
  \endgroup
}
\makeatother

%-------------------------- Bibliografía -----------------------%
\makeatletter
% Contador para las referencias
\newcounter{myref}
% Cabecera de la bibliografía, entra en el índice 'stc'
\newcommand{\MyBibliography}{%
  \clearpage
  \phantomsection
  {\centering\bfseries\Large Bibliografía y referencias\par}%
  \medskip
  \addcontentsline{stc}{section}{Bibliografía y referencias}%
  \setcounter{myref}{0}%
}
\newcommand{\MyBibItem}[4]{%
  \refstepcounter{myref}%
  \noindent[\themyref]\ #1.\ \emph{#2}. #3. #4\par
}
\makeatother

%--------- Establecer cuatro niveles de índice --------%
    \setcounter{tocdepth}{4}   
    \setcounter{secnumdepth}{4}% numera hasta \paragraph

%%%%%%%%%%%%%%%%%%%%%%%%%%%%%%%%

\makeatletter

    %--------- Interlineado Global --------%
    \edef\GlobalBaseLineStretch{\baselinestretch}
    
    %--------- Espaciado de seccinones --------%
    \renewcommand\section{\@startsection{section}
      {1}{\z@}
      {2.5ex \@plus 1ex \@minus .2ex} % espacio antes
      {1.5ex \@plus .2ex}             % espacio después
      {\normalfont\Large\bfseries}    % formato del título
    }
    \renewcommand\subsection{\@startsection{subsection}{2}{\z@}
      {3ex \@plus 0.8ex \@minus .2ex}
      {1.5ex \@plus .2ex}
      {\normalfont\large\bfseries}
    }
    \renewcommand\subsubsection{\@startsection{subsubsection}{3}{\z@}
      {3ex \@plus 0.5ex \@minus .2ex}
      {1ex \@plus .2ex}
      {\normalfont\normalsize\bfseries}
    }
    \renewcommand\paragraph{\@startsection{paragraph}{4}{\z@}%
    {-2ex \@plus -0.5ex \@minus -.2ex}% espacio antes
    {0.8ex \@plus .2ex}% espacio después
    {\normalfont\normalsize\itshape}}% ← cursiva en lugar de negrita

    %--------- Ecuaciones --------%   
    % Variables globales para almacenar títulos y notas
    \gdef\leq@current@section{}%
    \gdef\leq@current@subsection{}%
    \gdef\leq@last@printed@section{}%
    \gdef\leq@last@printed@subsection{}%
    
    % Comando para añadir nota (invisible en el cuerpo)
    \newcommand{\eqnote}[1]{%
      \addcontentsline{leq}{leqnote}{#1}%
    }
    
    % Comando para ecuaciones
    \newcommand{\eqblock}[2]{%
      % Verificar si necesitamos imprimir el título de sección
      \ifx\leq@current@section\leq@last@printed@section\else
        \addcontentsline{leq}{leqsection}{\leq@current@section}%
        \global\let\leq@last@printed@section\leq@current@section
        \gdef\leq@last@printed@subsection{}% Reset subsección al cambiar sección
      \fi
      % Verificar si necesitamos imprimir el título de subsección
      \ifx\leq@current@subsection\leq@last@printed@subsection\else
        \ifx\leq@current@subsection\@empty\else
          \addcontentsline{leq}{leqsubsection}{\leq@current@subsection}%
          \global\let\leq@last@printed@subsection\leq@current@subsection
        \fi
      \fi
      % Ecuación
      \refstepcounter{eqnum}%
      \begin{equation*}
        #1\tag{E.\theeqnum}
      \end{equation*}%
      \addcontentsline{leq}{leqt}{[E.\theeqnum] -- #2}%
      \addcontentsline{leq}{leqm}{#1}%
    }
    
    %%% ----- Índice de ecuaciones ----------%%%
    % Tamaño para []
    \newcommand{\leqIndexFont}{\normalsize}
    \newcommand{\leqTitleFont}{\tiny}
    \newcommand{\leqNoteFont}{\tiny}
    % Cabecera de sección 
    \newcommand*\l@leqsection[2]{%
      \addvspace{25pt}% Mayor separación antes de nueva sección
      \noindent\leqIndexFont
      \makebox[\linewidth]{\rule{.38\linewidth}{0.4pt}\ \ \textbf{  \MakeUppercase{#1}}\ \ \rule{.38\linewidth}{0.4pt}}%
      \par\vspace{-10pt}
    }
    % Cabecera de subsección 
    \newcommand*\l@leqsubsection[2]{%
      \par\addvspace{15pt}%
      \noindent\leqIndexFont #1
      \par\vspace{-7pt}
    }
    % Título de ecuación 
    \newcommand*\l@leqt[2]{%
    \par\addvspace{10pt}%
      \par\noindent\leqTitleFont #1\par
    }
    % Miniatura de ecuación
    \newcommand*\l@leqm[2]{%
      \vspace{0.3\baselineskip}%
      \par\scriptsize\noindent\hspace*{1.5em}\(#1\)\par%
    }
    % Nota de ecuación 
    \newcommand*\l@leqnote[2]{%
      \vspace{0.2\baselineskip}%
      \par\noindent\leqNoteFont\hspace*{1.5em}\textbf{Nota.--} #1\par\vspace{0.5\baselineskip}%
    }
    % Generación del índice
    \newcommand{\mylistofequations}{%
      \clearpage
      \phantomsection
      % Título visual de la lista (ajústalo a tu gusto):
      {\centering\bfseries\Large Índice de ecuaciones\par}
      % Entrada en nuestro índice 'stc':
      \addcontentsline{stc}{section}{Índice de ecuaciones}%
      % Llamada a tu lista real (usa el nombre exacto de tu macro):
      \@starttoc{leq}%
    }
    
    %--------- Captura de títulos de sección --------%
    \let\leq@original@section\section
    \let\leq@original@subsection\subsection
    % Flag para saber si estamos en una sección asterisco
    \newif\ifLeqStarSection
    \LeqStarSectionfalse
    
    % Redefinir \section CON CONTROL DE ASTERISCO
    \renewcommand{\section}{%
      \@ifstar{\leq@section@star}{\@dblarg\leq@section@nostar}%
    }
    \def\leq@section@star#1{%
      \global\LeqStarSectiontrue
      \gdef\leq@current@section{#1}%
      \gdef\leq@current@subsection{}%
      \leq@original@section*{#1}%
      \global\LeqStarSectionfalse
    }
    \def\leq@section@nostar[#1]#2{%
      \global\LeqStarSectionfalse
      \gdef\leq@current@section{#2}%
      \gdef\leq@current@subsection{}%
      \leq@original@section[#1]{#2}%
    }
    % Redefinir \subsection
    \renewcommand{\subsection}{%
      \@ifstar{\leq@subsection@star}{\@dblarg\leq@subsection@nostar}%
    }
    \def\leq@subsection@star#1{%
      \global\LeqStarSectiontrue
      \gdef\leq@current@subsection{#1}%
      \leq@original@subsection*{#1}%
      \global\LeqStarSectionfalse
    }
    \def\leq@subsection@nostar[#1]#2{%
      \global\LeqStarSectionfalse
      \gdef\leq@current@subsection{#2}%
      \leq@original@subsection[#1]{#2}%
    }

    % ----- Restauración de valores Globales de estilo ----------%
    % Flags
    \newif\ifInFootnote
    \InFootnotefalse
    \pretocmd{\@footnotetext}{\InFootnotetrue}{}{}
    \apptocmd{\@footnotetext}{\InFootnotefalse}{}{}
    % Profundidad de listas (soporta anidadas)
    \newcount\MyListDepth \MyListDepth=0
    \AtBeginEnvironment{la}{\advance\MyListDepth by 1}
    \AfterEndEnvironment{la}{\advance\MyListDepth by -1}
    \AtBeginEnvironment{lnum}{\advance\MyListDepth by 1}
    \AfterEndEnvironment{lnum}{\advance\MyListDepth by -1}
    \AtBeginEnvironment{itemize}{\advance\MyListDepth by 1}
    \AfterEndEnvironment{itemize}{\advance\MyListDepth by -1}
    \AtBeginEnvironment{enumerate}{\advance\MyListDepth by 1}
    \AfterEndEnvironment{enumerate}{\advance\MyListDepth by -1}
    % Restauración diferida: hazlo al INICIO del próximo párrafo real
    \newif\if@pendingRestore
    \@pendingRestorefalse
    \newcommand{\RequestRestore}{\global\@pendingRestoretrue}
    \everypar{%
      \if@pendingRestore
        \global\@pendingRestorefalse
        \ifInFootnote\else
          \ifnum\MyListDepth=0
            \setlength{\parskip}{\GlobalParskip}%
            \setstretch{\GlobalBaseLineStretch}%
          \fi
        \fi
      \fi
    }
    % Al final de bloques:
    \AfterEndEnvironment{equation}{\RequestRestore}
    \AfterEndEnvironment{equation*}{\RequestRestore}
    \AfterEndEnvironment{la}{\RequestRestore}
    \AfterEndEnvironment{lnum}{\RequestRestore}
    % Al final de macros:
    \apptocmd{\eqblock}{\RequestRestore}{}{}
    \apptocmd{\imagenfit}{\RequestRestore}{}{}
    
\makeatother

% ===================== ToC propio con 4 niveles (único bloque) =====================
\makeatletter

% --- Escritores: añadimos una línea al .stc en cada cabecera numerada ---
% Guardamos las versiones actuales para llamar al comportamiento normal
\let\MyTOC@orig@section\section
\let\MyTOC@orig@subsection\subsection
\let\MyTOC@orig@subsubsection\subsubsection
\let\MyTOC@orig@paragraph\paragraph

% SECTION
\renewcommand{\section}{%
  \@ifstar{\MyTOC@section@star}{\@dblarg\MyTOC@section@nostar}%
}
\def\MyTOC@section@star#1{%
  \MyTOC@orig@section*{#1}% sin entrada al .stc
}
\def\MyTOC@section@nostar[#1]#2{%
  \MyTOC@orig@section[{#1}]{#2}% comportamiento normal (escribe al .toc)
  % Escribimos también nuestra línea al .stc (texto con \numberline)
  \addcontentsline{stc}{section}{\protect\numberline{\thesection}#2}%
}

% SUBSECTION
\renewcommand{\subsection}{%
  \@ifstar{\MyTOC@subsection@star}{\@dblarg\MyTOC@subsection@nostar}%
}
\def\MyTOC@subsection@star#1{%
  \MyTOC@orig@subsection*{#1}%
}
\def\MyTOC@subsection@nostar[#1]#2{%
  \MyTOC@orig@subsection[{#1}]{#2}%
  \addcontentsline{stc}{subsection}{\protect\numberline{\thesubsection}#2}%
}

% SUBSUBSECTION
\renewcommand{\subsubsection}{%
  \@ifstar{\MyTOC@subsubsection@star}{\@dblarg\MyTOC@subsubsection@nostar}%
}
\def\MyTOC@subsubsection@star#1{%
  \MyTOC@orig@subsubsection*{#1}%
}
\def\MyTOC@subsubsection@nostar[#1]#2{%
  \MyTOC@orig@subsubsection[{#1}]{#2}%
  \addcontentsline{stc}{subsubsection}{\protect\numberline{\thesubsubsection}#2}%
}

% PARAGRAPH
\renewcommand{\paragraph}{%
  \@ifstar{\MyTOC@paragraph@star}{\@dblarg\MyTOC@paragraph@nostar}%
}
\def\MyTOC@paragraph@star#1{%
  \MyTOC@orig@paragraph*{#1}%
}
\def\MyTOC@paragraph@nostar[#1]#2{%
  \MyTOC@orig@paragraph[{#1}]{#2}%
  % Si no numeras los \paragraph, cambia \theparagraph por texto plano sin \numberline
  \addcontentsline{stc}{paragraph}{\protect\numberline{\theparagraph}#2}%
}

% --- Impresor del índice propio (lee el .stc) ---
\newcommand{\MyTOC}{%
  \par\bigskip
  {\centering\bfseries\Large Índice de contenido\par}%
  \medskip
  \begingroup
    \parindent=0pt\parskip=2pt
    \@starttoc{stc}%
  \endgroup
}

\makeatother
% ===================== fin ToC propio =====================


%%%%%%%%%%%%%%


% --- Datos del tema ---
\title{Unión Bancaria\\
\large Pilares y Código normativo único. La Únion de Mercados de Capitales}
\author{Marcelo Ortega}
\date{Marzo 2026}
\newcommand{\TemaCode}{3B44}

\oldtitle{
    B.46. El Sistema Europeo de Bancos Centrales y la política monetaria de la Eurozona: objetivos e instrumentos. La gobernanza económica del euro. La Unión Bancaria y otros desarrollos
}
\changerationale{
    Este tema versa sobre la regulación financiera en la UE. Se hace notar que, mientras que el tema B.27 tiene un enfoque de ciencia económica (teoría y evidencia empírica), este tiene una aproximación más jurídica e institucional.
}

\begin{document}


\maketitle
\changebox

\newpage

% --- Página de resumen y modificaciones ---
\puntosclave{ %% Escribir puntos clave Apuntes CECO
     Este título se desliga de los actuales 3B43 y 3B45. Sobre la relación con los dos primeros: el opositor ha de ser capaz de comprender cómo los tres están profundamente interrelacionados, en tanto que abordan las tres principales áreas de actuación conjunta en política macroeconómica de la UE (política monetaria, coordinación fiscal y construcción de un sistema financiero común y estable).

     Se recomienda dedicar 17-18 minutos a la Unión Bancaria y I O minutos a la Unión de los Mercados de Capitales, aproximadamente.
    }
\vspace{1cm}

\modificaciones{
    \textcolor{red}{\textbf{OJO ->}} Revisar:
    \begin{la}
        \item \href{https://www.bde.es/f/webbe/GAP/Secciones/Publicaciones/InformesBoletinesRevistas/RevistaEstabilidadFinanciera/26/4_REF50_RESOLUCION.pdf}{\textit{DIEZ AÑOS DE RESOLUCIÓN BANCARIA EN EUROPA: EVOLUCIÓN, LECCIONES APRENDIDAS Y RETOS FUTUROS}. CONCEPCIÓN y KRUSE, BdE (2025)}
        \item \href{https://www.bruegel.org/system/files/wp_attachments/pc_2013_05.pdf}{\textit{FROM SUPERVISION TO RESOLUTION: NEXT STEPSON THE ROAD TO EUROPEAN BANKING UNION}. Bruegel (2013)}
        \item \href{https://www.bruegel.org/system/files/wp_attachments/pc_2016_03.pdf}{\textit{SHOULD THE ‘OUTS’ JOIN THE EUROPEAN BANKING UNION?} Bruegel (2016)}
        \item \href{https://www.elsevier.es/en-revista-the-spanish-review-financial-economics-332-pdf-S2173126815000042}{\textit{A banking union for Europe: Making a virtue out of necessity}. ABASCAL et al., 2015}
        \item \href{https://shs.cairn.info/revue-l-europe-en-formation-2017-2-page-61?lang=en}{\textit{A brief overview of the European Banking Union}. GORTSOS (2017)}
        \item \href{https://www.google.es/url?sa=t&source=web&rct=j&opi=89978449&url=https://turia.uv.es/index.php/RDSFin/es/article/view/32884&ved=2ahUKEwiTzNexpf2VAxWnVaQEHUIgEE4QFnoECB4QAQ&usg=AOvVaw3fSPZWkT-Ugw0eY2muzjkQ}{\textit{LA UNIÓN BANCARIA Y SU ANSIADO TERCER PILAR: EL SISTEMA DE GARANTÍA DE DEPÓSITOS EUROPEO (Y EL FONDO ÚNICO DE RESOLUCIÓN COMO BROCHE DE LOS DOS PRIMEROS PILARES)}. VILLARRUBIA (2021)}
    \end{la}
    }

\newpage
\MyTOC
\newpage

\mylistofequations
\clearpage

\MyListOfFigures
\clearpage


\pagenumbering{arabic}
\setcounter{page}{1}


\section{Introducción}



\subsection{Contextualización}


\subsubsection{Relevancia}




\subsubsection{Problemática}



\subsection{Estructura de la exposición}







\clearpage
\section{Política Bancaria de la Unión: Código Normativo Único}
\puntosclave{
    La Unión Bancaria está más desarrollada y su impacto en términos de estabilidad financiera ha sido probablemente mayor.

    El Código Nonnativo Único pretende conseguir un sector bancario más resistente, más transparente y más eficiente y opera en dos grandes áreas: (i) requisitos prudenciales y (ii) prevención, gestión de las quiebras bancarias y protección de los depositantes. 
    
    Los requisitos prudenciales están recogidos en la CRR y en la CRD. La CRR recoge los requisitos mínimos de capital para el riesgo de crédito, el riesgo de mercado y el riesgo operacional, la Ratio de Cobertura de Liquidez (LCR, por sus siglas en inglés), la Ratio de Financiación Neta Estable (NSFR, por sus siglas en inglés), la ratio de apalancamiento, los requisitos de transparencia, de grandes exposiciones y de titulizaciones. La CRD recoge los requisitos gobernanza y control interno, incluyendo el \textit{fit and proper}, los requisitos de Pilar 2 y la autorización. 
    
    La prevención, gestión de las quiebras bancarias y protección de los depositantes se apoya nonnativamente en lo que se conoce como el marco de gestión de crisis y seguro de depósitos (CMDI), compuesto por tres normas: Directiva de Recuperación y Resolución Bancaria (BRRD), Reglamento del Mecanismo Único de Resolución (RMUR) y Directiva de Fondos de Garantía de Depósitos (DFGD), que se aplican en las diferentes etapas del ciclo de vida de los bancos en dificultades. El CMDI ha instaurado importantes principios, como la sustitución del "bail-out". o rescate de la entidad con cargo ª. los fondos del contribuyente, por el "bail-in", por el que el coste de la resolución de los bancos se cubre en primer lugar a través de los propios recursos del banco
}   


La política bancaria de la Unión Europea, en el sentido regulatorio que aquí interesa, es ante todo un ejercicio de armonización. 

Antes de la crisis financiera global, cada Estado miembro transponía a su manera la normativa prudencial internacional, con márgenes de discrecionalidad tan amplios que dos bancos con el mismo perfil de riesgo podían estar sujetos a exigencias de capital radicalmente distintas según el país en que operasen. 

Esa divergencia no era un simple defecto técnico: permitía el arbitraje regulatorio, distorsionaba la competencia dentro del mercado único y, como reveló la crisis, impedía cualquier respuesta coordinada cuando el riesgo se propagaba de un Estado a otro. El Código Normativo Único es la respuesta a ese diagnóstico, y conviene que lo entiendas ante todo como eso: un proyecto de armonización sustantiva (un mismo riesgo, una misma norma, una misma supervisión) antes que como una técnica jurídica concreta. La forma que adopta esa armonización, reglamento o directiva, es instrumental y varía según cuánto margen de adaptación nacional tolere el legislador en cada materia, pero lo que define al Código Normativo Único no es esa forma sino el grado de uniformidad que persigue: por eso se aplica a los veintisiete sin excepción, y por eso, como verás en el subepígrafe siguiente, la existencia de una arquitectura institucional distinta para la Unión Bancaria no rompe esa armonización sustantiva, sino que se superpone a ella únicamente en el plano de quién la hace cumplir.

\subsection{Historia: ¿Cómo hemos llegado hasta aquí?}
La relación entre los bancos y los soberanos europeos se remonta muy atrás en el tiempo. La mayoría de los grandes bancos del continente hunden sus orígenes en el siglo XIX y en los primeros años del XX, un período de nacionalismo económico en el que buena parte de ellos nacieron ligados a proyectos de desarrollo nacional o colonial y, con frecuencia, a un patrocinio estatal directo o indirecto. A cambio de ese patrocinio, los bancos aceptaron como algo natural financiar al Estado y dirigir el crédito hacia los sectores que el gobierno de turno quería favorecer, en un patrón de proximidad entre banca y poder político que la literatura económica bautizó como represión financiera, y que explica por qué, a diferencia de Estados Unidos, donde la banca creció como fenómeno privado y ascendente, en Europa la autoridad supervisora se concibió durante décadas como un asunto de soberanía nacional intocable.

Esa concepción resultó extraordinariamente resistente a la integración europea. El propio Tratado de Roma de 1957 exigía unanimidad para cualquier medida relativa a la protección del ahorro y al ejercicio de la profesión bancaria, un obstáculo que no existía para casi ningún otro sector de servicios. Las sucesivas iniciativas de la Comisión Europea —una primera tentativa en 1965, la Directiva de libertad de establecimiento de 1973, la primera Directiva bancaria de 1977— avanzaron poco, en buena medida porque la propia industria bancaria europea, a diferencia de la manufacturera, se mostró reacia a cualquier consolidación transfronteriza real y prefirió alianzas flexibles entre bancos nacionales que preservaran su identidad de origen. En paralelo, se sucedieron varios comités de coordinación supervisora —el Comité de Gobernadores de 1964, el Groupe de Contact de 1972, el Comité Consultivo Bancario de 1978, el Subcomité de Supervisión Bancaria creado por el propio Comité de Gobernadores entre 1988 y 1990— que respondían todos a un mismo patrón: despertaban en privado la esperanza de desembocar algún día en una supervisión europea integrada, pero esa esperanza se disipaba tras las primeras reuniones y el principio de que la supervisión correspondía en exclusiva al nivel nacional salía siempre reforzado.

\paragraph{La cláusula habilitante de Maastricht y una advertencia prematura.} 
El verdadero punto de inflexión jurídico, aunque tardaría veinte años en activarse, llegó con el Tratado de Maastricht de 1992. Sobre la base del Informe Delors de 1989, el Tratado introdujo, en su artículo 105(6) —hoy artículo 127(6) del TFUE tras la renumeración de Lisboa—, una cláusula habilitante que permitía atribuir al futuro Banco Central Europeo funciones de supervisión prudencial sin necesidad de una nueva reforma de los Tratados. En el momento de su redacción, esta cláusula no tenía vocación de aplicarse de inmediato: fue, más bien, una puerta trasera dejada entornada por precaución. Tommaso Padoa-Schioppa, miembro del Comité Delors y más tarde del primer Comité Ejecutivo del BCE, fue de los pocos que advirtió tempranamente del riesgo de crear una moneda única sin una supervisión bancaria igualmente única, y en 2004 acuñó la expresión que da nombre a buena parte de este epígrafe, single rulebook, para reclamar un marco regulatorio racionalizado y uniforme en toda la Unión. Sus advertencias, sin embargo, no encontraron eco político real durante casi dos décadas: como resumió el historiador financiero Harold James, los primeros diez años de vida del euro estuvieron marcados por la idea, después desmentida, de que la disciplina de mercado por sí sola bastaría para mantener a raya los desequilibrios bancarios nacionales.

\paragraph{La crisis financiera y el círculo vicioso banco-soberano.} 
La quiebra de Lehman Brothers en septiembre de 2008, precedida un año antes por el rescate del banco alemán IKB, desencadenó la crisis financiera global y, a partir de 2009-2010, su metamorfosis en crisis de deuda soberana de la zona euro. El mecanismo que convirtió una crisis bancaria en una crisis existencial para la moneda única fue el ya célebre círculo vicioso entre bancos y soberanos: los bancos, cargados de deuda pública de su propio Estado, se debilitaban a medida que esa deuda perdía valor, mientras que los gobiernos, obligados a rescatar a sus bancos, veían deteriorarse aún más sus propias finanzas públicas, en una espiral que se hizo especialmente visible en Irlanda en 2010, en España durante 2011-2012, y que llegó a poner en duda brevemente la solvencia bancaria francesa en el verano de 2011. El Informe de Larosière, publicado en febrero de 2009 por encargo de la Comisión Europea, diagnosticó con crudeza que la supervisión previa a la crisis había demostrado ser manifiestamente inadecuada, y su Recomendación 10 —que la futura legislación financiera europea se basara, siempre que fuera posible, en reglamentos de aplicación directa antes que en directivas de transposición dispar— es la semilla normativa directa de lo que hoy conocemos como Código Normativo Único. Sobre esa base, el Consejo Europeo creó en 2010 el Sistema Europeo de Supervisión Financiera, con sus tres nuevas Autoridades Europeas de Supervisión —entre ellas la Autoridad Bancaria Europea— encargadas precisamente de redactar ese cuerpo normativo común para los veintisiete.

\paragraph{El verano de 2012: de la crisis a la Unión Bancaria.} 
La expresión "unión bancaria" como tal fue introducida en el debate público por Nicolas Véron en diciembre de 2011, y se popularizó a partir de un artículo del Financial Times de abril de 2012. El propio Herman Van Rompuy, entonces presidente del Consejo Europeo, incorporó la idea a la hoja de ruta hacia una "auténtica Unión Económica y Monetaria". Pero fue la cumbre del Consejo Europeo de los días 28 y 29 de junio de 2012 la que convirtió el proyecto en decisión política irreversible, activando por fin, tras veinte años dormida, la cláusula habilitante de Maastricht. Los historiadores de la integración europea Nielsen y Smeets resumieron el episodio con una frase que conviene recordar: "no había nada inevitable en la unión bancaria"; hicieron falta, simultáneamente, una emergencia económica que ya no admitía demora —la crisis bancaria española estaba entonces en su momento más agudo—, una ventana de oportunidad electoral en Francia y Grecia, el reconocimiento generalizado del círculo vicioso banco-soberano como diagnóstico común, y un ejercicio de liderazgo político en el que Francia renunció a la soberanía nacional sobre su propia supervisión bancaria y Alemania aceptó que también sus bancos públicos y cooperativos quedaran bajo supervisión europea. El propio expresidente francés François Hollande recordaría después que aquella noche del 28 al 29 de junio de 2012 "salvó la moneda única, derrotó a los especuladores de los mercados" y allanó el camino para que, un mes más tarde, Mario Draghi pronunciara su célebre promesa de hacer "todo lo que fuera necesario" para preservar el euro. El Mecanismo Único de Supervisión echó a andar en noviembre de 2014, y el Mecanismo Único de Resolución un año después, completando en dos pilares operativos lo que aquella cumbre de verano había hecho posible.

\paragraph{La construcción paralela del Código Normativo Único.} 
Mientras se fraguaba la arquitectura institucional de la Unión Bancaria, avanzaba en paralelo la construcción de su contenido normativo. Sobre la base del proceso Lamfalussy, ya existente desde comienzos de siglo para el sector de valores, y del impulso decisivo del Informe Larosière, el legislador europeo revisó en 2011 la entonces vigente Directiva de Requisitos de Capital, escindiéndola deliberadamente en dos instrumentos de naturaleza distinta: un Reglamento de aplicación directa y una Directiva de transposición nacional, exactamente la misma técnica de "armonización máxima donde sea posible" que Larosière había recomendado. De ese proceso nació en 2013 el primer Paquete Bancario, CRR/CRD IV, que transpuso a derecho de la Unión el Acuerdo de Basilea III; le siguió en 2019 un segundo paquete, CRR II/CRD V; y más recientemente un tercero, CRR III/CRD VI, cuya negociación se cerró bajo Presidencia española del Consejo y que hoy, como ya sabes, es derecho plenamente vigente. Esta doble construcción simultánea —una Unión Bancaria institucional, nacida de una emergencia política concreta en el verano de 2012, y un Código Normativo Único de alcance más amplio, nacido de un diagnóstico técnico articulado ya desde 2009— es la que explica la arquitectura dual que desarrollamos a continuación.

\subsection{Contenido sustantivo: ¿Qué es y qué objetivos persigue?}
Antes de introducir el contenido material del Single Rulebook, debemos buscar una definición adecuada del concepto. 

Es cierto que el Single Rulebook no constituye per se un término jurídico. Se trata más bien de un concepto político y es poco probable que, por ejemplo, el TJUE lo reconociera como un concepto jurídico autónomo. En efecto, resulta relativamente difícil desarrollar una doctrina jurídica general y absoluta sobre lo que comprende el enfoque regulatorio del Single Rulebook. 

No obstante, sigue siendo posible examinar y explicar el funcionamiento de cada uno de los elementos que integran el Single Rulebook, incluida la forma en que interactúan entre sí. Una vez identificados estos mecanismos, pueden aplicarse de manera útil a todas las ramas del Single Rulebook, y no solo en el contexto de la supervisión de las entidades de crédito.

Existen, potencialmente, al menos dos perspectivas diferentes desde las que examinar el significado exacto del Single Rulebook.

\subsubsection{Aproximación literal: ¿Qué dice?}
Examinemos primero el nombre del concepto utilizado por los responsables políticos de la UE. Si el Single Rulebook es realmente único, ello implica que corresponde a un conjunto único de normas, posiblemente elaboradas en un único lugar. Si realmente es un código normativo (rulebook), debe abarcar un conjunto de normas precisas que se apliquen directamente a las entidades públicas y a los participantes del sector privado, regulando íntegramente sus actividades.

El primer elemento de esa unicidad procede, efectivamente, de la centralización de la producción normativa. La elaboración de normas ha pasado de 28 legisladores individuales a las instituciones y organismos de la Unión. Así, los actos jurídicos que forman parte del Single Rulebook son adoptados por el Consejo y el Parlamento, la Comisión y las ESAs. Desde esta perspectiva, las medidas nacionales adoptadas para aplicar directivas o directrices, o incluso reglamentos, no parecen, a primera vista, cumplir el criterio de unicidad. Sin embargo, como se explicará más adelante, estas medidas podrían y deberían considerarse parte del Single Rulebook.

Considerado como un código normativo (un cuerpo de normas jurídicas), comprende el Derecho sustantivo de la Unión en los ámbitos designados. Ello incluye todas las formas de actos jurídicos de la Unión previstas en el art. $288$ TFUE. Sin embargo, el análisis textual del término no responde a la cuestión de si los actos cuasilegislativos adoptados por organismos de la Unión (como las directrices y recomendaciones de las ESAs) forman parte del Single Rulebook.

Así pues, desde un punto de vista puramente textual, una definición útil del Single Rulebook sería considerarlo como un \textbf{cuerpo único de normas sustantivas de la Unión, adoptadas y aplicables a escala paneuropea}.

\subsubsection{Aproximación teleológica: ¿Qué finalidad persigue?}
Las referencias al Single Rulebook contenidas en la legislación de la Unión adoptada recientemente, así como los tabajos preparatorios y algunas declaraciones políticas, pueden ayudar a identificar los distintos elementos asociados al concepto. 

La lectura de los preámbulos de las diferentes normas de la Unión resulta útil para determinar la intención del legislador de la Unión respecto al significado del Single Rulebook. El tema recurrente es la armonización y el fortalecimiento del mercado único, pero también la garantía de unas condiciones de competencia equitativas entre las entidades financieras (\textit{level playing field}). Además, puede sostenerse que se ha puesto más énfasis en el elemento \textit{rulebook} que en el elemento \textit{single}. 

Así, el Rulebook puede entenderse como un \textbf{concepto general} utilizado por el legislador \textbf{para designar la regulación de las actividades de todos los participantes} en el sector de los servicios financieros.


\paragraph{¿No es el Single Rulebook simplemente una armonización máxima en todos los ámbitos?}
Con mucha frecuencia, el Single Rulebook se entiende simplemente como una armonización máxima en el sector de los servicios financieros, un concepto que suele asociarse con el traslado completo de la producción normativa al nivel de la Unión, eliminando de hecho todas las competencias de aplicación de los Estados miembros. Sin embargo, tal conclusión sería bastante imprecisa. En realidad, el concepto de armonización máxima es diferente.

La doctrina académica no utiliza de manera uniforme el concepto de armonización máxima, sino que emplea distintos términos. Las expresiones \textit{armonización exhaustiva}, \textit{armonización plena}, \textit{armonización total} y \textit{armonización completa} parecen hacer referencia a conceptos similares al de armonización máxima. 

Los colegisladores han empleado de forma constante el término \textit{armonización máxima} en el contexto de la elaboración del Single Rulebook, y una característica recurrente en todos los análisis sobre éste concepto es que se trata de un enfoque legislativo que implica el agotamiento completo del margen de discrecionalidad de un Estado miembro en la aplicación de una determinada norma de la Unión. Sin embargo, \textbf{sería incorrecto equiparar el concepto de Single Rulebook con el de armonización máxima}. El Single Rulebook constituye un enfoque regulatorio, pero no implica siempre una armonización máxima en todos los ámbitos, sino más bien un reequilibrio de los elementos de la construcción vertical de la aplicación del Derecho de la Unión (es decir, de la interacción entre la Unión y los Estados miembros). Esta distinción es importante no solo desde una perspectiva terminológica: el TJUE ha sostenido de manera constante que, en los ámbitos sujetos a armonización máxima, los Estados miembros no pueden mantener medidas más estrictas.

Este nuevo enfoque regulatorio no excluye por completo a los Estados miembros de su papel principal en la aplicación del Derecho de la Unión. Por ello, probablemente \textbf{resulte más adecuado vincular el Single Rulebook} no con la armonización máxima, sino con otro término utilizado en los Tratados: la \textbf{aproximación de las legislaciones aplicables}. A partir de ahí, el grado de aproximación que comporta el enfoque del Single Rulebook varía en función de sus distintos elementos.

Una razón adicional por la que el uso del término armonización máxima puede resultar engañoso en el contexto del concepto de Single Rulebook es que dicho término solo ha sido analizado por el TJUE en relación con las directivas de la Unión. Es como si el TJUE hubiera partido de la premisa de que los reglamentos, al ser directamente aplicables, apenas dejan margen para la actuación de los Estados miembros. Sin embargo, también puede ocurrir que los reglamentos dejen determinadas disposiciones a la aplicación por parte de los Estados miembros.

\subsection{Contenido material: ¿Qué forma el Código Normativo Único?}
El Single Rulebook podría considerarse simplemente como una colección de actos jurídicos individuales, cada uno con sus características particulares. Sin embargo, su estructura es, en cierta medida, \textbf{jerárquica}. 

Si se intenta definir el Single Rulebook atendiendo a su estructura y a las características de los distintos actos jurídicos que lo integran, la definición más adecuada es la siguiente: una \textbf{estructura regulatoria multinivel, compuesta tanto por actos de la Unión como por normas nacionales, que se aplican conjuntamente entre sí}.

Aquí tienes que ser cuidadosa porque "único" es engañoso:

CRR (Reglamento, no CRD) — directamente aplicable, idéntico letra por letra en los 27. Esta es la única parte genuinamente "única".
CRD IV/V, BRRD, DGSD — son Directivas: fijan un resultado armonizado pero se transponen a derecho nacional, y dejan opciones y discrecionalidades nacionales (las llamadas ONDs). Esto significa que, en sentido estricto, el "Código Normativo Único" no es único en su totalidad — solo la porción-Reglamento lo es sin fisuras. El BCE, dentro del SSM, ha intentado cerrar esas discrecionalidades mediante su Guía de opciones y discrecionalidades, pero eso solo vincula a las entidades significativas de la UB, no a los 27.

Los objetivos del Código Normativo Único son claros: \textbf{conseguir un bancario europeo más resistente, más transparente y más eficiente}. Con este objetivo en mente, su razón de ser es que la normativa se aplique homogéneamente, siempre que sea posible.

\subsubsection{Nivel 1: Legislación primaria}
En el nivel superior se encuentran los reglamentos y directivas, que contienen los principios generales y normas detalladas. La inmensa mayoría de instrumentos en este nivel se adoptan vía procedimiento legislativo ordinario (art. $289$ y $294$ TFUE) con base jurídica en el art. $114$ TFUE (armonización para el funcionamiento del mercado interior): codecisión plena Parlamento-Consejo, mayoría cualificada en Consejo, iniciativa exclusiva de la Comisión. No obstante, el \href{https://www.boe.es/buscar/doc.php?id=DOUE-L-2013-82128}{Reglamento (UE) nº 1024/2013 del Consejo, de 15 de octubre de 2013, que encomienda al BCE tareas específicas respecto de políticas relacionadas con la supervisión prudencial de las entidades de crédito} (de ahora en adelante, \textbf{Reglamento MUS}), rompe este patrón, aunque será tratado más adelante.\footnote{%
    Su base jurídica es el art. $127(6)$ TFUE, que exige unanimidad en el Consejo y solo consulta (no codecisión) al Parlamento Europeo. Esto no es un detalle menor — significa que la atribución de funciones de supervisión al BCE se hizo por una vía jurídica deliberadamente distinta y más restrictiva que el resto del CNU, precisamente porque el art. $127(6)$ permite atribuir tareas específicas de supervisión al BCE sin reformar los Tratados, evitando así abrir un proceso de reforma tratatista que en 2012-2013 era políticamente inviable.
}

Este nivel se organiza en dos bloques materialmente distintos según el momento del ciclo de vida de la entidad al que se dirigen: un bloque prudencial, que regula a la entidad en funcionamiento normal (lo que la literatura anglosajona llama \textit{going concern}), y un bloque preventivo o de gestión de crisis, que regula qué ocurre cuando la entidad \textbf{deja de ser viable} (\textit{gone concern}). Distinción que está en el propio diseño legislativo de la Unión, separando deliberadamente los instrumento que fijan cuánto capital y liquidez debe mantener un banco sano de los instrumentos que decides qué hacer con un banco que ya ha fracasado o está a punto de hacerlo.

\textbf{NOTA.-} Antes de describir los instrumentos de primer nivel, una precisión que no se puede perder de vista. Excluimos conscientemente el \href{https://www.boe.es/buscar/doc.php?id=DOUE-L-2014-81709}{Reglamento (UE) nº 806/2014 del Parlamento Europeo y del Consejo, de 15 de julio de 2014, por el que se establecen normas uniformes y un procedimiento uniforme para la resolución de entidades de crédito y de determinadas empresas de servicios de inversión en el marco de un Mecanismo Único de Resolución y un Fondo Único de Resolución y se modifica el Reglamento (UE) nº 1093/2010} (de ahora en adelante \textbf{Reglamento MUR}) pues lo trataremos más adelante.\footnote{%
    Se trata del mecanismo institucional de resolución centralizada del primer nivel porque ese mecanismo solo vincula a los Estados de la Unión Bancaria, pero no debemos pensar que la resolución bancaria como tal es un asunto exclusivo de la Unión Bancaria. La directiva que regula el contenido sustantivo de la resolución (qué es un instrumento de recapitalización interna, qué es el principio de que ningún acreedor reciba menos de lo que habría recibido en concurso, qué autoridad nacional debe designarse como autoridad de resolución) se aplica a los veintisiete, se haya adherido o no el Estado a la Unión Bancaria. Lo que cambia para los Estados de la Unión Bancaria no es la norma sustantiva sino quién toma la decisión última: para la Unión Bancaria decide una junta europea con sede en Bruselas; para el resto, su propia autoridad nacional designada conforme a esa misma directiva.
}

\paragraph{Regulación prudencial: Reglamento (UE) 575/2013 y Directiva 2013/36}
El bloque prudencial está formado por dos instrumentos que nacieron unidos en 2013 como un único paquete legislativo, pero que adoptan formas jurídicas distintas por una razón deliberada. 

El primero es el \href{https://www.boe.es/buscar/doc.php?id=DOUE-L-2013-81261}{Reglamento (UE) n ° 575/2013 del Parlamento Europeo y del Consejo, de 26 de junio de 2013 , sobre los requisitos prudenciales de las entidades de crédito y las empresas de inversión} (de ahora en adelante, \textbf{Reglamento de Requisitos de Capital}), de aplicación directa e idéntica en los veintisiete sin necesidad de transposición, porque el legislador quiso que el núcleo duro del cálculo de solvencia (fondos propios, ratios de capital, tratamiento del riesgo de crédito, de mercado y operacional, ratio de apalancamiento, requisitos de liquidez) no dejara ningún margen de interpretación nacional que pudiera generar arbitraje regulatorio entre bancos de distintos Estados compitiendo en el mismo mercado interior. 

El segundo es la \href{https://www.boe.es/buscar/doc.php?id=DOUE-L-2013-81262}{Directiva 2013/36/UE del Parlamento Europeo y del Consejo, de 26 de junio de 2013, relativa al acceso a la actividad de las entidades de crédito y a la supervisión prudencial de las entidades de crédito y las empresas de inversión} (de ahora en adelante, \textbf{Directiva de Requisitos de Capital}), que regula todo aquello que el legislador consideró que requería sensibilidad a las estructuras institucionales y de supervisión nacionales: la autorización y el \textbf{pasaporte único de entidades de crédito}, los requisitos de gobierno corporativo y de idoneidad de administradores y responsables de funciones clave, el proceso de revisión y evaluación supervisora que cada autoridad competente realiza sobre cada entidad, y los colchones de capital de naturaleza macroprudencial.\footnote{%
    El pasaporte único (o pasaporte comunitario) es un mecanismo legal de la Unión Europea que permite a una entidad de crédito autorizada en un Estado miembro operar en el resto de la UE y del Espacio Económico Europeo. Esto se puede hacer mediante la creación de sucursales o bajo el régimen de libre prestación de servicios, sin requerir nuevas autorizaciones de los reguladores nacionales de cada país de acogida
} Actualizada por CRD V (2019/878) y ahora CRD VI (Directiva 2024/1619): exigió transposición con fecha límite el 11 de enero de 2026, refuerza sustancialmente la evaluación de idoneidad de los miembros del órgano de dirección y titulares de funciones clave, y la ABE desarrollará normas técnicas para armonizar estas evaluaciones y reducir la divergencia actual entre supervisores. Las disposiciones sobre sucursales de terceros países se retrasan hasta el 11 de enero de 2027, con un régimen transitorio para contratos anteriores al 11 de julio de 2026. 

El fundamento último de este reparto, tal como aparece en los considerandos de ambos textos, es alcanzar lo que el legislador llama expresamente el mismo riesgo, la misma norma, la misma supervisión, como respuesta directa al diagnóstico de fragmentación regulatoria que hizo el Informe de LAROSIÈRE en 2009 tras la crisis financiera global: antes de este paquete, cada Estado miembro transponía Basilea II con su propio grado de rigor, y eso permitió a ciertas entidades desplazar riesgo hacia las jurisdicciones más permisivas. Atendiendo, en este sentido, a los principios de armonización máxima, en el CRR, y armonización de mínimos/resultado en la CRD.

En cuanto a sujetos, el bloque prudencial se dirige primariamente a las entidades de crédito, entendidas como aquellas empresas cuya actividad consiste en recibir depósitos u otros fondos reembolsables del público y conceder créditos por cuenta propia, y se extiende por vía de consolidación a las sociedades matrices, las sociedades financieras de cartera y las sociedades financieras mixtas de cartera que encabezan un grupo bancario, aunque estas últimas no sean entidades de crédito en sí mismas.\footnote{
Las empresas de servicios de inversión, que hasta hace pocos años quedaban también dentro del ámbito de este reglamento y esta directiva, han migrado en su mayor parte hacia un régimen prudencial propio y paralelo pensado específicamente para ellas, de modo que hoy el solapamiento entre ambos regímenes es mucho menor que en la década pasada. 
}

En cuanto al procedimiento de supervisión que este bloque articula, la pieza central es el proceso de revisión y evaluación supervisora, mediante el cual la autoridad competente (el BCE si la entidad es significativa y pertenece a la Unión Bancaria, o la autoridad nacional en cualquier otro caso) evalúa anualmente los riesgos de cada entidad y puede imponerle, además del mínimo legal, un requerimiento adicional de capital propio de esa entidad, lo que se conoce como exigencia de Pilar dos, distinta del mínimo uniforme de Pilar uno que fija el reglamento.

\begin{specialblock}{Marco actual: ¿Qué obligaciones concretas impone la normativa sobre la regulación prudencial?}
    La regulación prudencial bancaria tiene por objeto asegurar que las entidades operan con recursos propios suficientes para poder asumir los riesgos que se derivan de su actividad financiera, contribuyendo así a la estabilidad del sistema financiero. 

    A pesar de que reflejar todo el nuevo contenido normativo en un Reglamento habría sido mucho más conveniente desde el punto de vista de aplicación homogénea de la normativa, no es una opción legalmente factible. En efecto, desde un punto de vista jurídico, no todas las normas bancarias pueden adoptar la forma de un Reglamento. Según el artículo 53 del TFUE, las normas que regulan el acceso a una actividad o su ejercicio en el mercado interior deben adoptarse mediante Directivas.
    
    • los requisitos mínimos de capital para el riesgo de crédito, el riesgo de mercado y el riesgo operacional, el riesgo de contraparte y el riesgo de CVA (riesgo de ajuste de valoración del crédito): ratio de capital resulta de dividir el capital regulatorio entre los activos ponderados por riesgo. Los activos de un banco suelen consistir en efectivo, bonos, acciones y préstamos. A cada clase de activo se le asigna una ponderación por riesgo, en función de cuán arriesgado sea para el banco detentar ese activo. Lógicamente, los bancos necesitan menos capital para cubrir exposiciones a activos de menor riesgo. A mayor cantidad de activos ponderados por riesgo, mayor capital regulatorio resultará necesario. El recuadro I explica los requisitos de capital regulatorios mínimos vigentes en la Unión Europea (Pilar 1 , siguiendo la terminología de Basilea III).\footnote{No se han de confundir con los pilares de la Unión Bancaria.
    } El riesgo de crédito es el principal riesgo al que se enfrentan las entidades bancarias. A este respecto, hay dos grandes enfoques para el cálculo de los activos ponderados por riesgo: el enfoque estandarizado y los modelos internos. Bajo el enfoque estandarizado, los bancos aplican las ponderaciones por riesgo definidas por el supervisor, mientras que en los modelos internos, las entidades calculan sus propias ponderaciones bajo una serie de condiciones. Tal y como se explica en el recuadro 2, éste ha sido uno de los principales elementos modificados bajo la reforma de 2017 de Basilea III y el Paquete Bancario, que transpone dicha reforma al marco regulatorio de la UE.\footnote{%
        Los Acuerdos de Basilea se adoptan en el seno del Comité de Supervisión Bancaria de Basilea, presidido en estos momentos por D. Pablo Hernández de Cos. gobernador del Banco de España, y contienen los estándares internacionales en materia prudencial que luego han de adaptarse a las respectivas jurisdicciones. En el caso de la UE, esto se lleva a cabo por medio de Directivas y Reglamentos.
    }
    
    • la Ratio de Cobertura de Liquidez (LCR, por sus siglas en inglés) y la Ratio de Financiación Neta Estable (NSFR, por sus siglas en inglés), esta última introducida en CRR2. La LCR es un requisito de liquidez orientado al corto plazo, por el cual la entidad tiene que mantener una reserva suficiente de activos líquidos de alta calidad (HQLA, por sus siglas en inglés) que les permita sobrevivir a un periodo de tensión de liquidez significativa de 30 días naturales. La LCR es el resultado de dividir los HQLA entre las salidas netas de efectivo en 30 días y tiene que ser de al menos el 100%. Por otro lado, la NSFR tiene una orientación a mayor plazo (más de un año) y se define como como la cantidad de financiación•estable disponible en relación con la cantidad de financiación estable requerida, que deberá ser superior también al 100%.
    • la ratio de apalancamiento, esta última introducida en CRR2: el apalancamiento se produce cuando los activos de un banco superan su base de capital. La CRR no impide el apalancamiento (intrínseco al propio negocio bancario), pero sí pretende reducir el apalancamiento excesivo que, a la larga, puede comprometer la solvencia del banco. La ratio de apalancamiento viene definida como el capital Tier I de un banco dividido por sus activos totales consolidados medios. Los bancos deberán tener una ratio de apalancamiento mínima del 3'!/o sobre el capital Tier 1. Esta ratio da una medida de la capacidad de un banco para hacer frente a sus obligaciones financieras a largo plazo.
    • los requisitos de transparencia (también conocidos como Pilar 3)
    • los requisitos de grandes exposiciones 
    • los requisitos de retención. para titulizaciones

    Por su parte, CRD recoge:
    • los requisitos gobernanza y control interno, incluyendo el "fit and proper"
    • los requisitos de Pilar 2 (es decir, los requisitos por encima del mínimo regulatorio,
    explicados en el recuadro 1)
    • la autorización
    Al tratarse de una Directiva, es necesario que haya una norma nacional que la transponga y que será la aplicable por el Mecanismo Único de Supervisión.
    
    Por último, hay una serie de cuestiones que son de competencia estrictamente nacional y que, por tanto, escapan tanto a la normativa UE como al Mecanismo Único de Supervisión. Se trata de:
    • Participación de los bancos en el sector no financiero
    • Normativa bancaria de insolvencia
    • Ciertas herramientas macroprudenciales (por ejemplo, LTV, LTI)
    
    Todos estos requisitos son aplicables directamente por el Mecanismo Único de Supervisión, objeto de análisis en la siguiente sección.
\end{specialblock}


\paragraph{Regulación preventiva: }
El bloque preventivo, por su parte, está formado por otros dos instrumentos, ambas directivas y no reglamentos, lo que ya dice algo sobre la voluntad del legislador de dejar aquí más margen de adaptación a las estructuras concursales y administrativas nacionales que en el bloque prudencial. 

La primera es la \href{https://www.boe.es/buscar/doc.php?id=DOUE-L-2014-81284}{Directiva 2014/59/UE del Parlamento Europeo y del Consejo, de 15 de mayo de 2014, por la que se establece un marco para la recuperación y la resolución de entidades de crédito y empresas de servicios de inversión} (de ahora en adelante, DRRB). La finalidad declarada, tal como se lee en sus considerandos, es romper el circulo vicioso entre Estado-Banca que se mencionó al inicio: el principio rector que articula todo el instrumento es que las pérdidas se absorben siguiendo estrictamente el orden de prelación de los instrumentos de capital y deuda de la entidad, empezando por los accionistas (\textit{bail in}), antes de que se pueda recurrir a cualquier fondo público. Además, se sitúa en un orden superior el principio que actúa como garantía última frente a este mecanismo (instrumento jurídico de reoslución): que ningún acreedor puede acabar recibiendo, como consecuencia de la resolución, menos de lo que habría recibido si la entidad hubiera sido liquidada por el procedimiento concursal ordinario (\textit{No Creditor Worse Off},  o NCWO).\footnote{%
    Este segundo principio no es una simple declaración de intenciones sino una garantía exigible ante los tribunales, y de hecho ha sido el fundamento de varios litigios reales contra decisiones de resolución ya adoptadas.
} En cuanto a su estructura, impone a cada entidad la obligación de elaborar, con carácter previo y periódico, un plan de recuperación que ella misma debe preparar para escenarios de estrés severo, y exige a la autoridad de resolución designada en cada Estado que elabore, en paralelo, un plan de resolución para cada entidad, que incluye la determinación de si esa entidad es o no viable en resolución sin apoyo público extraordinario. Solo si la entidad entra en una situación de inviabilidad y concurre un interés público en optar por la resolución en lugar de la liquidación concursal ordinaria, se activan los instrumentos de resolución propiamente dichos, entre los que la recapitalización interna es el más característico pero no el único.

La segunda, la \href{https://www.boe.es/buscar/doc.php?id=DOUE-L-2014-81282}{Directiva 2014/49/UE del Parlamento Europeo y del Consejo, de 16 de abril de 2014, relativa a los sistemas de garantía de depósitos} (de ahora en adelante, DSGD). Completa este bloque preventivo desde el lado del depositante minorista más que desde el lado de la entidad: su finalidad declarada es prevenir el pánico bancario protegiendo el ahorro de los ciudadanos hasta una cuantía determinada por depositante y por entidad, y exige a cada Estado miembro designar uno o varios sistemas de garantía financiados ex ante por las propias entidades adheridas mediante aportaciones periódicas calculadas en función de su perfil de riesgo, de modo que el coste de proteger a los depositantes recaiga sobre la propia industria bancaria y no sobre el contribuyente. 

La reforma (reforma \textit{Crisis Management and Deposit Insurance} o CMDI) que ambas directivas han sufrido en 2026 no altera sus principios, sino que corrige lo que el propio proceso legislativo reconoce como una laguna de aplicación práctica: durante los primeros diez años de vigencia del marco, la inmensa mayoría de las entidades pequeñas y medianas que fracasaron no fueron sometidas al procedimiento de resolución sino liquidadas por vías concursales nacionales dispares y descoordinadas entre sí, precisamente porque el filtro del interés público en la práctica excluía casi automáticamente a estas entidades de menor tamaño. Con esto en mente, la reforma amplía el uso efectivo de la resolución a estas entidades y da a los fondos de garantía de depósitos un papel más amplio como fuente de financiación de la resolución, sin tocar ni la cuantía garantizada al depositante ni el principio de que nadie puede salir peor parado que en un concurso ordinario.

\begin{specialblock}{Marco actual: ¿Qué obligaciones concretas impone la normativa?}
    se apoya normativamente en lo que se conoce como el marco de gestión de crisis y seguro de depósitos (CMDI), compuesto por tres normas: Directiva de Recuperación y Resolución Bancaria (BRRD), Reglamento del Mecanismo-Único de Resolución (RMUR) y Directiva de Fondos de Garantía de Depósitos (DFGD), que se aplican en las diferentes etapas del ciclo de vida de los bancos en dificultades. 
    
    El CMDI ha instaurado importantes principios, como la sustitución del \textit{bail-out}, o rescate de la entidad con cargo a los fondos del contribuyente, por el \textit{bail-in}, por el que el coste de la resolución de los bancos se cubre en primer lugar a través de los propios recursos del banco (es decir, accionistas y acreedores). Para facilitar la aplicación del \textit{bail-in}, se exige a las entidades el cumplimiento con el requerimiento mínimo de fondos propios y pasivos elegibles (MREL, por sus siglas en inglés), que es adicional a los requisitos prudenciales de capital. El Código normativo único de la UE para la prevención y gestión de las quiebras bancarias se ve complementado por los marcos regulatorios nacionales para liquidación de entidades. España no cuenta con regulación específica para la liquidación de entidades, aplicándose la normativa concursal general.
    
    ll.lf. lV. la protección de los depositantes
    Los sistemas de garantía de depósitos (en adelante, SGD) actuales son fruto de una Directiva del año 20 14. Se trata de una ambiciosa reforma con respecto a la Directiva anterior, del año 1994, con el objeto principal encaminada a eliminar las diferencias entre las legislaciones de los Estados Miembros en lo referente a los SGD.
    
    Los sistemas funcionap en el ámbito de cada Estado miembro, y su función es reembolsar a los depositantes hasta un límite definido en caso de que el banco depositario quiebre y los depósitos no estuvieran disponibles.
    
    Los principales elementos de los SGD que resultan armonizados por la Directiva son los siguientes:
    • los SGD contarán con un fondo;
    • la cantidad cubierta por el fondo será de 100.000€ por titular y entidad;
    • se reduce el plazo de pago a siete días para 2024;
    • se mejora la infom1ación para los depositantes;
    • se introducen acuerdos de financiación ex ante fijados, en general, en el 0,8\% de los depósitos con cobertura;
    • se garantiza que sea el sector bancario el que financie los fondos de los sistemas de garantía de depósitos, configurándose un sistema de financiación que responde al siguiente esquema:
    o Las aportaciones en el fondo tienen que ser ex ante.
    o En el supuesto de que el fondo de un Estado no tenga la suficiente liquidéz para proteger los depósitos, para cubrir el posible déficit puede gravar de manera extraordinaria a sus miembros.
    o Una de las exigencias a los Estados, es que el fondo tenga distintas fom1as de financiarse para poder reembolsar o proteger a los depositantes. Los fondos de los distintos f:stados miembros pueden prestarse entre sí.
    o Cuando el fondo paga directamente a los depositantes, asume sus derechos y puede recuperar parte del dinero del patrimonio del banco que ha quebrado (dependiendo eso sí, de la cantidad de pérdidas y recuperaciones).
\end{specialblock}

\begin{specialblock}{¿Cómo se define y calibra el requisito de capital? El MREL}
    II.II.II — El MREL -- Fundamento: BRRD (define y calibra el requisito) + SRMR (para entidades significativas de zona euro, fijado por la SRB). Qué hace: exige a las entidades mantener fondos propios y pasivos elegibles para absorber pérdidas y recapitalizarse, facilitando el bail-in. Estados: aplica a los 27 (toda entidad de la UE tiene requisito MREL); la fijación específica por la SRB solo en zona euro. Última reforma: la CMDI 2026 revisa el uso de fondos de garantía de depósitos para cubrir parte del gap de MREL en bancos medianos/pequeños — antes esto no era así, es contenido nuevo de la reforma de 2026, no de 2019 (BRRD2/SRMR2, que fue la anterior gran calibración de MREL).

    En aras de evitar que los contribuyentes tuvieran que volver a afrontar por medio de dinero público recapitalizaciones bancarias y facilitar así el paso del "bail-out" al "bail-in", el Financia/ Stability Board (FSB) aprobó en 2015 los estándares Total loss Absorbing Capacity (TLAC), refrendados por el G-20. Estos estándares establecen que las entidades globalmente sistémicas (GSlls, por sus siglas en inglés) deben contar con un nivel mínimo de fondos propios y de pasivos con capacidad para absorber pérdidas. Este nivel se fijaría en un mínimo de $18\%$, a cumplir fundamentalmente con capital y deuda subordinada, sobre el total de activos ponderados por riesgo (APR) para 2022.

    El objetivo del TLAC (y del MREL, como transposición de esta figura al ordenamiento jurídico europeo) es que los accionistas y acreedores de las entidades, con algunas excepciones, sean los primeros en asumir las pérdidas de una entidad inviable y, en su caso, de procurar su recapitalización. Así, se requiere a las entidades mantener un importe suficiente de instrumentos de capital y deuda para absorber pérdidas y, en su caso, convertirse en capital, de forma que permita la implementación de la herramienta de resolución que las autoridades determinen que es la más apropiada. Se asegura que dichos instrumentos siempre están disponibles en caso de ser necesaria su utilización, y es independiente de los requerimientos de capital a los que también están sujetas las entidades y que detem1ina su supervisor. Existe una jerarquía para la absorción de pérdidas, que sigue este orden: CET I , Additional Tier 1, Tier 2, deuda subordinada y deuda senior non-preferred. No obstante, este orden no se ha respetado en todas las jurisdicciones a nivel internacional. como se explica en el recuadro 3 . 

    A pesar de que el espejo europeo del TLAC es el MREL, existen importantes diferencias entre ambos, en particular:
    • El TLAC afecta solo a las GSII, mientras que el MREL es exigible a todas las entidades en la UE.
    • El requerimiento mínimo de TLAC es igual para todas las entidades ( 18% de los APR), mientras que el MREL se fija para cada entidad, a partir de las fórmulas y valoraciones previstas en la BRRD2.
    • El TLAC ha de cumplirse con instrumentos subordinados (capital, deuda subordinada y deuda sénior no preferente, con excepciones menores) mientras que el MREL se puede cumplir parcialmente con otros instrumentos. 

    El MREL se debe ��umplir en principio con capital, deuda subordinada, deuda sénior (preferente o no preferente) y depósitos no preferentes ni garantizados con vencimiento a más de un año. No se puede cumplir, en cambio, con depósitos garantizados, derivados o títulos garantizados. No obstante, se pide que una proporción relevante del MREL se cumpla específicamente con instrumentos subordinados: capital, deuda subordinada y deuda sénior no preferente. Esta exigencia de instrumentos subordinados se debe a que estos absprberían pérdidas más fácilmente y con menor litigiosidad legal, a diferencia del resto de pasivos (deuda sénior, depósitos corporativos).

    Para el cumplimiento con el requisito de subordinación, se han clasificado las entidades con arreglo a tres grupos, con mayores requisitos de subordinación para las entidades de mayor tamaño:
    • Grupo 1: GSIIs (destinatarias originales de las obligaciones del TLAC).
    • Grupo 2: grandes entidades: entidades de grupos de resolución con un activo total que
    supere los 100.000 millones de euros.
    • Grupo 3: resto de entidades. 
    No obstante, las autoridades de resolución pueden decidir que ciertas entidades del Grupo 3 reciban el tratamiento más severo del Grupo 2, cuando consideren probable que su inviabilidad plantee un riesgo sistémico. El MREL presenta un componente de absorción de pérdidas, uno de recapitalización y otro para dar confianza al mercado. La lógica subyacente es contar con suficientes recursos corno para absorber las pérdidas de la entidad y recapitalizarla, otorgando suficiente confianza al mercado acerca del nivel de solvencia de la entidad. El I de enero de 2024 es la fecha límite para cumplir con los requerimientos de MREL por parte de las entidades conforme a esta nueva directiva. al margen de las obligaciones de MREL que ya han empezado a imponer, conforme a la anterior directiva, las autoridades de resolución, como la Junta Única de Resolución para las entidades significativas de la unión bancaria. Antes de esa fecha, el I de enero de 2022, las entidades deberán cumplir con unos objetivos intermedios a  determinar por las autoridades de resolución. Las GSII, las grandes entidades y las entidades "asimiladas" deberán cumplir para entonces con los requisitos mínimos de subordinación. Asimismo, las autoridades podrán prorrogar el plazo más allá de 2024 teniendo en cuenta la situación financiera de la entidad, la propia capacidad de sustituir emisiones que vayan venciendo y la capacidad de cumplir con los requisitos en plazo. En caso de incumplimiento del requerimiento de MREL, las autoridades de resolución podrán prohibir a las entidades la retribución de instrumentos de CETl y ATI. Las sanciones pueden imponerse desde el momento del incumplimiento, en proporción al nivel del mismo y en todo caso a partir de los seis meses, salvo perturbación grave de los mercados financieros, entre otras circunstancias. Este tipo de sanciones ya. se preveían en caso de incumplimiento de los requisitos prudenciales e impactan especialmente en las entidades, pudiendo reducir el atractivo de sus emisiones.
\end{specialblock}

\subsubsection{Nivel 2: Actos delegados y/o de ejecución}
Por debajo del Nivel 1 se sitúan las medidas de ejecución de la Unión (Nivel 2): los actos delegados y los actos de ejecución de la Comisión, que desarrollan y aplican esos principios generales y, además, encontraríamos aquí una categoría específica de estos actos de la Comisión: las normas técnicas (\textit{technical standards}), que debemos considera igualmente de segundo nivel. 

Encuentran su fundamento en dos preceptos del TFUE que conviene distinguir con precisión porque responden a lógicas distintas. Uno de ellos habilita a la Comisión para adoptar actos que completan o modifican elementos no esenciales de un acto legislativo, y se activa cuando lo que se necesita es desarrollar sustantivamente una norma que el legislador ha dejado deliberadamente incompleta. Por otro lado, otro habilita a la Comisión para adoptar actos cuya única finalidad es garantizar que una norma ya completa se aplique en condiciones uniformes en todos los Estados miembros, y aquí el contenido normativo no se añade sino que simplemente se estandariza su forma de aplicación, típicamente mediante plantillas y formatos de reporting idénticos para todas las entidades. 

Esta distinción explica por qué a este segundo nivel se le suele llamar, según de qué tipo de acto se trate, normas técnicas de regulación en el primer caso y normas técnicas de ejecución en el segundo, aunque ambas comparten el mismo procedimiento de elaboración.

\paragraph{Autoridad Bancaria Europea: Technical Standards}
El procedimiento arranca siempre en la autoridad técnica correspondiente: la Autoridad Bancaria Europea (ABE).\footnote{%
    Aunque el mismo esquema se replica para la autoridad de valores y mercados y para la autoridad de seguros y pensiones de jubilación en sus respectivos ámbitos materiales
}

La ABE elabora un proyecto de norma, lo somete obligatoriamente a consulta pública y a un análisis de coste y beneficio, y lo remite después a la Comisión Europea. A partir de ahí la Comisión dispone de un plazo, habitualmente de tres meses, para aprobar el proyecto tal cual, rechazarlo, o modificarlo motivadamente antes de adoptarlo como acto delegado o de ejecución propio. La práctica demuestra que la Comisión ha ejercido este poder de modificación cuando el contenido técnico propuesto por la autoridad tenía implicaciones políticas o de coste regulatorio que la Comisión no compartía. 

Una vez adoptado el acto por la Comisión, tanto el Parlamento Europeo como el Consejo disponen todavía de un plazo de objeción, normalmente de tres meses y prorrogable otros tres, durante el cual pueden oponerse al acto en su conjunto, pero no tienen la posibilidad de enmendar el contenido: su control es un veto binario, aceptar o rechazar en bloque, nunca una negociación línea a línea como ocurre en el procedimiento legislativo ordinario del primer nivel. 

La finalidad última de todo este diseño, tal como se explica en los propios reglamentos fundacionales de las autoridades europeas de supervisión, es permitir que el detalle técnico y los formatos de aplicación puedan actualizarse con la rapidez que exige un mercado financiero cambiante, sin necesidad de reabrir cada vez el procedimiento legislativo ordinario completo con su lentitud y su exposición a la negociación política plena entre veintisiete delegaciones nacionales, el Parlamento y la Comisión.


El Single Rulebook también incluye las leyes nacionales que transponen y dan efecto a las directivas pertinentes de la Unión. Habida cuenta de la prohibición del efecto directo vertical inverso, es necesario incluir las leyes que establecen las obligaciones concretas de los participantes en el mercado. Las leyes nacionales también forman parte del fundamento del Derecho sustantivo que debe aplicarse en el marco de la Unión Bancaria de la Unión, tanto en el MUS como en el Mecanismo Único de Resolución (MUR). Por último, las leyes nacionales pueden, en principio, aplicar actos adoptados por las instituciones y organismos de la Unión en cualquiera de los tres niveles.

Antes de continuar, cuestiona tú misma un supuesto que suele darse por bueno sin más: que el segundo nivel es "mero detalle técnico" y el tercero "mera recomendación blanda" son afirmaciones que conviene matizar antes de memorizarlas como si fueran verdades planas, porque en la práctica ocurre algo contraintuitivo: el segundo nivel, pese a llamarse técnico, incorpora decisiones de calado político porque quien tiene la última palabra no es el regulador técnico sino la Comisión, y el tercer nivel, pese a no ser jurídicamente vinculante en sentido estricto, produce en la práctica un grado de cumplimiento tan alto que muchos autores lo consideran derecho duro disfrazado de derecho blando. Si en el examen presentas la jerarquía de los cuatro niveles como una jerarquía también de importancia práctica decreciente, estarás repitiendo una simplificación que no resiste el contraste con cómo funciona realmente el sistema.

\subsubsection{Nivel 3: Autoridad Bancaria Europea}
El tercer nivel abandona por completo la forma de acto jurídico vinculante y adopta la forma de directrices y recomendaciones, cuyo fundamento se encuentra en las disposiciones del reglamento fundacional de la Autoridad Bancaria Europea que regulan sus funciones generales de convergencia supervisora. 

Resumen general que pediste: la ABE, ESMA y EIOPA (y el JERS para riesgo sistémico) emiten guías de aplicación "cumplir o explicar", recomendaciones, y responden preguntas técnicas (Q&A) que sin ser jurídicamente vinculantes en sentido estricto generan una fuerte presión de convergencia supervisora ("soft law" con efecto cuasi-normativo). El instrumento más relevante para tu pregunta anterior sobre divergencia es la Guía del BCE sobre opciones y discreciones (2016, actualizada), que solo vincula a entidades significativas bajo el SSM, no al resto de la UE.

\paragraph{Directrices y recomendaciones: comply or explain}
La naturaleza jurídica de estos instrumentos es lo que la doctrina anglosajona denomina \textit{comply or explain}: la autoridad nacional competente destinataria de una directriz no está obligada a cumplirla en sentido estricto, pero si decide no hacerlo debe notificarlo expresamente a la autoridad europea, habitualmente en un plazo de dos meses desde la publicación, explicando y motivando las razones de su incumplimiento.\footnote{%
    La consecuencia de no cumplir no es una sanción jurídica directa, pero sí una exposición reputacional pública ante el resto de supervisores europeos y ante el mercado, y esta exposición ha resultado en la práctica tan disuasoria que el grado de cumplimiento efectivo de estas directrices ronda niveles muy cercanos a la totalidad.
} 

Las directrices pueden dirigirse indistintamente a las autoridades nacionales competentes o directamente a las propias entidades financieras y demás participantes del mercado, y suelen someterse también a consulta pública previa salvo que concurra una razón de urgencia que lo justifique.

Si el TJUE tuviera que pronunciarse sobre si las directrices y recomendaciones de las ESAs forman parte del Single Rulebook, probablemente se centraría especialmente en la naturaleza de estos actos. Dado que no son jurídicamente vinculantes, sería difícil interpretarlos como parte de un único cuerpo de Derecho de la Unión. Al mismo tiempo, deben tenerse en cuenta el procedimiento específico de \textit{cumplir o explicar} (comply or explain) y el papel especial que ocupan las ESAs en el proceso de toma de decisiones del sector financiero. 

El mero hecho de que corresponda a un Estado miembro (o a su administración) justificar por qué ha adoptado medidas divergentes o se ha abstenido completamente de actuar puede conducir a la conclusión de que la solución propuesta en una directriz o recomendación, aun sin ser jurídicamente vinculante, constituye el estándar efectivo. 

Las directrices y recomendaciones de las ESAs influyen sobre la legislación nacional de una manera muy similar a las directivas. Debido a su objeto y a su relación directa con la legislación de Nivel 1 y Nivel 2, estos actos no pueden ser fácilmente ignorados por los Estados miembros. La elección binaria entre «cumplir o explicar» sitúa a las directrices y recomendaciones de las ESAs cerca de las directivas de armonización mínima, en la medida en que buscan garantizar un nivel mínimo de convergencia de las prácticas. Además, recientemente el legislador de la Unión ha optado por reforzar el peso de estos actos frente a las instituciones de la Unión, más allá de la mera obligación de «cumplir o explicar».

Por las razones expuestas, debemos considerar las directrices y recomendaciones de las ESAs como elementos integrantes del Single Rulebook que, por razones de comodidad, se las denominará actos de Nivel 3, ya que corresponden también al trabajo realizado anteriormente por los Comités de Nivel 3, cuando estos únicamente disponían de facultades de \textit{soft law} en el marco del proceso LAMFALUSSY. 

Clasificar las directrices y recomendaciones un nivel por debajo de las normas técnicas también refleja adecuadamente la jerarquía existente entre estos diferentes actos. Resulta únicamente desafortunado que la base jurídica para la adopción de las directrices y recomendaciones se encuentre normalmente (aunque no siempre) en el acto legislativo de Nivel 1. En este sentido, los actos de Nivel 3 no requieren necesariamente la existencia previa de actos de Nivel 2 para poder adoptarse.

\paragraph{Preguntas y respuestas}
Dentro de este mismo tercer nivel conviene situar también las respuestas a preguntas técnicas que la autoridad europea publica periódicamente, mediante las cuales resuelve dudas interpretativas concretas planteadas por el mercado o por los propios supervisores nacionales sobre la aplicación de una norma ya vigente; su función es cerrar, de manera más ágil y casuística que una directriz general, la brecha que inevitablemente se abre entre lo que la norma escrita dice y lo que la práctica supervisora efectivamente hace en cada Estado miembro, brecha que es precisamente el objeto de buena parte de la literatura sobre divergencia que ya te comenté. 

\paragraph{Acuerdos de mediación}
Existe, finalmente, un instrumento de este tercer nivel que rompe la regla general de la falta de fuerza vinculante y que por eso merece mención separada: cuando dos autoridades nacionales competentes de Estados miembros distintos discrepan sobre un caso transfronterizo concreto y no logran ponerse de acuerdo dentro de un plazo determinado, la ABE puede adoptar una decisión de mediación que resulta jurídicamente vinculante para ambas autoridades nacionales, sustituyendo su desacuerdo por una decisión centralizada; a diferencia de una directriz general, esta decisión no se dirige a todo el sistema sino que resuelve un conflicto individual y concreto, y constituye, dentro de todo el tercer nivel, el único instrumento que se acerca genuinamente al derecho duro en sentido propio.



\textcolor{red}{Las Conclusiones del Consejo Europeo de junio de 2012 (la hoja de ruta hacia la "Genuine EMU", informe Van Rompuy) son el primer texto político que distingue explícitamente entre el rulebook ya existente/en marcha y la necesidad adicional de una "unión bancaria" con supervisión y resolución centralizadas — ahí sí encontrarás la dualidad enunciada como tal, aunque en lenguaje de conclusiones del Consejo, no como norma jurídica vinculante.




\subsection{Ámbito subjetivo: Una arquitectura institucional dual}
Como vemos, el CNU es un cuerpo normativo de alcance uniforme: se aplica, en su contenido sustantivo, a los veintisiete Estados miembros de la Unión sin distinción alguna, sea cual sea su relación con la moneda única.

Sin embargo, ese mismo cuerpo normativo no se aplica de la misma manera en todos los Estados, porque quién verifica su cumplimiento (quién ejerce, en último término, la función de supervisión y, en su caso, de resolución) varía radicalmente según el Estado miembro de que se trate. Esta disociación entre uniformidad de la norma sustantiva y dualidad de su aplicación institucional es la clave que permite entender por qué la Unión Bancaria y el Código Normativo Único, pese a solaparse en buena parte de su contenido, no son ni pueden presentarse como la misma cosa: el primero es una capa institucional de aplicación centralizada que se superpone, solo para determinados Estados, sobre el segundo, que es el sustrato normativo común a todos.

Esta arquitectura dual no es una construcción doctrinal a posteriori, sino que está anclada en dos bases jurídicas del Tratado de naturaleza deliberadamente distinta. Por un lado, la legislación de contenido sustantivo del Código Normativo Único se adopta bajo el artículo $114$ del TFUE, mediante procedimiento legislativo ordinario con codecisión plena de Parlamento y Consejo, y su ámbito de aplicación territorial es, por defecto, el conjunto de la Unión. Por otro lado, la atribución de funciones de supervisión directa al BCE se realiza bajo un fundamento jurídico distinto y más restrictivo, el artículo $127(6)$ del TFUE (la \textit{cláusula habilitante}), que exige unanimidad en el Consejo y mera consulta, no codecisión, al Parlamento Europeo, precisamente porque no crea norma sustantiva nueva sino que traslada a una institución concreta, el BCE, una función que de otro modo correspondería a cada Estado.\footnote{%
    Esta segunda base jurídica no tiene vocación de aplicarse a los veintisiete: por su propia naturaleza de atribución de una tarea concreta a una institución de la eurozona, su ámbito natural es la propia eurozona, ampliable únicamente mediante el mecanismo de cooperación estrecha que regula el art. $7$ del Reglamento (UE) n.º 1024/2013, por el que un Estado no euro puede solicitar voluntariamente su incorporación al Mecanismo Único de Supervisión sin necesidad de adoptar la moneda única, categoría que, como acabamos de corregir, no tiene hoy ningún miembro activo tras la incorporación plena de Bulgaria a la eurozona.
}

Esta diferencia de base jurídica se traduce en dos redes institucionales distintas que conviene no confundir. 

\subsubsection{Normativa: Sistema Europeo de Supervisión Financiera (2010)}
La primera, de alcance pleno a los veintisiete Estados miembros, es el Sistema Europeo de Supervisión Financiera, creado en 2010 a raíz precisamente del Informe LAROSIÈRE, y compuesto por:
\begin{la}
    \item \textbf{Junta Europea de Riesgo Sistémico}. La Junta Europea de Riesgo Sistémico, encargada de la supervisión macroprudencial mediante recomendaciones de cumplir o explicar; 
    \textbf{Autoridades Europeas de Supervisión}. Las tres Autoridades Europeas de Supervisión: la Autoridad Bancaria Europea para el sector bancario, la Autoridad Europea de Valores y Mercados para los mercados de capitales, y la Autoridad Europea de Seguros y Pensiones de Jubilación para seguros y pensiones; 
    \textbf{Autoridades Nacionales Competentes}. Las autoridades nacionales competentes de los veintisiete; y,
    \item \textbf{Comité Conjunto}. Un Comité Conjunto de las tres Autoridades para cuestiones transversales, como los conglomerados financieros. 
\end{la}

Este sistema es el que redacta, mediante los Niveles 2 y 3, el contenido técnico del Código Normativo Único, y su ámbito de aplicación coincide exactamente con el de ese Código: los veintisiete, sin excepción. La Autoridad Bancaria Europea, en particular, es una autoridad de la Unión con competencia sobre el conjunto del mercado interior, y su función es idéntica se trate de un Estado de la Unión Bancaria o no (redactar normas técnicas, emitir directrices, mediar entre supervisores nacionales).

\paragraph{Superpuesta: Instituciones de la Unión Bancaria}
La segunda red institucional, de alcance restringido, es la propia Unión Bancaria: el Mecanismo Único de Supervisión y el Mecanismo Único de Resolución que veremos más adelante, más el todavía inexistente Esquema Europeo de Garantía de Depósitos. 

Esta segunda red no sustituye a la primera sino que se superpone a ella, exclusivamente para los Estados participantes en la Unión Bancaria: dentro de ese subconjunto de Estados, la función de supervisión directa que en el resto de la Unión desempeña la autoridad nacional competente pasa a ser ejercida, para las entidades significativas, por el Banco Central Europeo, y la función de resolución que en el resto de la Unión desempeña la autoridad nacional de resolución pasa a ser ejercida, para esas mismas entidades, por la Junta Única de Resolución. 

La diferencia entre ambos grupos de Estados no es, por tanto, de contenido normativo sino de arquitectura de aplicación.

\subsubsection{Supervisión: ¿Cómo se traduce esta dualidad en la supervisión?} 
Dentro de la Unión Bancaria, el Código Normativo Único se aplica mediante una supervisión que hemos calificado ya de centralizada aunque de intensidad variable: el Banco Central Europeo ejerce supervisión directa sobre las entidades significativas a través de los Equipos Conjuntos de Supervisión, y supervisión indirecta (mediante reglamentos, orientaciones e instrucciones generales, y la facultad de avocación en cualquier momento) sobre las entidades menos significativas, que siguen siendo supervisadas en el día a día por su autoridad nacional competente pero ya no de forma autónoma, sino integrada en un mecanismo único bajo la dirección última del propio Banco Central Europeo. 

De manera equivalente, la resolución de estas entidades, cuando resulta necesaria, se decide de forma centralizada por la Junta Única de Resolución, con el respaldo financiero del Fondo Único de Resolución, en lugar de depender exclusivamente de la autoridad nacional de resolución y de un fondo nacional de garantía de depósitos aislado.

Fuera de la Unión Bancaria, en el resto de la Unión, el mismo Código Normativo Único se aplica bajo una arquitectura completamente distinta: la supervisión y la resolución siguen siendo competencia exclusiva de la autoridad nacional competente y de la autoridad nacional de resolución respectivamente, sin que exista ninguna instancia europea con poder de decisión directa sobre una entidad concreta. 

Lo que impide que esta descentralización total degenere en la misma fragmentación regulatoria que existía antes de la crisis financiera global es, precisamente, el papel de la Autoridad Bancaria Europea como garante de la aplicación coherente del Código Normativo Único en todo el mercado interior. Ahora bien, la Autoridad Bancaria Europea no sustituye a la autoridad nacional en ningún momento, pero limita el margen de divergencia práctica entre ellas de un modo que no es en absoluto desdeñable.


Conviene resaltar que la línea que separa a los Estados de la Unión Bancaria del resto de la Unión no traza una frontera de todo o nada entre quienes están "dentro" del Código Normativo Único y quienes están "fuera" de él, porque todos los veintisiete están igualmente dentro de él en cuanto a su contenido sustantivo. Lo que traza esa línea es, exclusivamente, el grado de centralización institucional con el que ese contenido común se hace efectivo: centralización plena, con capacidad decisoria europea directa sobre entidades concretas, dentro de la Unión Bancaria; coordinación y mediación, sin capacidad decisoria directa sobre ninguna entidad concreta, fuera de ella. Esta distinción entre uniformidad normativa y dualidad institucional es, en definitiva, la clave conceptual, jurídica y económica que permite entender por qué la Unión Bancaria puede describirse con propiedad como "el Código Normativo Único bajo una arquitectura institucional de aplicación centralizada", una fórmula que, como verás en el epígrafe siguiente, es exactamente la que emplearemos para introducir la propia Unión Bancaria una vez fijados estos límites conceptuales.




































\clearpage
\section{Unión Bancaria: Sistema Bancario de la Eurozona}
\puntosclave{
    El Mecanismo Único de Supervisión (MUS) es el sistema de supervisión bancaria en buena parte de los Estados miembros de la UE (los países participantes en el MUS son todos los de la zona euro y Bulgaria). Está formado por el BCE y las autoridades nacionales de supervisión de los países participantes. En 2024, el MUS celebrará su décimo aniversario. 
    
    El BCE supervisa directamente los 11O bancos más importantes de los países participantes en el MUS, que poseen casi el $82\%$ de los activos bancarios de estos países. 
    
    Los Equipos Conjuntos de Supervisión (JST, por sus siglas en inglés) son una de las principales formas de cooperación entre el BCE y los supervisores nacionales. Los JST llevan a cabo la supervisión continuada de las entidades de crédito significativas. Los bancos que no se consideran significativos se conocen como entidades "menos significativas". Siguen siendo supervisadas por sus supervisores nacionales, en estrecha colaboración con el BCE. 
    
    El MUS es el elemento más sólido de la Unión Bancaria hasta la fecha y el que ha permitido, por ejemplo, que el sistema bancario europeo haya resistido holgadamente las turbulencias financieras desatadas en marzo de 2023 en EE.UU. y Suiza. 
    
    El SREP es la actividad nuclear de la supervisión bancaria europea, que proporciona a los supervisores un conjunto de herramientas para evaluar de manera homogénea los riesgos a los que se enfrentan las entidades de crédito y comprobar que éstas los gestionan adecuadamente sobre la base de cuatro criterios: ( 1 ) capital, ( 2) liquidez, (3) modelo de negocio y (4) gobemanza.

    El Mecanismo Único de Resolución (MUR) está compuesto por: (1) la Junta Única de Resolución (SR.B, por sus siglas en inglés), ( 2) el Fondo Único de Resolución (SRF, por sus siglas en inglés) y (3) las Autoridades Nacionales de Resolución. 
    
    Además, está previsto que una vez se ratifique la reforma del Tratado del Mecanismo Europeo de Estabilidad, a fecha de redacción de este tema pendiente de ratificación por parte del Parlamento italiano, el SRF cuente con un respaldo fiscal o \textit{common backstop}. 
    
    El SRB cuenta con una Sesión Plenaria y una Sesión Ejecutiva. En la Sesión Plenaria del SRB participan todas las Autoridades Nacionales de Resolución de los países participantes en la Unión Bancaria, junto con los seis miembros permanentes del SRB. En la Sesión Ejecutiva del SRB participan los seis miembros pennanentes y los observadores antes mencionados. 
    
    El SRF tiene como objetivo asegurar la eficiente aplicación de las herramientas de resolución y facilitar la ejecución de las medidas acordadas por el SRB. El objetivo para el SRF en 2024 es alcanzar al menos el $1\%$ de los depósitos cubiertos de la eurozona (estimado actualmente por encima de los 70.000 millones de euros). El SRF se financia con contribuciones ex ante de las entidades de crédito y de algunas sociedades de inversión. En el supuesto de que un caso de resolución (o varios) consuma todos los recursos disponibles en el SRF y sea necesaria más financiación, se recaudarán contribuciones ex post, soportadas por las mismas entidades financieras. Normalmente, esas contribuciones ex post no están disponibles o no son accesibles de inmediato, por lo que el SRF deberá recurrir a operaciones de endeudamiento, ya sean privadas o públicas. Solo en el supuesto de que estos recursos no fueran suficientes, se activaría el "common backstop", cuando esté disponible en el futuro.

    En el año 2015, la Comisión Europea presentó su propuesta normativa para establecer, progresivamente, un Esquema de Garantía de Depósitos Único (EDIS, por sus siglas en inglés). La propuesta de EDIS se divide en tres etapas sucesivas: (i) reaseguro, (ii) coaseguro y (iii) seguro pleno. 
    
    Durante la fase de reaseguro, los SGD participantes podrían obtener financiación en aquellos casos en los que presentaran falta de liquidez, bien debida a un evento de desembolso o por tener que contribuir a una resolución. Asimismo, se podría llegar a la asunción de pérdidas, pero por un importe máximo del $20\%$ del exceso de pérdidas. 
    
    En la fase de coaseguro, con una duración esperada de 4 años, se aumenta el grado de mutualización de riesgos: ya no sería necesario agotar en un primer momento los recursos del SGD nacional, para acudir al EDIS.
    
    En la fase de seguro pleno, el EDIS ofrecería a los SGD nacionales la financiación íntegra de la necesidad de liquidez y cubriría todas las pérdidas derivadas de un evento de desembolso o una solicitud de contribución a una resolución. El mecanismo sería el mismo que en la etapa de coaseguro, con la diferencia de que el SEGD cubría un $100\%$. 
    
    La financiación del EDIS se llevaría a cabo mediante aportaciones ex ante de los bancos participantes y llegaría a cubrir hasta el $0,8\%$ de los depósitos garantizados. La gestión del EDIS correspondería al SRB, aunque la figura de los SGD nacionales y de las correspondientes autoridades competentes para su gestión no desaparecería, sino que sus funciones se adaptarían. 
    
    Las negociaciones sobre el EDIS se encuentran totalmente bloqueadas.
}

Antes de desarrollar en detalle esa arquitectura centralizada, conviene justificar por qué este tema va a detenerse en ella con tanta extensión y no en la arquitectura descentralizada que rige para el resto de la Unión, ya suficientemente descrita en el epígrafe anterior.

La pertenencia a la Unión Bancaria no es, para un Estado miembro, una opción de política económica entre otras: es la consecuencia automática de una obligación que sí forma parte del acervo comunitario, la de adoptar el euro, recogida sin excepción en los Tratados de Adhesión de todo Estado que se haya incorporado a la Unión Europea desde el Tratado de Maastricht. El único Estado miembro que cuenta con una exclusión formal y permanente respecto de esta obligación es Dinamarca, en virtud del Acuerdo de Edimburgo de 1992.\footnote{%
    Suecia se encuentra en una posición distinta que conviene no confundir con la danesa: carece de un opt-out jurídico formal, pero evita voluntariamente cumplir los criterios de convergencia y permanecer fuera del Mecanismo de Tipos de Cambio II, lo que le permite posponer de facto, aunque no de iure, una obligación de adopción del euro que en teoría sigue vigente para ella.
}. Todos los demás Estados miembros están destinados, tarde o temprano, a adoptar la moneda única y, con ella, a incorporarse automáticamente al ámbito de aplicación del Mecanismo Único de Supervisión, sin necesidad de acto de adhesión distinto alguno. La condición de Estado ajeno a la Unión Bancaria es, por tanto, y salvo para Dinamarca, una situación transitoria.

Sentada esta premisa, el desarrollo que sigue se limita, por tanto, al Código Normativo Único tal como se aplica bajo la arquitectura institucional propia de la Unión Bancaria, es decir, a sus tres pilares: el Mecanismo Único de Supervisión, el Mecanismo Único de Resolución, y el todavía inexistente Esquema Europeo de Garantía de Depósitos.


\textcolor{red}{Desde el punto de vista de la política monetaria, en mayo de 2010, el BCE (BCE) puso en marcha el \textit{Securities Market Program} (SMP), diseñado para comprar temporalmente deuda soberana. Este programa fue sustituido en agosto de 2012 por el programa \textit{Outright Monetary Transactions} (OMT), por el que el BCE se comprometió a comprar deuda soberana de los países de la zona del euro en dificultades, cumpliendo una serie de condiciones, en los mercados secundarios. Otras medidas encaminadas a aumentar el crédito al sector privado y reducir los costes crediticios fueron las \textit{Targeted long Term Refinancing Operations}, los tipos negativos en las facilidades de depósito y un programa ampliado de compra de activos, que comenzaron a aplicarse en 2014.

En cuanto a la supervisión y regulación bancaria, la Autoridad Bancaria Europea publicó los primeros resultados oficiales de las pruebas de resistencia de los bancos en octubre de 2011. A su vez, las autoridades de la zona del euro anunciaron en junio de 2012 la futura creación de la Unión Bancaria Europea. La Unión Bancaria pretende evitar tanto la fragmentación bancaria europea como los rescates públicos a los bancos en quiebra. Para ello, la Unión Bancaria se basa en un Código normativo único y en diferentes pilares. En primer lugar, el Mecanismo Único de Supervisión comenzó a funcionar en noviembre de 2014, con el objetivo de garantizar la am10nización de las prácticas de supervisión bancaria en todos los Estados miembros que participan en la UEM El segundo pilar, el Mecanismo Único de Resolución, inició su andadura en 2015. El tercer y último pilar, el Fondo de Garantía de Depósitos Único, aún no existe y las negociaciones políticas para su creación están totalmente estancadas.}

 



\subsection{Pilar 1: Mecanismo Único de Supervisión}
El Mecanismo Único de Supervisión es el sistema de supervisión bancaria común a los Estados de la Eurozona y a cualquier Estado no perteneciente a la eurozona que haya establecido una cooperación reforzada con el BCE.\footnote{%
    Tras la incorporación plena de Bulgaria a la eurozona el 1 de enero de 2026, y con ella al Mecanismo Único de Supervisión como miembro pleno tras su periodo de cooperación reforzada iniciado en 2020, esta segunda categoría ha quedado vacía en la práctica: hoy los 21 Estados participantes en el Mecanismo Único de Supervisión coinciden exactamente con los 21 Estados de la eurozona.
} Es decir, el sujeto obligado a aplicar la regulación contenida en el Código Normativo Único, verificando que las entidades de crédito la cumplen, sin que ello suponga en ningún caso una potestad para crear esa norma sustantiva. 

Su fundamento jurídico constitutivo se encuentra en el Reglamento (UE) n.º 1024/2013 del Consejo, adoptado sobre la base del art. $127(6)$ del TFUE, que atribuye al BCE las funciones específicas de supervisión prudencial de las entidades de crédito establecidas en los Estados participantes.\footnote{%
    Distinta de esta norma constitutiva, y subordinada a ella, es el Reglamento (UE) n.º 468/2014 del BCE, de 16 de abril de 2014, por el que se establece el marco de cooperación en el Mecanismo Único de Supervisión entre el BCE y las autoridades nacionales competentes y con las autoridades nacionales designadas, comúnmente denominado Reglamento Marco del Mecanismo Único de Supervisión. 

No debe confundirse con el acrónimo del Mecanismo Único de Resolución, con el que comparte semejanza fonética en español. El Reglamento Marco no es una norma del legislador europeo sino un reglamento adoptado por el propio BCE en ejercicio de una potestad normativa instrumental, distinta de la potestad legislativa ordinaria del primer nivel del Código Normativo Único y de la potestad delegada del segundo nivel: se trata de una tercera vía normativa, específica del ámbito supervisor, cuya única función es desarrollar las modalidades operativas de cooperación entre el BCE y las autoridades nacionales dentro del marco ya creado por el Reglamento 1024/2013, y de establecer la metodología de clasificación de las entidades como significativas o menos significativas.
} En consecuencia, el MUS está formado por el BCE y las ANC de los países participantes, y sus objetivos declarados son \textbf{garantizar la seguridad y solidez del sistema bancario europeo, aumentar la integración y la estabilidad financieras, y garantizar una supervisión coherente entre entidades y entre Estados}.

\subsubsection{Arquitectura institucional y ámbito subjetivo}
Los principales objetivos de la supervisión bancaria europea, ya enunciados, se despliegan a través de un conjunto de facultades atribuidas al BCE: realizar revisiones supervisoras, inspecciones in situ e investigaciones; conceder o retirar licencias bancarias; evaluar la adquisición y enajenación de participaciones cualificadas por parte de los bancos; garantizar el cumplimiento de las normas prudenciales de la Unión; y fijar requisitos de capital más elevados, los llamados colchones, para contrarrestar cualquier riesgo financiero detectado.

Dentro de esta arquitectura conviene detenerse en cómo se articula internamente la toma de decisiones, porque es un elemento que suele omitirse y que resulta indispensable para entender el funcionamiento real del mecanismo. 

Sobre esta base institucional, el Mecanismo Único de Supervisión despliega su actividad diferenciando, como veremos a continuación, entre entidades significativas y menos significativas, ejerciendo sobre ellas facultades de intensidad distinta y canalizando su núcleo de actividad a través del Proceso de Revisión y Evaluación Supervisora.

\textcolor{red}{El BCE supervisa directamente, a fecha de la última actualización de la lista de entidades supervisadas, en torno a ciento diez a ciento catorce bancos de los países participantes en el Mecanismo Único de Supervisión (además realiza evaluaciones puntuales cuando se produce una autorización, una revocación de licencia, o un cambio relevante en la estructura de un grupo), que en conjunto poseen aproximadamente el ochenta y dos por ciento de los activos bancarios del conjunto de los Estados participantes. Los criterios para determinar si una entidad de crédito se considera significativa, y por tanto queda sometida a la supervisión directa del BCE, se establecen en el Reglamento del Mecanismo Único de Supervisión y en su Reglamento Marco, y son los que ya trabajamos en la pregunta anterior sobre entidades significativas y menos significativas. Los Equipos Conjuntos de Supervisión están formados por personal del BCE y de los supervisores nacionales pertinentes: en concreto, cada equipo cuenta con un coordinador del BCE, subcoordinadores nacionales, y un equipo de expertos, y su tamaño, composición general y organización se adaptan al tamaño, al modelo de negocio y al perfil de riesgo del banco supervisado.}

Así que la cadena completa, si quieres memorizarla como tal, es: el Equipo Conjunto de Supervisión hace el trabajo continuo de supervisión y genera la propuesta técnica; el Consejo de Supervisión delibera sobre esa propuesta y la convierte en una decisión formal; y el Consejo de Gobierno del Banco Central Europeo la adopta jurídicamente mediante el silencio positivo del procedimiento de no objeción, salvo que decida oponerse expresamente dentro del plazo. [Probable] Si en el examen te preguntan quién supervisa a una entidad significativa concreta, la respuesta más precisa es "el Equipo Conjunto de Supervisión correspondiente, bajo la dirección última y la responsabilidad decisoria del Consejo de Supervisión del Banco Central Europeo" — mencionar solo uno de los dos deja fuera la mitad de la arquitectura que tú misma pediste que incluyera en el subepígrafe correspondiente.

Como regla general, las decisiones de supervisión ordinarias sobre una entidad menos significativa —el resultado de su proceso de revisión y evaluación supervisora, las medidas cualitativas que deba adoptar, las inspecciones que se le practiquen— las toma la propia autoridad nacional competente de forma autónoma, sin necesidad de elevarlas al Consejo de Supervisión del Banco Central Europeo caso por caso. Ahora bien, hay un conjunto de decisiones que la doctrina y la propia normativa denominan procedimientos comunes, y que se reservan siempre al Banco Central Europeo con independencia de si la entidad afectada es significativa o menos significativa, fundamentalmente: la autorización para constituir una entidad de crédito, la retirada de esa autorización, y la evaluación de la adquisición o el incremento de participaciones cualificadas en una entidad de crédito. En estos tres supuestos concretos, la autoridad nacional competente no decide por sí misma ni siquiera si la entidad es menos significativa: lo que hace es instruir el procedimiento, recabar la información necesaria, y elevar un proyecto de decisión al Banco Central Europeo, que es quien adopta la decisión final mediante el mismo procedimiento de no objeción del Consejo de Gobierno que ya conoces, a propuesta del Consejo de Supervisión.\footnote{%
    La razón de fondo de esta excepción es que conceder o retirar una licencia bancaria, o permitir que alguien tome el control de una entidad de crédito, son decisiones que afectan a la integridad del sistema bancario en su conjunto y a la aplicación coherente de estándares en todo el mercado único, con independencia del tamaño de la entidad concreta afectada — no tendría sentido que la entrada o salida de un actor del sistema bancario europeo quedara sujeta a criterios puramente nacionales cuando todo lo demás en el marco prudencial persigue precisamente evitar esa fragmentación.
}

\paragraph{Órgano Ejecutivo: Consejo de Supervisión}
El Consejo de Supervisión no ejecuta esa supervisión continuada; su función es distinta y se sitúa un escalón por encima, en el nivel decisorio. Es el órgano que prepara las decisiones formales de supervisión —una exigencia de capital adicional de Pilar dos, la retirada de una licencia, el resultado del proceso de revisión y evaluación supervisora— a partir precisamente del trabajo que le eleva el Equipo Conjunto de Supervisión responsable de esa entidad, y las somete después al procedimiento de no objeción del Consejo de Gobierno del Banco Central Europeo, que ya vimos.

Las decisiones de supervisión se preparan en el seno de un \textbf{Consejo de Supervisión}, integrado por representantes del BCE y de las ANCs, y se elevan después al Consejo de Gobierno del propio BCE mediante un \textbf{procedimiento de no objeción}: el Consejo de Gobierno no aprueba activamente cada decisión de supervisión, sino que dispone de un plazo para oponerse a la propuesta del Consejo de Supervisión, y si no lo hace, la decisión se entiende adoptada.\footnote{%
    No obstante, ningún ejercicio de una potestad pública de esta intensidad queda exento de control, y el Mecanismo Único de Supervisión cuenta con dos vías de fiscalización que conviene distinguir con precisión porque operan en planos distintos. La primera es un control administrativo previo, encarnado en el Comité Administrativo de Revisión, un órgano interno al propio BCE pero dotado de independencia funcional, que examina, a instancia de la entidad afectada, la legalidad procedimental y sustantiva de las decisiones de supervisión antes de que estas puedan ser impugnadas ante la jurisdicción de la Unión. El dictamen del Comité Administrativo de Revisión no es vinculante para el Consejo de Supervisión, que puede apartarse de él motivadamente al elaborar una nueva propuesta de decisión, pero en la práctica constituye un filtro que reduce sustancialmente el volumen de decisiones que terminan siendo recurridas ante el Tribunal de Justicia de la Unión Europea. La segunda es una rendición de cuentas de naturaleza política y periódica, ejercida ante el Parlamento Europeo, mediante comparecencias regulares de la presidencia del Consejo de Supervisión ante su Comisión de Asuntos Económicos y Monetarios, de las que la comparecencia de Claudia Buch en marzo de 2026 constituye el ejemplo más reciente, y que exige al BCE justificar públicamente el ejercicio de una función supervisora que, no debe olvidarse, actúa sobre un sector cuya estabilidad afecta directamente a la protección del ahorro de millones de ciudadanos europeos.
}

Esta arquitectura responde a un principio estructural conocido como principio de separación (\textit{Chinese Walls}), que exige mantener funcionalmente distintas, dentro de una misma institución, la función de supervisión y la función de Política Monetaria, precisamente para evitar que los objetivos de estabilidad de precios contaminen las decisiones supervisoras concretas sobre una entidad, y viceversa, que consideraciones supervisorias sobre una entidad concreta condicionen indebidamente la Política Monetaria general de la eurozona.\footnote{%
    El fundamento último de este principio de separación es evitar el llamado riesgo de conflicto de intereses institucional: un banco central que fija tipos de interés y simultáneamente decide si una entidad concreta debe elevar su capital podría verse tentado a suavizar exigencias supervisoras si estas amenazan la transmisión de su política monetaria, o a la inversa. El Consejo de Supervisión cuenta además con un Comité Directivo que apoya sus actividades y prepara sus reuniones, sin capacidad decisoria propia.
} 

\paragraph{Órgano operativo: Equipos Conjuntos de Supervisión}
El BCE supervisa directamente, a fecha de la última actualización de la lista de entidades supervisadas, en torno a ciento diez a ciento catorce bancos de los países participantes en el Mecanismo Único de Supervisión (además realiza evaluaciones puntuales cuando se produce una autorización, una revocación de licencia, o un cambio relevante en la estructura de un grupo), que en conjunto poseen aproximadamente el ochenta y dos por ciento de los activos bancarios del conjunto de los Estados participantes.\footnote{%
    Estas cifras deben tratarse como aproximadas y sujetas a revisión periódica. La incorporación plena de Bulgaria en enero de 2026 ha añadido entidades a la lista de supervisión directa sin alterar sustancialmente ni el número total ni el porcentaje de activos concentrados, dado el tamaño relativamente reducido del sistema bancario búlgaro respecto al conjunto de la eurozona.
} 

La principal forma de cooperación operativa entre el BCE y los supervisores nacionales son los Equipos Conjuntos de Supervisión, que llevan a cabo la supervisión continuada de las entidades de crédito significativas.\footnote{
Las entidades que no se consideran significativas se conocen como entidades menos significativas y siguen siendo supervisadas en el día a día por sus ANCs, en estrecha colaboración con el BCE, que conserva sobre ellas una función de vigilancia general y la facultad de asumir en cualquier momento su supervisión directa si lo considera necesario para garantizar la aplicación coherente de estándares de supervisión estrictos, sin que la autoridad nacional pueda oponerse a esa avocación de competencia.
} 

La supervisión continuada del día a día —el seguimiento permanente de la entidad, la ejecución del programa de examen supervisor, las inspecciones in situ, el análisis constante de su situación de riesgo— la realizan los Equipos Conjuntos de Supervisión. Son ellos quienes están, en sentido operativo, "encima" de la entidad significativa de manera constante a lo largo de todo el año: es el nivel de ejecución.

El coordinador del equipo, que generalmente no procede del país en el que se encuentra el banco supervisado (salvaguarda deliberada frente a la captura regulatoria nacional), dirige el equipo y orienta sus actividades de supervisión, y es nombrado por un periodo de tres a cinco años; los subcoordinadores nacionales son responsables de áreas temáticas o geográficas claramente definidas y apoyan al coordinador en la supervisión en curso, y en los equipos más grandes, el coordinador y los subcoordinadores nacionales forman lo que se conoce como un equipo conjunto central, que distribuye las tareas entre el resto de miembros. 

Las principales tareas de estos equipos son llevar a cabo el Proceso de Revisión y Evaluación Supervisora, proponer el programa de examen supervisor, incluido un plan de inspecciones in situ, aplicar el programa de examen supervisor ya aprobado y cualquier decisión de supervisión derivada de él, y garantizar la coordinación con los equipos de inspección in situ, sirviendo de enlace permanente con los supervisores nacionales.

\subsubsection{Proceso de Revisión y Evaluación Supervisora}
El Proceso de Revisión y Evaluación Supervisora (SREP), es la actividad nuclear de la supervisión bancaria europea: \textbf{proporciona a los supervisores un conjunto de herramientas para evaluar de manera homogénea los riesgos} a los que se enfrenta cada entidad de crédito y comprobar que esta los gestiona adecuadamente, sobre la base de cuatro criterios: (i) capital; (ii) liquidez; (iii) modelo de negocio; y, (iv) gobernanza.\footnote{%
    El origen del Proceso de Revisión y Evaluación Supervisora como marco común no está en el Reglamento del Mecanismo Único de Supervisión ni en su Reglamento Marco, sino en las Directrices sobre el propio proceso que emite la Autoridad Bancaria Europea, es decir, en el tercer nivel del Código Normativo Único que trabajamos hace varias preguntas, bajo el mecanismo de cumplir o explicar. Esas directrices vinculan, en el sentido débil-pero-eficaz que ya describimos, a todas las autoridades competentes de los veintisiete Estados miembros, sean o no de la Unión Bancaria, precisamente porque el proceso de revisión y evaluación no es una institución exclusiva del Mecanismo Único de Supervisión sino un desarrollo general de la Directiva de Requisitos de Capital, que ya vimos que se aplica a los veintisiete. Así que en ese sentido amplio tienes razón: el marco metodológico común procede de la Autoridad Bancaria Europea y alcanza a cualquier autoridad competente de la Unión, dentro o fuera de la Unión Bancaria.

Dentro del Mecanismo Único de Supervisión específicamente, sin embargo, el grado de uniformidad no es idéntico entre entidades significativas y menos significativas, y aquí es donde tu pregunta necesita matiz. Para las entidades significativas, el Banco Central Europeo no se limita a aplicar las directrices generales de la Autoridad Bancaria Europea sino que ha desarrollado su propia metodología de supervisión, más detallada y específica, que aplican directamente los Equipos Conjuntos de Supervisión bajo la dirección última del Consejo de Supervisión, con la decisión final adoptada, como ya sabes, mediante el procedimiento de no objeción del Consejo de Gobierno. Para las entidades menos significativas, es la autoridad nacional competente quien aplica el proceso de revisión y evaluación, partiendo también de las directrices de la Autoridad Bancaria Europea pero con un grado de proporcionalidad considerablemente mayor: no todas las entidades menos significativas reciben el mismo nivel de escrutinio, sino que existe dentro de esta categoría una subdistinción entre entidades menos significativas de alto impacto y el resto, de modo que las autoridades nacionales pueden aplicar un proceso simplificado, con menor frecuencia y menor profundidad de análisis, a las entidades menos significativas de menor tamaño relativo.
}

Al final de cada ejercicio anual del proceso, el supervisor comunica a la entidad los resultados, estableciendo una serie de objetivos que debe corregir dentro de un plazo determinado; en el marco de este proceso se fijan, de manera individualizada para cada entidad, los requisitos de capital de Pilar dos, que complementan los requisitos mínimos generales de Pilar uno para cubrir riesgos no contemplados por este último, y se comunican también las directrices de Pilar dos, orientadas a que la entidad esté en condiciones de afrontar eventuales tensiones financieras adicionales. 

Como parte de este mismo proceso, el supervisor analiza además los procesos internos de evaluación de la adecuación del capital y de la liquidez que cada entidad debe realizar sobre sí misma, conocidos respectivamente como ICAAP e ILAAP.

Este proceso, que durante casi una década funcionó con un diseño prácticamente inalterado, fue sometido a evaluación externa por encargo del propio BCE, que en septiembre de 2022 encomendó a cinco expertos de alto nivel en supervisión bancaria un análisis independiente sobre su eficacia y eficiencia y sobre su relación con el resto de procesos de supervisión. Tras siete meses de acceso a documentación del BCE y más de setenta reuniones, el grupo de expertos produjo una evaluación cuyos resultados se publicaron en abril de 2023, y que confirmaba que, desde 2014, el BCE se había consolidado con éxito como un supervisor eficaz y respetado, sobre todo teniendo en cuenta la elevada heterogeneidad de partida de los procesos supervisores nacionales previos a la Unión Bancaria, al tiempo que consideraba que la organización era ya lo suficientemente sólida y madura como para agilizar sus procesos y mejorar su priorización basada en riesgo. 

\subsubsection{Balance: }
Las recomendaciones del informe se articulaban en tres grandes áreas: cultura supervisora, notas del proceso de revisión, y medidas cualitativas. En materia de cultura supervisora, el informe sugería pasar de un enfoque rígido a otro basado en riesgos, aumentar la eficiencia del proceso reduciendo su duración entre tres y cuatro meses, integrar mejor en el análisis los resultados de otras actividades supervisoras, e incrementar la claridad respecto a las expectativas del supervisor. En cuanto a las notas del proceso, los expertos recomendaban reducir el sesgo mecanicista existente, evitando una traslación automática entre la nota obtenida y el nivel de requisito de capital de Pilar dos, y favoreciendo en su lugar un enfoque riesgo a riesgo con un uso más limitado del proceso interno de evaluación de capital. Por último, el informe apoyaba aumentar la importancia relativa de las medidas cualitativas frente a las puramente cuantitativas, empujando al conjunto del sistema a dejar de girar exclusivamente en torno al capital y reforzando los lazos entre esas medidas cualitativas y las notas del proceso.

El informe fue acogido positivamente por quien entonces presidía el Consejo de Supervisión del Mecanismo Único, el italiano Andrea Enria, quien declaró que las aportaciones de los expertos reforzaban la convicción de que la supervisión debía ser más adaptable, intrusiva y centrada en el riesgo\footnote{%
    Enria ejerció la presidencia del Consejo de Supervisión desde el 1 de enero de 2019 hasta el 31 de diciembre de 2023. Su sucesora, Claudia Buch, anteriormente vicepresidenta del Bundesbank, fue nombrada por el Consejo de la Unión Europea el 19 de octubre de 2023 y asumió el cargo el 1 de enero de 2024 para un mandato de cinco años. Buch compareció ante la Comisión de Asuntos Económicos y Monetarios del Parlamento Europeo el 18 de marzo de 2026 en el marco de la rendición de cuentas periódica del Mecanismo Único de Supervisión ante dicha institución.
}. Contrariamente a lo que podría sugerir una lectura desactualizada de este episodio, la implementación de las recomendaciones de los expertos no ha quedado pendiente de una futura presidencia, sino que ha sido efectivamente acometida desde 2024: un primer conjunto de medidas de reforma del proceso de revisión se decidió en 2024 y se encuentra ya sustancialmente implantado, con plena aplicación prevista para 2026, y en 2025 se inició además un segundo conjunto de reformas bajo el nombre de proyecto de supervisión de nueva generación, sustentado a su vez en una iniciativa de cultura supervisora del Mecanismo Único lanzada en febrero de 2025 con un plan de activación para el bienio 2025-2026\footnote{%
    Los resultados agregados del ejercicio 2025 del Proceso de Revisión y Evaluación Supervisora, publicados en noviembre de ese año, confirman beneficios ya tangibles de la reforma: las decisiones provisionales se notificaron a las entidades antes de finales de julio y las decisiones finales antes de finales de octubre, frente a plazos anteriores más dilatados; las decisiones son ahora más breves y se centran en los requisitos, el resultado principal y las preocupaciones supervisoras clave, remitiendo el detalle adicional a un anexo; se ha reducido el número de medidas cualitativas impuestas, concentrando la escalada de exigencias en las deficiencias más materiales; y para entidades con perfil de riesgo estable y sin necesidad de nuevos requisitos, el resultado del proceso se transmite mediante una carta operativa en lugar de una decisión formal completa, actualizándose de forma bienal en vez de anual.
}.

En definitiva, el pilar supervisor de la Unión Bancaria ha funcionado de manera muy exitosa y ha sido, hasta la fecha, el eslabón que mejor ha operado dentro de toda la arquitectura de la Unión Bancaria, hasta el punto de haber permitido que el sistema bancario europeo resistiera con holgura las turbulencias financieras desatadas en marzo de 2023 en Estados Unidos y Suiza. A pesar de este buen funcionamiento, el BCE ha optado por seguir adaptando el mecanismo a los nuevos tiempos, como demuestra la doble oleada de reformas de 2024 y 2025 que se acaba de describir, en lugar de conformarse con un diseño ya considerado satisfactorio.

\subsection{Pilar 2: Mecanismo Único de Resolución}
El Mecanismo Único de Resolución es el sistema centralizado de gestión de crisis bancarias en los Estados de la Unión Bancaria, es decir, el segundo pilar de su arquitectura institucional, encargado de decidir y ejecutar la resolución ordenada de una entidad inviable cuando concurre un interés público en optar por ese cauce en lugar de la liquidación concursal ordinaria de cada Estado.

Su fundamento jurídico se encuentra en el Reglamento (UE) n.º 806/2014, adoptado mediante procedimiento legislativo ordinario sobre la base del art. $114$ del TFUE, a diferencia del Mecanismo Único de Supervisión, y desarrolla a nivel institucional el contenido sustantivo que ya establece la Directiva de Reestructuración y Resolución Bancaria que trabajamos al hablar del Código Normativo Único: lo que el Mecanismo Único de Resolución añade no es una norma sustantiva nueva sobre qué es la resolución o cómo se absorben las pérdidas, sino quién decide y ejecuta esa resolución para las entidades de la Unión Bancaria. 

El mecanismo está compuesto por tres elementos que conviene distinguir desde el principio porque cada uno cumple una función distinta dentro del conjunto: la Junta Única de Resolución, que es la autoridad de resolución a nivel europeo y el órgano decisorio central del sistema; el Fondo Único de Resolución, que es el instrumento financiero que sostiene la ejecución de las decisiones de la Junta; y las Autoridades Nacionales de Resolución de cada Estado participante, que ejecutan sobre el terreno buena parte de las decisiones adoptadas a nivel central y conservan competencia directa sobre las entidades de menor tamaño relativo, en una lógica de reparto de intensidad que, al igual en el pilar supervisor, aquí se reproduce con matices propios.\footnote{%
    Está previsto además que, una vez se ratifique la reforma del Tratado del Mecanismo Europeo de Estabilidad, el Fondo Único de Resolución cuente con un respaldo fiscal común o \textit{common backstop}. A fecha de redacción de esta actualización, julio de 2026, esa ratificación sigue pendiente: el Parlamento italiano rechazó formalmente la ratificación del Tratado reformado el 21 de diciembre de 2023, y pese a la presión reiterada de otras instituciones europeas (entre ellas el vicepresidente del BCE, Luis de Guindos, que en octubre de 2025 advirtió públicamente de que no se puede esperar a la próxima crisis para cerrar esta laguna) Italia continúa siendo, a día de hoy, el único de los veinte miembros del MEDE que no ha ratificado la reforma, lo que mantiene bloqueada su entrada en vigor. Constituye por tanto una laguna estructural persistente de la Unión Bancaria, que ya ha cumplido más de una década de existencia sin que su segundo pilar esté completo en este extremo.
} 

La misión principal del mecanismo es asegurar una resolución ordenada de un banco inviable con el mínimo impacto posible sobre la economía real y las finanzas públicas de los Estados miembros de la Unión Bancaria y del resto de países afectados, asegurando la estabilidad financiera del conjunto, y aunque está concebido para actuar en épocas de crisis o ante situaciones delimitadas a entidades individuales, su actividad exige una preparación continua y anticipada (planes de resolución).

\subsubsection{Composición institucional y órganos de decisión}
La Junta Única de Resolución es una agencia europea dotada de personalidad jurídica propia, y su estructura de gobierno se articula en dos formaciones distintas que conviene no confundir porque tienen composición, competencias y lógica de funcionamiento diferentes. 

La Sesión Plenaria reúne a todas las Autoridades Nacionales de Resolución de los países participantes en la Unión Bancaria junto con los seis miembros permanentes de la Junta, y cuenta como observadores permanentes, sin derecho de voto, a la Comisión Europea y al Banco Central Europeo, en una lógica de participación institucional cruzada, donde el Banco Central Europeo nunca actúa en un vacío institucional sino siempre en cooperación o supervisión mutua con otras instituciones europeas. Las competencias de esta Sesión Plenaria son de naturaleza fundamentalmente presupuestaria y de gobierno general del sistema: la aprobación anual del presupuesto, del programa de trabajo y de las cuentas de la Junta; la aprobación de un esquema de resolución particular cuando este exige recurrir a fondos del Fondo Único de Resolución por un importe superior a cinco mil millones de euros, o cuando dentro de un periodo de doce meses varios esquemas de resolución sucesivos hayan superado en conjunto ese mismo importe; la autorización de la recaudación de contribuciones extraordinarias ex post a las entidades, de préstamos entre los distintos compartimentos nacionales del Fondo, y de vías de financiación alternativa con terceros; la aprobación de la política de inversiones del propio Fondo; y la aprobación del marco general de cooperación entre la Junta y las Autoridades Nacionales de Resolución.

La Sesión Ejecutiva, en cambio, reúne únicamente a los seis miembros permanentes de la Junta, junto con los mismos observadores permanentes ya mencionados, y es el órgano que ejerce la actividad decisoria ordinaria y más frecuente del sistema: aprueba los planes de resolución de las entidades bajo su responsabilidad directa, incluyendo la fijación de su requisito mínimo de fondos propios y pasivos admisibles, y aprueba los esquemas de resolución concretos cuando una entidad deviene inviable, esquemas que en todo caso requieren el respaldo posterior de la Comisión Europea y, en su caso, la aprobación del Consejo, lo que introduce un elemento de control político sobre una decisión que en su origen es técnica, de manera estructuralmente paralela al papel que ya vimos que desempeña la Comisión Europea sobre las normas técnicas de la Autoridad Bancaria Europea en el segundo nivel del Código Normativo Único. 

Existe además una modalidad reforzada, la Sesión Ejecutiva Extendida, que se activa para determinadas decisiones y en la que participan también las Autoridades Nacionales de Resolución de los Estados afectados; cuando en esa formación extendida no se alcanza un consenso entre todos los participantes, el reglamento prevé que la decisión se adopte finalmente solo por los miembros permanentes de la Junta, mediante mayoría simple entre ellos, de modo que la participación nacional amplía la deliberación pero no puede en último término bloquear ni diluir la decisión, que sigue perteneciendo a la instancia centralizada.\footnote{%
    La Junta Única de Resolución opera desde 2015. Está presidida por Dominique Laboureix, anteriormente secretario general de la autoridad francesa de supervisión y resolución, la ACPR, nombrado por el Consejo de la Unión Europea el 25 de noviembre de 2022 y que tomó posesión de su cargo en enero de 2023. Laboureix es, desde marzo de 2025, además presidente del Grupo Directivo de Resolución del Consejo de Estabilidad Financiera a nivel internacional, lo que da al Mecanismo Único de Resolución una proyección añadida en la coordinación de estándares de resolución fuera de la propia Unión.
}

\subsubsection{Funciones del mecanismo y relación entre la Junta y las Autoridades Nacionales de Resolución}
Las tareas del MUR son muy variadas. En primer lugar, recae sobre éste la elaboración de los planes de resolución de todos los bancos de la Unión Bancaria, sin excepción alguna en cuanto al deber de tenerlos, aunque el grado de detalle e intensidad de revisión varíe según el tamaño de la entidad. También el establecimiento del nivel mínimo de pasivos elegibles de cada entidad, es decir, del requisito conocido como MREL, que fija cuánto pasivo susceptible de absorber pérdidas debe mantener cada banco para que su resolución sea viable sin recurrir a fondos públicos. Por otro lado, la valoración de la resolubilidad efectiva de cada entidad y la identificación y eliminación de los obstáculos que puedan impedirla en la práctica, la aprobación y ejecución de esquemas de resolución concretos cuando una entidad deviene inviable y la gestión del propio Fondo de Resolución. 

La distribución de estas tareas entre la Junta Única de Resolución y las Autoridades Nacionales de Resolución reproduce, con sus propios matices, la misma lógica del Mecanismo Único de Supervisión: la Junta asume competencia de resolución directa, incluida la aprobación de planes y esquemas de resolución, sobre las entidades significativas y sobre los grupos bancarios con actividad transfronteriza relevante dentro de la Unión Bancaria, mientras que las Autoridades Nacionales de Resolución conservan competencia directa sobre el resto de entidades, si bien siempre dentro del marco normativo fijado por la Junta y bajo su supervisión última, ya que es la propia Junta quien aprueba también los planes de resolución de estas entidades de menor tamaño y puede intervenir si lo considera necesario para garantizar la coherencia del sistema en su conjunto.

Hasta la aprobación del paquete de revisión del marco de gestión de crisis y garantía de depósitos que modifica simultáneamente la Directiva de Reestructuración y Resolución Bancaria, el propio Reglamento del Mecanismo Único de Resolución y la Directiva de Sistemas de Garantía de Depósitos, la distinción práctica entre entidades grandes y entidades pequeñas o medianas no residía tanto en la titularidad formal de la competencia de resolución como en el hecho de que el filtro del interés público excluía casi sistemáticamente a estas últimas del procedimiento de resolución propiamente dicho, remitiéndolas a procedimientos concursales nacionales dispares entre sí. La reforma de 2026 amplía el uso efectivo de la resolución precisamente para estas entidades de menor tamaño, recalibrando los criterios de valoración del interés público y ampliando el papel de los fondos de garantía de depósitos como fuente adicional de financiación de la resolución cuando el nivel mínimo de pasivos elegibles de una entidad pequeña resulta insuficiente por sí solo, con una fecha general de aplicación fijada en el 11 de mayo de 2028 para las disposiciones del Reglamento y un plazo de transposición de veinticuatro meses para las dos directivas del paquete.\footnote{%
    Esta reforma no ha alterado la arquitectura institucional que se acaba de describir sino el umbral práctico de activación del procedimiento de resolución frente a la liquidación concursal ordinaria, y el papel de los fondos de garantía de depósitos dentro de la financiación de ese procedimiento.
}

\subsubsection{El Fondo Único de Resolución: financiación, usos y respaldo fiscal común}
El Fondo Único de Resolución tiene por objeto asegurar la eficiente aplicación de las herramientas de resolución y facilitar la ejecución de las medidas acordadas por la Junta Única de Resolución. Como ya se ha explicado, la crisis de un banco debe financiarse en primer lugar con cargo a sus propios accionistas y acreedores, mediante el mecanismo de absorción interna de pérdidas, y solo si esto resulta insuficiente puede recurrirse al Fondo, que a su vez está financiado enteramente por las propias entidades financieras de la industria bancaria, nunca directamente por el contribuyente; únicamente de manera excepcional, y tras haber agotado el uso del Fondo, podría llegar a recurrirse a fondos públicos nacionales, en última instancia y bajo condiciones estrictas.

El objetivo estatutario del Fondo, fijado por el art. $69$ del Reglamento del MUR, es alcanzar al menos el uno por ciento del importe de los depósitos cubiertos de todas las entidades de crédito autorizadas en los Estados participantes, un umbral que debía cumplirse al término de un periodo inicial de acumulación de ocho años, comprendido entre 2016 y 2023. La Junta confirmó que el Fondo alcanzó ese nivel objetivo a 31 de diciembre de 2023, con un importe disponible de setenta y ocho mil millones de euros frente a un umbral calculado en torno a los setenta y cinco mil millones sobre la base de unos depósitos cubiertos estimados entonces en siete billones y medio de euros en el conjunto de la Unión Bancaria.\footnote{%
    El Fondo se situó en ochenta mil millones de euros a 31 de diciembre de 2024 y en algo más de ochenta y uno mil millones a 31 de diciembre de 2025, superando en ambos casos el umbral estatutario, motivo por el cual la Junta ha confirmado que, por tercer año consecutivo, las entidades no tendrán que realizar contribuciones ordinarias ex ante al Fondo para 2026, ya que el nivel objetivo se sigue considerando alcanzado. La siguiente verificación anual está prevista para comienzos de 2027.
} Se financia mediante contribuciones ex ante de las entidades de crédito y de determinadas empresas de servicios de inversión, calculadas conforme a criterios de riesgo individual de cada entidad; en el supuesto de que uno o varios casos de resolución agoten los recursos disponibles y sea necesaria financiación adicional, se recaudarían contribuciones ex post a cargo de las mismas entidades, y dado que estas últimas normalmente no están disponibles ni son accesibles de forma inmediata en el momento en que se necesitan, el Fondo debería recurrir en ese escenario a operaciones de endeudamiento, ya sean privadas o públicas, para cubrir el desfase temporal.

Solo en el supuesto de que ni las contribuciones ex post ni el endeudamiento puente resultaran suficientes se activaría el respaldo fiscal común o \textit{common backstop}, que en el caso de la zona euro consistiría en una línea de crédito proporcionada por el MEDE, diseñada para duplicar aproximadamente la capacidad efectiva del Fondo. Como ya se ha explicado, este respaldo sigue sin estar disponible porque continúa pendiente la ratificación del Tratado reformado del MEDE por parte del Parlamento italiano, de modo que, pese a que el Fondo se encuentra hoy correctamente capitalizado según sus propios objetivos estatutarios, la Unión Bancaria continúa operando, más de una década después de su creación, sin esta pieza adicional de respaldo público mutualizado que estaba prevista desde el diseño original del sistema.

Entre los posibles usos del Fondo se cuentan la concesión de garantías sobre activos y pasivos de una entidad tras su resolución; la concesión de préstamos a una entidad en resolución, a sus filiales, a un banco puente o a una sociedad de gestión de activos; la compra de cualquier tipo de activos de la entidad en resolución; contribuciones, en sentido amplio, a un banco puente y a una sociedad de gestión de activos; compensaciones a accionistas y acreedores que hayan soportado en resolución más pérdidas de las que habrían soportado en una liquidación concursal ordinaria, en aplicación directa del principio de que ningún acreedor puede salir peor parado que en concurso, que ya trabajamos al hablar de la Directiva de Reestructuración y Resolución Bancaria; y contribuciones a la absorción de pérdidas y a la recapitalización de una entidad, sustituyendo a determinados acreedores que hayan sido excluidos de un procedimiento de recapitalización interna. 

Para hacer operativo el Fondo en su fase de transición hacia la plena mutualización, los Estados miembros firmaron un Acuerdo Intergubernamental, instrumento jurídico de derecho internacional público y no de derecho de la Unión, cuya función es regular la transferencia progresiva de las contribuciones nacionales recaudadas hacia los compartimentos mutualizados del propio Fondo\footnote{%
    El recurso a un acuerdo intergubernamental, y no a un reglamento de la Unión, para regular este aspecto concreto de la mutualización financiera, responde a las mismas cautelas políticas y constitucionales que explican por qué el respaldo fiscal común depende de una reforma del Tratado del Mecanismo Europeo de Estabilidad, y no de legislación ordinaria de la Unión: la mutualización de recursos financieros entre Estados miembros para el rescate de entidades privadas ha sido, desde el origen mismo de la crisis del euro, la dimensión políticamente más sensible de toda la Unión Bancaria, mucho más que la centralización de la supervisión o de la decisión de resolución en sí mismas.
}.









Hasta la fecha, sólo ha habido dos casos de resolución: el del Banco Popular Español y el de las filiales eslovena y croata de Sberbank (se decidió que no era necesario resolver la matriz austríaca de Sberbank). Pero ello no es debido a que no hubiera más entidades en crisis. Al contrario, otras entidades han sido declaradas inviables, en particular, los bancos letones ABLV y PNB Banka y los bancos italianos Banca Popolare di Vicenza y Veneto Banca. Sin embargo, al considerar que no concurría interés público en estos casos, estas entidades fueron liquidadas por las autoridades nacionales con arreglo a sus respectivos marcos jurídicos nacionales, y en algunos casos con un uso de los fondos del contribuyente más generoso, en lugar de resueltas por la europea Junta Única de Resolución sobre la base de la normativa europea común. Como reacción a estos casos, la reforma buscaba ensanchar el concepto de interés público y aumenta la carga de la prueba sobre la autoridad de resolución para demostrar que no concurre interés público.

Desde el lanzamiento del segundo pilar de la Unión Bancaria, se han registrado también casos de entidades que han recibido fondos públicos sin ni siquiera haber sido declaradas inviables o con probabilidad de serlo (FOLTF, por sus siglas en inglés). Se trata, por ejemplo, del banco italiano Monte dei Paschi di Siena, beneficiario de la denominada recapitalización precautoria. La propuesta de CMDI busca restringir el uso del instrumento de recapitalización precautoria a casos en los que la entidad sea claramente solvente y que el apoyo sea estrictamente temporal, aumentando los requisitos para determinar desde el inicio la duración del apoyo y la estrategia de salida.



En los momentos de redactar este tema, los colegisladores (Consejo de la UE y Parlamento Europeo) se encuentran debatiendo sobre la base de la propuesta normativa de la Comisión Europea. Con carácter general, la propuesta de la Comisión es positiva y adecuada, tratando de actualizar el marco jurídico sobre la base de las experiencias acumuladas en los últimos años. No obstante, el sector bancario europeo necesita ir un paso más allá y contar, por fin, con un Fondo de Garantía de Depósitos Único, cuya última fase, !a de seguro pleno, debería comenzar en 2024 de acuerdo con la propuesta normativa de la Comisión Europea de 20 15, tal y como se explica en una sección posterior. Y además, es necesario diseñar un buen sistema de provisión de liquidez en resolución, como se detalla en la siguiente sección. Hasta entonces, aunque se formulen propuestas de mejora del marco de gestión de crisis bancarias, no se estará completando la Unión Bancaria, respondiendo a lo que el sector bancario realmente necesita.

ll.l/. VI. la provisión de liquidez en resolució11

La resolución del Banco Popular en 20 1 7 ha dejado claro que la liquidez en resolución, incluyendo la disponibilidad de colateral suficiente, es un tema prioritario. En la resolución de un gran banco, la credibilidad (en cuanto al tamaño y la rapidez) en el suministro de liquidez solo es posible con la intervención de un banco central. El Fondo Único de Resolución (SRF, por sus siglas en inglés), incluso con un backstop operativo (más adelante se explica en qué consiste el backstop), podría ser suficiente para afrontar la crisis de un banco mediano o pequeño, pero su capacidad limitada es un hándicap a la hora de gestionar la crisis de un banco sistémico. En algunos países, como Estados Unidos y el Reino Unido, el banco central suministra la liquidez en resolución con una garantía pública. Esta es la alternativa que debería ser explorada en la- Unión Bancaria. El SRB está siguiendo el nivel de activos comprometidos (encwnbered) para valorar la capacidad de los bancos para acceder a financiación en el mercado o en el banco central. Asimismo, el SRB está cooperando con el BCE en el diseño de escenarios de estrés de liquidez, para poder estimar necesidades potenciales de liquidez y planificar acciones.



\subsection{Pilar 3: ¿Esquema Europeo de Garantía de Depósitos (EDIS)?}
El Esquema Europeo de Garantía de Depósitos es el tercer pilar, todavía no implementado, de la arquitectura institucional de la Unión Bancaria, concebido para mutualizar a nivel europeo la financiación y, en su fase final, también la cobertura de pérdidas de los fondos nacionales de garantía de depósitos (FNGD, contemplados en la Directiva de Sistemas de Garantía de Depósitos), aplicable a los veintisiete Estados miembros. 

\subsubsection{Arquitectura por fases: Propuesta de 2015}
En noviembre de 2015 la Comisión Europea presentó su propuesta normativa, formalmente un reglamento que modificaría el propio Reglamento del MUR para incorporar estas nuevas funciones a la Junta Única de Resolución, con el objetivo de completar así la arquitectura de tres pilares que ya habían anticipado tanto el Informe de los Cuatro Presidentes de 2012 como el Informe de los Cinco Presidentes de 2015. La propuesta articulaba una implantación progresiva en tres etapas sucesivas, durante las cuales el esquema europeo cubriría las necesidades de financiación que pudieran tener los FNGD, tanto para compensar a los depositantes como para contribuir a un procedimiento de resolución, así como, en niveles crecientes, las pérdidas en que definitivamente pudiera incurrir el fondo nacional.

\paragraph{Primera Fase: Reaseguro}
La primera fase, denominada de reaseguro y prevista con una duración de tres años, permitiría a los FNGD obtener financiación del esquema europeo en aquellos casos en los que presentaran falta de liquidez, ya fuera por un evento de desembolso a depositantes o por tener que contribuir a un procedimiento de resolución, con un límite del $20\%$ de los fondos europeos disponibles o diez veces los recursos del propio fondo nacional, \textbf{lo que resultara menor}. 

Un evento de desembolso generador de falta de liquidez se producía cuando la cuantía de los depósitos cubiertos de la entidad inviable superaba tanto el importe de recursos financieros que el fondo nacional debía mantener conforme a la Directiva de Sistemas de Garantía de Depósitos como el importe adicional que ese fondo pudiera recaudar mediante aportaciones extraordinarias ex post en el plazo de tres días desde el propio evento; y de manera paralela, la falta de liquidez por contribución a una resolución se producía cuando la cuantía exigida al fondo nacional para esa contribución superaba sus recursos financieros disponibles en ese momento. En esta fase de reaseguro, además de la financiación puramente transitoria, el esquema europeo asumiría también una cobertura directa del exceso de pérdidas del fondo nacional participante equivalente al veinte por ciento de ese exceso, y de manera simultánea se iniciaría un traspaso progresivo de los fondos ya acumulados por los sistemas nacionales hacia el fondo europeo, fijado en un veinte por ciento del total de recursos nacionales durante esta primera etapa.

\paragraph{Segunda Fase: Coaseguro}
La segunda fase, de coaseguro y con una duración esperada de cuatro años adicionales, elevaría sustancialmente el grado de mutualización de riesgos: dejaría de ser necesario que el fondo nacional agotara primero sus propios recursos antes de acudir al esquema europeo, de modo que ambos, fondo nacional y esquema europeo, se responsabilizarían conjuntamente y desde el primer momento de los costes derivados de un evento de desembolso o de una contribución a resolución. 

La participación proporcional del esquema europeo en esos costes aumentaría de manera progresiva año a año dentro de esta fase, comenzando en el treinta y seis por ciento durante el primer año y alcanzando el $84\%$ en el cuarto y último año de la etapa.

\paragraph{Tercera Fase: Mutualización plena}
La tercera fase, denominada de seguro pleno o de mutualización plena, pretendía alcanzarse en el año 2024.\footnote{%
    Este calendario, formulado en la propuesta original de 2015, no se ha cumplido ni remotamente: a fecha de esta actualización, julio de 2026, la propuesta legislativa continúa clasificada en el registro de seguimiento legislativo del Parlamento Europeo con el estatus de \textit{bloqueada}, categoría que se aplica a aquellos expedientes en los que no se ha producido ningún avance sustancial durante al menos los nueve meses previos a la última revisión de dicho registro, en enero de 2026.
} 

En esta etapa, el esquema europeo ofrecería a los fondos nacionales la financiación íntegra de cualquier necesidad de liquidez y cubriría la totalidad de las pérdidas derivadas de un evento de desembolso o de una contribución a resolución, funcionando con el mismo mecanismo operativo que la fase de coaseguro pero con una participación del $100\%$ en lugar de una participación parcial creciente.

\subsubsection{¿Cómo se estructuraría? Financiación y gestión institucional}
La financiación del esquema europeo, tal como se diseñó en la propuesta, se realizaría mediante aportaciones ex ante de los bancos participantes, hasta alcanzar un nivel de recursos equivalente al $0,8\%$ de los depósitos garantizados del conjunto de la Unión Bancaria, exactamente el mismo porcentaje objetivo que el fijado para los FNGD considerados individualmente. 

La gestión del esquema correspondería a la \textbf{Junta Única de Resolución}, que de este modo vería ampliadas sus funciones más allá de la resolución bancaria propiamente dicha para asumir también la administración de este fondo mutualizado de garantía de depósitos, sin que ello supusiera la desaparición de los FNGD ni de sus autoridades competentes, cuyas funciones se adaptarían para coexistir con el nuevo nivel europeo en lugar de ser sustituidas por él.

\subsubsection{Estado de las negociaciones}
Como decíamos, las negociaciones se encuentran totalmente bloqueadas en el seno del Consejo de la Unión Europea. La oposición principal sigue procediendo de Alemania, respaldada por otros Estados del norte de la Unión, cuya reticencia de fondo es la que ya articuló públicamente el Ministerio de Finanzas alemán desde el origen mismo del debate: la exigencia de que, antes de aceptar cualquier mutualización de riesgos de depósitos entre Estados, se complete primero una reducción efectiva del riesgo en los balances bancarios nacionales, en particular en lo relativo a la exposición de los bancos a la deuda soberana de su propio Estado y a los activos dudosos heredados de crisis anteriores, bajo la lógica de que compartir riesgos antes de haberlos reducido suficientemente equivaldría a que los contribuyentes de los Estados más solventes terminaran, de facto, subvencionando la fragilidad bancaria de otros. 

Francia, por su parte, tampoco ha mostrado gran apetito por avanzar en este esquema, aunque por una razón distinta: la fuerte concentración y el gran tamaño relativo de sus propias entidades de crédito hace que el diseño de una mutualización plena no le resulte tan claramente beneficioso como a Estados con sistemas bancarios más fragmentados o de menor tamaño individual.\footnote{%
    Diversas propuestas alternativas de diseño han intentado desbloquear este estancamiento sin éxito hasta la fecha, entre ellas una propuesta franco-alemana que planteaba un esquema institucionalmente integrado pero financiado de manera diferenciada por país, vinculado además a la introducción de recargos de capital por concentración de deuda soberana en el balance de cada entidad, precisamente para abordar la preocupación alemana de fondo sin renunciar a la integración institucional. Ninguna de estas propuestas ha logrado hasta ahora superar el bloqueo político básico.
}

\paragraph{Reforma del Marco de Gestión de Crisis: Sustituto parcial}
En este contexto, la Comisión ha optado por una reforma parcial del Marco de Gestión de Crisis que ya introujimos brevemente en el apartado anterior. En abril de 2026, se reforma el Reglamento del MUR, la Directiva de Reestructuración y Resolución y la Directiva sobre Esquemas de Garantía de Depósitos.\footnote{%
    El 20 de abril de 2026 publicó en el DOUE el paquete legislativo definitivo de reforma del marco de gestión de crisis y garantía de depósitos, formado por el Reglamento (UE) 2026/808, que modifica el Reglamento del Mecanismo Único de Resolución, la Directiva (UE) 2026/806, que modifica la Directiva de Reestructuración y Resolución Bancaria, y la Directiva (UE) 2026/804, que modifica específicamente la Directiva de Sistemas de Garantía de Depósitos que aquí más interesa.
}

Conviene subrayar con toda claridad que esta reforma no es el Esquema Europeo de Garantía de Depósitos ni lo sustituye institucionalmente (ni crea ningún fondo mutualizado ni avanza en esa dirección), lo que hace es una intervención mucho más modesta: reforzar y flexibilizar el uso de los FNGD, cada uno de ellos financiado y respaldado únicamente por su propio Estado, dentro de los procedimientos de resolución y liquidación, precisamente porque estos FNGD no habían podido hasta ahora desplegar todo su potencial como herramienta de gestión de crisis.

El razonamiento que sostiene esta reforma parte de dos problemas prácticos identificados en la aplicación del marco de 2014. Por un lado, en un procedimiento de liquidación concursal ordinaria, el proceso de transferencia de los activos de la entidad inviable hacia un adquirente puede resultar excesivamente largo, lo que se traduce en un deterioro material del valor de esos activos y, en consecuencia, en una pérdida añadida para los acreedores de la entidad; permitir que los recursos del FNGD faciliten económicamente esa transferencia acelera el proceso y reduce ese deterioro, en beneficio del conjunto del sistema financiero, si bien la reforma final condiciona expresamente este uso a que el coste de la contribución del fondo sea inferior al que habría supuesto un pago directo a los depositantes cubiertos, de modo que la operación deba resultar más barata que la alternativa de desembolso directo. Por otro lado, la reforma aclara y refuerza que las aportaciones del fondo nacional de garantía de depósitos cuentan a efectos del umbral mínimo de recapitalización interna previa del $8\%$ de los pasivos totales, umbral cuyo cumplimiento es condición para poder acceder a la financiación del Fondo Único de Resolución dentro de la Unión Bancaria, facilitando así indirectamente que entidades de menor tamaño, para las que alcanzar ese umbral únicamente con sus propios accionistas y acreedores resulta más difícil, puedan también beneficiarse de la financiación centralizada del segundo pilar. Esta es, de hecho, la misma ampliación del papel de los FNGD como fuente de financiación de la resolución que mencionamos en el MUR, y que queda plenamente contextualizada desde el lado del tercer pilar.

En cuanto al alcance de la protección al depositante, la reforma mantiene en cien mil euros por titular y por entidad adherida la cuantía protegida por los fondos nacionales de garantía de depósitos, sin alterar este umbral que ya conoces del Código Normativo Único. Lo que sí amplía es el ámbito subjetivo de esa protección, para cubrir también a determinadas entidades públicas, como hospitales, centros educativos y ayuntamientos, así como el dinero de clientes depositado en determinados tipos de fondos de clientes, entre ellos los de sociedades de inversión y de entidades de pago y de dinero electrónico. Con esta ampliación limitada, y no con una cobertura general de todos los depósitos no garantizados, el legislador europeo ha optado por mantener el sistema existente antes de las turbulencias financieras de marzo de 2023 en Estados Unidos, cuando las autoridades estadounidenses activaron la denominada excepción sistémica para que el Fondo de Garantía de Depósitos Federal cubriera la totalidad de los depósitos, incluidos los no garantizados por encima del umbral ordinario, de Silicon Valley Bank y Signature Bank; la Unión Europea, a diferencia de esa respuesta estadounidense ad hoc, ha preferido ampliar de forma tasada y permanente el círculo de sujetos protegidos antes que introducir un mecanismo excepcional de cobertura ilimitada activable discrecionalmente ante una crisis sistémica futura.\footnote{%
    Esta decisión puede entenderse como una toma de postura implícita sobre el debate de fondo que sigue bloqueando el propio Esquema Europeo de Garantía de Depósitos: mientras que la respuesta estadounidense de 2023 mostró la disposición a mutualizar riesgo de manera amplia y rápida ante una crisis concreta, la Unión ha optado, también en esta reforma de 2026, por una vía de ampliación gradual y predefinida por ley, coherente con la misma cautela alemana hacia la mutualización de riesgos que explica el bloqueo del propio esquema europeo.
}








\clearpage
\section{Unión de Mercado de Capitales}
\puntosclave{
    La Unión de Mercado de Capitales está compuesta de un conjunto muy elevado y disperso de piezas normativas sobre temas muy dispares, a diferencia de la Unión Bancaria (que sí cuenta con una estructura muy definida). Se recomienda por tanto centrar la atención en las iniciativas más relevantes de la Unión de los Mercados de Capitales, y priorizar una visión ordenada, de conjunto, y entendiendo la motivación para cada norma y su relación con las demás: el seguimiento actualizado de todas es complejo y probablemente innecesario en el marco de la preparación de las oposiciones a TCEE.

    A diferencia de la Unión Bancaria, la CMU es un proyecto con una fuerte dispersión nonnativa y que ha ido evolucionando a lo largo de los años. 
    
    Tras el lanzamiento de la iniciativa ,CMU en 2015, se han elaborado varios planes de acción, que incluyen tanto medidas legislativas como no legislativas, y la Comisión ha adoptado varios paquetes de medidas legislativas: (i) el plan de acción CMU inicial publicado en septiembre de 2015; (ii) la revisión ya ampliación el plan de acción en junio de 2017, (iii) el plan de acción de septiembre de 2020, (iv) cuatro iniciativas legislativas presentadas por la Comisión Europea en noviembre de 2021; y, (v) otras tres iniciativas legislativas presentadas por la Comisión Europea en diciembre de 2022. 
    
    El marco revisado de titulización entró en vigor en 2019, con el objetivo principal deestablecer un mercado de tituliz,ación de la UE que ayude a financiar la economía sin crear riesgos para la estabilidad financiera. El marco se compone de (1) el Reglamento sobre titulización, que crea una categorja de productos de titulización simples, transparentes y normalizados (STS), (2) una modificación de la CRR, que hace que el tratamiento del capital de las titulizaciones para los bancos sea más sensible al riesgo e introduce un tratamiento preferencial específico para las titulizaciones STS y (3) una modificación de las normas prudenciales relativas a las inversiones de las aseguradoras en titulizaciones establecidas en la Directiva Solvencia. 
    
    En 2017, se aprobó una nueva Directiva de Folletos, que sustituye a la anterior del año 2003, con el objetivo de simplificar las obligaciones administrativas para la publicación de folletos, pero de un modo que garantice que los inversores siguen estando correctamente informados y protegidos. 
    
    La Directiva y el Reglamento de bonos garantizados buscan establecer unos requisitos mínimos de armonización que deberán cumplir todos los bonos garantizados comercializados en la UE. Esto aumentará la seguridad para los inversores y abrirá nuevas oportunidades, en particular allí donde los mercados están menos desarrollados. 
    
    El Reglamento sobre proveedores europeos de servicios de financiación participativa (crowdfuding) establece normas uniformes en toda la UE para la prestación de servicios de financiaci(m participativa basados en inversiones y préstamos relacionados con la financiación empresarial. Permite a las plataformas solicitar un pasaporte de la UE basado en un conjunto único de normas, lo que les facilita ofrecer sus servicios en toda la UE con una única autorización. 
    
    El producto paneuropeo de pensiones personales (PEPP) es un plan voluntario de pensiones personales que complementará los sistemas públicos y profesionales de pensiones existentes, así como los planes nacionales de pensiones privadas. Los PEPP están regulados por un Reglamento de 20 19. Estas normas tienen como objetivo dar a los ahorradores más opciones y ofrecerles productos de pensiones personales más competitivos a la hora de ahorrar para la jubilación, al tiempo que disfrutan de una sólida protección de los consumido res. 
    
    Hasta el 25 de junio de 2021, las normas prudenciales aplicables a las empresas de inversión formaban parte de la CRR y la CRD. El 26 de junio de 202 1 , la mayoría de las empresas de inversión pasaron a estar sujetas a un nuevo marco prudencial, compuesto por el Reglamento y la Directiva sobre las empresas de inversión (IFR e IFD, respectivamente). El IFR y el IFD se aplican a las empresas de inversión consideradas suficientemente pequeñas y no interconectadas (denominadas empresas de "clase 3 ") y a las empresas de inversión que no entran en ninguna de las otras categorías (denominadas empresas de "clase 2"). La gran mayoría de las empresas de inversión de la UE pertenecen a estas dos categorías. 
    
    El SME Growth Market, un nuevo tipo de sistema multilateral de negociación, fue introducido por la MiFID II para mejorar el acceso de las PYME a la financiación basada en el mercado. Este régimen se ha mejorado a través de varias iniciativas legislativas desde su creación, como el SME listing Act, adoptado en 2019, y el Listing Package, propuesto por la Comisión Europea y en negociación por los colegisladores de la UE en los momentos de redactar este tema. 
    
    En mayo de 2023 , los colegisladores de la UE alcanzaron un acuerdo para crear el ESAP. El acceso gratuito, fácil de usar, centralizado y digital a la información financiera y relacionada con la sostenibilidad hecha pública por las empresas europeas, incluidas las pequeñas empresas, facilitará el proceso de toma de decisiones a una amplia gama de inversores, incluidos los inversores minoristas. El ESAP no impone a las empresas europeas ningún requisito adicional de notificación de infonnación, ya que simplemente facilitará el acceso a la infonnación ya hecha pública en aplicación de las directivas y reglamentos europeos pertinentes.
}

Frente a la Unión Bancaria, cuya pertenencia es la consecuencia automática de una obligación de Tratado que solo Dinamarca puede eludir de forma permanente, la Unión del Mercado de Capitales es una estrategia política de la Comisión Europea (un conjunto de planes de acción sucesivos, revisados y reformulados con cada nueva Comisión desde 2015) que avanza exclusivamente instrumento a instrumento, mediante la base jurídica general de armonización del mercado interior. 

Es un hecho bien conocido que los mercados de capitales de la Unión en su conjunto siguen estando poco desarrollados en comparación con los de otras grandes economías, como EE.UU. o el Reino Unido.\footnote{%
    Avanzar en la CMU se ha vuelto más relevante si cabe desde la salida de Reino Unido de la Unión Europea. Sin embargo, el progreso realizado en la CMU es relativamente reducido: no tanto por razones políticas, como en el caso de la Unión Bancaria, sino por razones técnicas. A diferencia de la Unión Bancaria, la CMU es un proyecto con una fuerte dispersión normativa que ha ido evolucionando a lo largo de los años. Mientras que los tres pilares de la Unión Bancaria estaban claros desde el principio, las medidas adoptadas para la CMU han sido mucho menos claras y fruto de constantes revisiones.
} Aunque las diferencias estructurales hacen que los EE.UU. no puedan considerarse como referencia directa para la UE, el potencial de desarrollo de los de los mercados de capitales de la Unión sigue siendo significativo, especialmente en el espacio minorista y para la financiación de las PYME. Por ello, la Comisión presentó en 2015 su primera estrategia para la Unión del Mercado de Capitales.

Lo primero que debemos tener en cuenta es que, a partir de marzo de 2025, cuando la Comisión Europea publicó la comunicación \textit{Savings and Investments Union — A Strategy to Foster Citizens' Wealth and Economic Competitiveness in the EU}, sustituye formalmente la estrategia para la Unión del Mercado de Capitales por una estrategia rebautizada, la Unión del Ahorro y las Inversiones, que absorbe y reformula toda la agenda anterior. La sustitución no es meramente nominal: \textbf{el cambio de nombre incorpora un cambio real de énfasis}. Donde la Unión del Mercado de Capitales ponía el acento casi exclusivamente en la integración de los mercados de valores como fin en sí mismo, la Unión del Ahorro y las Inversiones parte de un diagnóstico más amplio y, en cierto modo, más urgente: los hogares europeos mantienen alrededor de diez billones de euros en depósitos bancarios, un ahorro seguro y accesible pero que rinde sistemáticamente menos que la inversión en mercados de capitales, mientras que la Unión necesita movilizar entre setecientos cincuenta y ochocientos mil millones de euros adicionales al año hasta 2030 para financiar sus prioridades estratégicas (competitividad, transición ecológica y digital, y, de forma creciente desde 2025, defensa).\footnote{%
    Esta reformulación no surgió de la nada, sino que responde a tres informes encargados o impulsados por la propia Comisión durante 2024. El primero, presentado por LETTA el 17 de abril de 2024 bajo el título \textit{Much More than a Market}, fue el que propuso explícitamente, por primera vez, la creación de una Unión del Ahorro y las Inversiones como transformación de la Unión del Mercado de Capitales, señalando que esta última, diez años después de su lanzamiento, no había logrado completarse. El segundo, publicado por NOYER el 25 de abril de 2024 bajo el título \textit{Developing European capital markets to finance the future}, se centró específicamente en las barreras a la profundidad de los mercados de capitales europeos. El tercero, y el más citado políticamente, es el Informe DRAGHI , \textit{The future of European competitiveness}, que cuantificó precisamente esa necesidad de inversión adicional de setecientos cincuenta a ochocientos mil millones de euros anuales y señaló el infradesarrollo de los sistemas de pensiones de segundo y tercer pilar en la mayoría de los Estados miembros como una oportunidad perdida concreta para movilizar ahorro hacia inversión productiva a largo plazo, el mismo diagnóstico que, como ya viste, subyace al fracaso comercial del PEPP.
} 

No obstante, y a diferencia de la Unión Bancaria, el catálogo de instrumentos de este proyecto ha ido cambiando con cada revisión y constituye un ámbito de acción muy heterogeneo. 

Para simplificar este hecho, y en aras de estructurarlo de manera coherente, consideraremos la estrategia de la Comisión (articulada en torno a cuatro líneas de actuación), como una estrategia que incide en dos ámbitos distintos: el mercado y sus participantes.\footnote{%
    La estrategia original de la Comisión considera medidas dirigidas a los ciudadanos, a las empresas, a los propios mercados, y a la centralización de la supervisión, esta última la línea de trabajo más directamente heredera de la vieja ambición de integración institucional de la Unión del Mercado de Capitales.
}

\subsection{Historia: ¿Cómo hemos llegado a donde estamos?}
El primer hito histórico en la formación de la UMC fue la ya mentada iniciativa de 2015. Ésta establecía $33$ acciones que incluían propuestas normativas para relanzar la titulización y promover la inversión a largo plazo en infraestructuras, esto último mediante la reforma de la normativa de seguros Solvencia II. Asimismo, se anunciaba la próxima presentación de una propuesta normativa para simplificar los folletos de emisión y se lanzaban dos consultas públicas en relación con los bonos garantizados y los fondos de capital riesgo, preguntando en este último caso si tenía sentido introducir modificaciones específicas a la normativa reguladora de los fondos EuVECA (European Venture Capital Funds) y EuSEF (European Social Entrepreneurship Funds).

Desde entonces, se han elaborado varios planes de acción, incluyendo medidas legislativas y no legislativas. La primera de estas revisiones tuvo lugar en junio de 2017. La Comisión Europea actualizó y complementó el Plan de Acción de 2015, estableciendo una lista de nueve prioridades. Entre ellas, destacan el fortalecimiento de los poderes de la \textit{European Securities and Market Agency} (ESMA), el establecimiento de un marco regulatorio facilitador de la cotización en mercado de las PYME, la revisión del tratamiento prudencial de las firmas de inversión o medidas para favorecer los mercados secundarios de préstamos fallidos. Asimismo, se anunció que la Comisión Europea presentaría propuestas normativas sobre el producto paneuropeo de pensiones personales (PEPP), los bonos garantizados y el aumento de la seguridad jurídica en la tenencia de valores en un contexto transfronterizo. 

A pesar de la importante mejora del marco legislativo del mercado de capitales, la sensación general en 2020 entre los participantes era de lentitud y desaliento: aún quedaba mucho camino por recorrer y lo estábamos caminando, no corriendo. Ante este esceneraio, la Comisión creó un Foro de Alto Nivel, introduciendo a entidades de la sociedad civil y expertos, con el fin de presentar propuestas que pudieran relanzar la UMC. A raíz de este foro, surgieron un conjunto de medidas centradas en cuatro ámbitos de acción (financiación de las empresas, las infraestructuras de mercado, la inversión minorista y el mercado interior) que fueron integradas por la Comisión en el tercer Plan de Acción de la UMC (2020). Plan de Acción que contenía $16$ acciones (legislativas y no legislativas), entre las que destaca la creación de un punto de acceso único de información sobre empresas para inversores (\textit{European Single Access Point}, ESAP). Formalmente, la Comisión Europea presentó en 2021 cuatro iniciativas legislativas con el objetivo de avanzar en la cronología establecida por el Plan: (i) la creación del \textit{European Single Access Point} (ESAP); (ii) la revisión del \href{https://www.boe.es/buscar/doc.php?id=DOUE-L-2015-80967}{Reglamento (UE) 2015/760 del Parlamento Europeo y del Consejo, sobre los fondos de inversión a largo plazo europeos} (de ahora en adelante: Reglamento ELTIF; (iii) revisión de la \href{https://www.boe.es/doue/2011/174/L00001-00073.pdf}{Directiva 2011/61/UE del Parlamento Europeo y del Consejo, relativa a los gestores de fondos de inversión alternativos} (de ahora en adelante, Directiva AIFM); y, (4) la revisión del \href{https://www.boe.es/doue/2014/173/L00084-00148.pdf}{Reglamento (UE) nº 600/2014 del Parlamento Europeo y del Consejo, relativo a los mercados de instrumentos financieros} (de ahora en adelante, Reglamento MiFI).

Más adelante, en diciembre de 2022, la Comisión presentó tres nuevas iniciativas legislativas para aumentar la resiliencia y atractivo de los servicios de compensación y liquidación de valores de la Unión, armonizar algunas normas de insolvencia empresarial en la Unión y un nuevo \textit{Listing Act} para aliviar las cargas de salir a cotizar en mercado, sobre todo de las empresas de menor tamaño.

Como se puede ver, el número de iniciativas legislativas, tanto adoptadas como en proceso de negociación, son muy elevadas. Por ello, a continuación se incluye una breve explicación únicamente de algunas de las iniciativas que, a juicio de quien escribe este tema, son más relevantes.


\subsection{Instrumentos que inciden sobre el Mercado}
Bajo esta categoría se agrupamos aquellos instrumentos de la UMC cuyo objeto regulado no es un sujeto concreto sino el propio mercado, entendido en un sentido deliberadamente amplio que incluye tanto el instrumento financiero negociable en sí mismo como la infraestructura o plataforma a través de la cual ese instrumento cambia de manos entre dos partes distintas del sujeto que lo creó originalmente. 

La razón de fondo que explica por qué el legislador europeo ha dedicado un esfuerzo regulatorio tan intenso precisamente a esta dimensión estructural del mercado, y no solo a los sujetos que en él participan, se encuentra en el diagnóstico mismo que dio origen a toda la UMC: los mercados de capitales de la Unión están comparativamente poco desarrollados frente a otras grandes economías, y buena parte de esa infradesarrollo se debe a que esta infraestructura necesaria o bien no existían de forma armonizada, o bien había quedado desprestigiados y mal calibrados prudencialmente tras la crisis financiera global. 

\subsubsection{Titulización}
La titulización es la operación mediante la cual una entidad de crédito u otra entidad originadora agrupa un conjunto de préstamos que mantiene en su balance y los transforma en valores negociables que después vende a inversores, sustituyendo así su posición de acreedor original, y las rentas periódicas que esa posición lleva aparejadas, por una cantidad fija y pactada de dinero recibida de manera anticipada. 

La operación cumple, en esencia, una doble función: por un lado, transfiere el riesgo de crédito de esos préstamos desde el balance de la entidad originadora hacia los inversores que adquieren los valores; por otro lado, libera precisamente por esa vía un margen operativo en el balance de la entidad originadora, equivalente aproximadamente al valor de los activos transferidos, que le permite volver a conceder nuevos préstamos sin necesidad de captar capital adicional. 

Este segundo efecto es el que explica el interés regulatorio de la titulización desde la perspectiva del legislador: el espacio de balance así liberado resulta un instrumento especialmente valioso para financiar los grandes retos de inversión que enfrenta la Unión: la transición ecológica, el impulso a la competitividad, la economía digital y, más recientemente, el refuerzo de la industria de defensa, sin que ello exija necesariamente una expansión paralela del capital regulatorio del sector bancario.

\paragraph{Estructura actual: ¿Cómo se regulan los intrumentos de titulización?} 
El marco jurídico que regula esta operación está compuesto por tres instrumentos que actúan de manera coordinada:\footnote{%
    Este diseño normativo, en vigor desde 2019, responde directamente a la experiencia de la crisis financiera global de 2008, en la que un uso desregulado de la titulización, caracterizado por la práctica conocida como \textit{originar para distribuir} (conceder préstamos con la intención de titulizarlos y venderlos de inmediato), contribuyó de manera decisiva a la propagación del riesgo por todo el sistema financiero internacional, precisamente porque los originadores carecían de incentivo alguno para evaluar con rigor la calidad crediticia de los préstamos que titulizaban.
}
\begin{la}
    \item \textbf{Reglamento (UE)}. El primero es el Reglamento sobre Titulización, que introduce normas comunes de diligencia debida, de retención de riesgo por parte del originador y de transparencia informativa para todas las titulizaciones realizadas en la Unión, y que crea, dentro de ese marco general, una categoría reforzada de productos denominados titulizaciones simples, transparentes y normalizadas, conocidas por sus siglas en inglés como STS, que se benefician de un tratamiento regulatorio más favorable precisamente por cumplir criterios adicionales de calidad estructural;
    \item \textbf{Reglamento (UE) 575/2013}. El segundo es la modificación correspondiente del Reglamento de Requisitos de Capital, que ya tratamos, y que ajusta el tratamiento del capital regulatorio que las entidades de crédito deben mantener frente a sus posiciones de titulización, haciéndolo más sensible al riesgo efectivo de la operación e introduciendo un tratamiento preferente específico para las titulizaciones que ostentan la etiqueta STS; y,
    \item \textbf{Directiva Solvencia II}. El tercero es la modificación paralela de las normas prudenciales de la Directiva Solvencia II, que ajusta de manera equivalente el tratamiento de las inversiones que las compañías aseguradoras realizan en instrumentos de titulización, extendiendo así la misma lógica de sensibilidad al riesgo desde el sector bancario hacia el asegurador.
\end{la}  

El objetivo declarado del marco vigente es doble y deliberadamente equilibrado: por un lado, preservar la titulización como herramienta legítima y valiosa de financiación de la economía real, evitando repetir el error de prohibirla o penalizarla en exceso; por otro lado, eliminar específicamente las prácticas que en el pasado generaron riesgo para la estabilidad financiera, mediante la retención obligatoria de riesgo por parte del originador y la exigencia de transparencia frente al inversor. En 2021, el ámbito de la etiqueta STS se amplió además a las titulizaciones sintéticas dentro de balance, una modalidad en la que el originador no transfiere los activos subyacentes sino únicamente el riesgo de crédito asociado a ellos mediante CDS o CDO (\textit{Collateral Debt Security} o \textit{Collateral DebtObligation}), reconociendo así que esta variante, correctamente estructurada, podía beneficiarse de las mismas ventajas prudenciales que la titulización tradicional con transferencia efectiva de activos.

\paragraph{Reforma en curso (2025-2026): Flexibilidad y coherencia} 
El marco descrito, pese a su solidez prudencial, ha sido objeto en los últimos años de una crítica reiterada por parte de la industria y de la propia Comisión: el mercado de titulización de la Unión no ha logrado recuperar, ni de lejos, el volumen que tenía antes de la crisis financiera, en buena medida porque los requisitos de capital resultantes del marco resultaron excesivamente conservadores en relación con el riesgo real de muchas operaciones, y porque las cargas operativas de cumplimiento resultaron desproporcionadas para el tamaño de buena parte de las emisiones. 

Por ello, el 17 de junio de 2025 la Comisión presentó un paquete de revisión que constituye la primera iniciativa legislativa concreta de la nueva estrategia de la Unión del Ahorro y las Inversiones que ya mencionamos antes, y que propone modificar simultáneamente el Reglamento sobre Titulización y el Reglamento de Requisitos de Capital. Su contenido más relevante consiste en sustituir los suelos de ponderación de riesgo actualmente fijos ($10\%$ para tramos senior de operaciones STS y el $15\%$ para las que no lo son, con independencia de la calidad real de los activos subyacentes),  por un enfoque más granular y sensible al riesgo, en el que el \textbf{requisito de capital se vincula directamente a la calidad crediticia media de la cartera subyacente}, si bien manteniendo un suelo mínimo, situado entre el $5-15\%$, para evitar una descapitalización excesiva y preservar la coherencia con los estándares internacionales de Basilea III. 

La propuesta reduce además, del $100\to70\%$, el umbral de homogeneidad de préstamos a pequeñas y medianas empresas que una cartera debe cumplir para poder acogerse al régimen STS, permitiendo así que carteras mixtas con un componente mayoritario, pero no exclusivo, de financiación a PYME se beneficien de este tratamiento preferente. Introduce tambié una nueva categoría de titulización especialmente segura, denominada \textit{resiliente}, que se sumaría a la etiqueta STS ya existente para ofrecer un ahorro de capital todavía mayor a las estructuras de menor riesgo; y permite que la verificación del cumplimiento de los criterios STS pueda realizarse por verificadores externos independientes, aliviando así la carga administrativa actualmente asociada a esa comprobación.

No obstante, a fecha de esta actualización no es todavía derecho vigente sino que se encuentra en un estadio avanzado del procedimiento legislativo ordinario: el Consejo alcanzó su posición negociadora general el 19 de diciembre de 2025, el PE aprobó formalmente su propia posición el 5 de mayo de 2026, y las negociaciones tripartitas entre ambas instituciones y la Comisión para acordar el texto definitivo se esperan durante la segunda mitad de 2026.\footnote{%
    Uno de los puntos de fricción que persisten en esas negociaciones tripartitas es hasta qué punto debe ampliarse la protección crediticia elegible para las titulizaciones sintéticas STS: tanto el Consejo como el Parlamento apoyan, en sus respectivas posiciones, incluir garantías no financiadas otorgadas por compañías de reaseguro dentro de esa protección elegible, pese a que el BCE emitió en noviembre de 2025 un dictamen mostrando reservas significativas sobre este punto concreto, por el riesgo de que una garantía de reaseguro resulte, en la práctica, menos líquida y menos verificable que una garantía o colateral puramente financiero.
} 

\subsubsection{Bonos garantizados}
Los bonos garantizados son instrumentos de deuda emitidos por entidades de crédito que ofrecen a sus titulares lo que la doctrina denomina protección de doble recurso:\footnote{%
    Esta doble capa de protección es lo que distingue estructuralmente al bono garantizado tanto de un bono de deuda ordinaria, que solo otorga un derecho general sobre el patrimonio del emisor, como de un valor de titulización, en el que los activos subyacentes salen efectivamente del balance del originador y el inversor asume el riesgo de esos activos de forma aislada, sin recurso adicional contra el patrimonio general del originador.
}
\begin{la}
    \item \textbf{Cobertura de activos}. Por un lado, el titular del bono ostenta un derecho directo y preferente sobre un conjunto de activos específicamente asignados a garantizar la emisión (conjunto o cobertura de activos), que permanece segregado del resto del balance del emisor; y,
    \item \textbf{Derecho ordinario}. Por otro lado, y de manera adicional, conserva también un derecho ordinario, no preferente, sobre el resto de los activos del emisor en caso de que la cobertura preferente resultara insuficiente para satisfacer su crédito. 
\end{la} 

Los bonos garantizados constituyen, por ello, una \textbf{fuente de financiación especialmente eficiente para las entidades de crédito} y ofrecen un elevado nivel de certidumbre a los inversores, lo que explica que sus mercados se encuentren particularmente desarrollados en aquellos Estados miembros que cuentan con regímenes nacionales de larga tradición, entre ellos Alemania, Dinamarca, Francia, España, Italia, Luxemburgo y Suecia.\footnote{%
    A cierre de diciembre de 2015, el volumen vivo de bonos garantizados emitidos por entidades con sede en la Unión alcanzaba los 2,1 billones de euros, lo que representaba el ochenta y cuatro por ciento del volumen total de este instrumento a escala mundial, un dato que por sí solo explica por qué el legislador europeo consideró prioritario armonizar un instrumento en el que la Unión ya ostentaba, de facto, una posición dominante a nivel global antes incluso de contar con un marco jurídico común.
}

\paragraph{Estructura actual: ¿Cómo se regulan los bonos garantizados?} 
El marco jurídico armonizador está compuesto por una Directiva y un Reglamento que actúan conjuntamente: la Directiva sobre la emisión de bonos garantizados y su supervisión pública, y el Reglamento que modifica correspondientemente el Reglamento de Requisitos de Capital en lo relativo al tratamiento prudencial de las exposiciones en forma de bonos garantizados. 

Ambos textos fueron presentados por la Comisión Europea en marzo de 2018 y acordados por los colegisladores a comienzos de 2019, formalizándose como Directiva (UE) 2019/2162 y su Reglamento correspondiente el 27 de noviembre de 2019. El objetivo declarado de este marco no es crear el instrumento, que ya existía y funcionaba con éxito en varios Estados miembros desde hace décadas, sino establecer unos requisitos mínimos de armonización que todos los bonos garantizados comercializados en la Unión deben cumplir, precisamente para que la seguridad del instrumento se extienda también a aquellos Estados donde el mercado estaba menos desarrollado o carecía de un régimen legal específico, abriendo así nuevas oportunidades de financiación allí donde antes no existía un marco de confianza equivalente.

El marco armonizado proporciona, en primer lugar, una \textbf{definición común y vinculante} de qué constituye un bono garantizado a efectos de todo el derecho de la Unión, sustituyendo así la disparidad de definiciones nacionales previamente existente. En segundo lugar, define las \textbf{características estructurales exigibles} al instrumento, entre ellas los criterios de elegibilidad de los activos que pueden formar parte de la cobertura, los requisitos de sobrecolateralización y los mecanismos de gestión de la liquidez de la cobertura ante eventuales desfases de vencimiento. En tercer lugar, define con precisión las \textbf{tareas y responsabilidades de supervisión pública de los bonos garantizados}, exigiendo a cada Estado miembro designar una autoridad competente específica para esta función, distinta en su ámbito material de la supervisión prudencial general de la entidad emisora aunque pueda residir en la misma institución. En cuarto lugar, establece las \textbf{condiciones que permiten a un emisor utilizar la etiqueta bonos garantizados europeos}, una denominación protegida que solo puede emplearse cuando la emisión cumple efectivamente con la totalidad de los requisitos armonizados, de modo que el inversor pueda confiar en ella como sello de calidad reconocible en toda la Unión sin necesidad de verificar caso por caso el régimen nacional aplicable. Y en quinto lugar, \textbf{refuerza y clarifica las condiciones bajo las cuales un bono garantizado puede beneficiarse del trato prudencial preferente} que ya le reconocía el Reglamento de Requisitos de Capital, vinculando ese trato preferente al cumplimiento efectivo de los nuevos requisitos armonizados y no ya, como ocurría antes de 2019, a una remisión más laxa a las tradiciones regulatorias nacionales.

\paragraph{Revisión periódica (Comisión): En curso} 
El marco de 2019 es plenamente aplicable desde el 8 de julio de 2022, fecha límite que se dio a los Estados miembros para transponer la Directiva a sus respectivos ordenamientos internos, y desde entonces la etiqueta de bono garantizado europeo convive con los regímenes nacionales preexistentes, ahora armonizados en sus elementos mínimos. 

La propia Directiva contempla, en su artículo 31, un mandato de revisión periódica: la Comisión debía presentar, antes del 8 de julio de 2025, un informe al PE y al Consejo sobre la aplicación práctica de la Directiva desde la perspectiva del nivel de protección otorgado al inversor y sobre la evolución general del mercado de bonos garantizados en la Unión, informe que puede incluir recomendaciones para una acción normativa ulterior. 

Para preparar ese informe, la Comisión solicitó en julio de 2023 el asesoramiento técnico de la Autoridad Bancaria Europea, que hizo público su dictamen de respuesta el 23 de septiembre de 2025, elaborado en colaboración con las autoridades nacionales competentes de los Estados miembros. A fecha de esta actualización (julio 2026), la Comisión continúa trabajando en la elaboración de su propio informe final a partir de ese dictamen, sin que se haya presentado todavía ninguna propuesta legislativa concreta de reforma.\footnote{%
    La Directiva contempla además un segundo mandato de revisión más específico, referido a los bonos garantizados con estructuras de vencimiento extensible (aquellos en los que el vencimiento del bono puede prorrogarse automáticamente ante determinadas circunstancias de estrés de liquidez del emisor), respecto de los cuales la Comisión debía encargar un estudio de riesgos y beneficios y, en su caso, acompañar su informe de una propuesta legislativa si lo consideraba oportuno. 

Este proceso de revisión de los bonos garantizados es, por tanto, considerablemente menos avanzado que el que ya vimos para la titulización: mientras que esta última cuenta ya con una propuesta legislativa concreta en fase de negociación tripartita, los bonos garantizados se encuentran todavía en la fase previa de diagnóstico técnico, sin que exista aún un texto normativo de reforma sobre el que negociar.
} 

\subsubsection{Crowdfunding}
El mercado de financiación participativa, o crowdfunding, de la Unión ha sido tradicionalmente un mercado poco desarrollado en comparación con el de otras grandes economías mundiales, y el diagnóstico que llevó al legislador europeo a intervenir en este ámbito no señala como causa principal ni la falta de demanda de financiación por parte de las pequeñas empresas ni la falta de apetito inversor por parte de los ahorradores, sino un \textbf{obstáculo puramente regulatorio}: durante muchos años, las plataformas de crowdfunding que pretendían ofrecer sus servicios a escala transfronteriza se encontraron con requisitos de concesión de licencias divergentes entre Estados miembros y con una ausencia casi total de normas comunes en toda la Unión, lo que generó unos costes operativos y de cumplimiento tan elevados que impidió a las plataformas expandir eficazmente su actividad más allá de su mercado nacional de origen. 

Esta fragmentación regulatoria tuvo, además, un efecto negativo de doble sentido: por el lado de la oferta de financiación, las pequeñas empresas dispusieron de menos opciones de financiación alternativa a la bancaria; y por el lado de la demanda de inversión, los inversores se enfrentaron a menos opciones de colocación de su ahorro y a una mayor incertidumbre a la hora de invertir de forma transfronteriza, precisamente por la disparidad de estándares de protección entre los distintos regímenes nacionales.

\paragraph{Estructura actual: ¿Cómo se regula el mercado de crowdfunding?} 
La respuesta del legislador es el Reglamento (UE) 2020/1503, sobre proveedores europeos de servicios de financiación participativa, que establece normas uniformes y directamente aplicables en toda la Unión para la prestación de servicios de financiación participativa basados tanto en la inversión como en el préstamo, siempre que ambos estén vinculados a la financiación de una actividad empresarial. El Reglamento es plenamente aplicable desde noviembre de 2021, con un régimen transitorio para las plataformas ya autorizadas conforme a regímenes nacionales previos que se extendió hasta noviembre de 2023, fecha a partir de la cual toda plataforma que preste estos servicios en la Unión debe contar necesariamente con la autorización armonizada. 

El elemento central de su diseño es el \textbf{pasaporte único}: una plataforma que obtiene autorización como proveedor de servicios de financiación participativa conforme a un único conjunto de normas puede, a partir de esa autorización inicial, ofrecer sus servicios en cualquier otro Estado miembro.

El régimen se apoya en tres elementos de protección que operan de manera acumulativa. En primer lugar, \textbf{normas claras de divulgación} de información, exigibles tanto a los promotores de los proyectos que buscan financiación como a las propias plataformas, de modo que el inversor cuente con información homogénea con independencia del Estado miembro de origen de la plataforma o del proyecto. En segundo lugar, \textbf{normas de gobernanza y de gestión de riesgos exigibles directamente a las plataformas} de financiación participativa, orientadas a asegurar que su propio funcionamiento interno no introduzca riesgos adicionales para los inversores que a través de ellas colocan su capital. Y en tercer lugar, competencias de \textbf{supervisión sólidas y armonizadas} atribuidas a las autoridades nacionales competentes encargadas de vigilar el funcionamiento de estas plataformas, de modo que la supervisión, aunque siga siendo nacional en cuanto a la autoridad que la ejerce, responda a un mismo estándar mínimo en toda la Unión.

\paragraph{Estado actual del mercado y de la revisión del Reglamento.} 
El mercado se ha consolidado de forma modesta pero constante: en torno a doscientas cincuenta plataformas operan hoy bajo esta autorización armonizada en el conjunto de la Unión, habiéndose superado ya los $4.200\text{M}€$ captados acumulados bajo este régimen desde su entrada en aplicación. 

Sin embargo, la propia finalidad armonizadora del Reglamento (pasaporte único) dista todavía de haberse alcanzado en la práctica: según datos recopilados a comienzos de 2025, solo una minoría de las plataformas autorizadas había solicitado efectivamente el pasaporte para operar en otros Estados miembros, y entre quienes lo han solicitado existe una disparidad considerable en el número de países de destino declarados, lo que sugiere que persisten fricciones prácticas, lingüísticas y de interpretación divergente de la norma entre autoridades nacionales, pese a existir ya un reglamento único y directamente aplicable.\footnote{%
    Distintos informes independientes de seguimiento del sector, elaborados durante 2025 a partir del registro público de la Autoridad Europea de Valores y Mercados, han documentado además un grado desigual de cumplimiento de las obligaciones de transparencia y divulgación, con una puntuación notablemente mejorable y una minoría de plataformas alcanzando un nivel de cumplimiento calificado como excelente. De nuevo, vemos como contar con un reglamento único no garantiza, por sí solo, una aplicación uniformemente rigurosa por parte de todos los operadores del mercado, y la Autoridad Europea de Valores y Mercados no dispone todavía de un marco de reporte estandarizado que permita comparar de forma homogénea variables clave del sector, como las tasas de impago o de recuperación entre plataformas.
}

El propio Reglamento contempla una cláusula de revisión periódica (art. $45$), actualmente en curso de evaluación por parte de la Comisión, en la que uno de los debates más relevantes es si debe elevarse el umbral máximo de financiación por proyecto y periodo de doce meses, actualmente fijado en cinco millones de euros, hasta una cifra sensiblemente superior, con el fin de permitir que operaciones de mayor envergadura, hoy excluidas del régimen armonizado por superar ese umbral, puedan beneficiarse también del pasaporte único en lugar de tener que acudir a regímenes nacionales o a otros instrumentos de financiación de mercado ya vistos, como el propio folleto o la titulización.

\subsubsection{SME Growth Market}
Cotizar en bolsa puede dar un impulso significativo a las pequeñas y medianas empresas, por varias vías acumulativas: reduce su dependencia de la financiación bancaria, diversifica la base de inversores a la que pueden acudir, facilita un acceso más ágil a capital propio adicional y a financiación de deuda en rondas sucesivas, y proporciona un mayor perfil público y reconocimiento de marca frente a clientes, proveedores y potenciales socios. 

Pese a estas ventajas, los mercados públicos de la Unión dedicados a las pequeñas y medianas empresas llevan tiempo enfrentando serias dificultades para atraer nuevos emisores: Europa produce en la actualidad solo la mitad de las ofertas públicas iniciales de pequeñas y medianas empresas que generaba antes de la crisis financiera global, un dato que por sí solo evidencia que el problema no es coyuntural sino estructural. 

El legislador ha diagnosticado este problema como un fallo de mercado con dos caras que se retroalimentan mutuamente. Por el lado de la oferta de valores, los emisores se enfrentan a unos costes de cumplimiento normativo elevados en relación con su tamaño, lo que hace que cotizar resulte proporcionalmente mucho más gravoso para una pequeña empresa que para una gran cotizada. Por el lado de la demanda de esos valores, la insuficiente liquidez de estos mercados de menor tamaño perjudica simultáneamente a los tres tipos de participantes: a los propios emisores, que soportan un mayor coste de capital precisamente por cotizar en un mercado poco líquido; a los inversores, que se muestran reacios a colocar su ahorro en valores de baja liquidez y mayor riesgo de volatilidad asociado; y a los intermediarios financieros, cuyos modelos de negocio dependen de un flujo constante de órdenes que solo existe en mercados suficientemente líquidos. Se trata, por tanto, de una externalidad de red clásica: la falta de liquidez ahuyenta a los participantes, y su ausencia perpetúa a su vez la falta de liquidez.

\paragraph{Instrumento legal.} 
La respuesta normativa a este diagnóstico es la creación del SME Growth Market como una nueva categoría, dentro de los sistemas multilaterales de negociación ya existentes, introducida por la Directiva de Mercados de Instrumentos Financieros, conocida como MiFID II. 

La opción de diseño del legislador es significativa: en lugar de crear un mercado regulado tradicional, se opta por habilitar una submodalidad de sistema multilateral de negociación, sujeta ya de por sí a requisitos regulatorios más ligeros que un mercado regulado, y sobre esa submodalidad ya aligerada se construyen a su vez requisitos específicamente calibrados para el tamaño de las empresas que en ella cotizan, en un ejercicio de proporcionalidad regulatoria en dos niveles sucesivos. 

Este régimen se ha ido perfeccionando desde su creación mediante sucesivas iniciativas legislativas, entre las que destaca el llamado SME Listing Act, formalmente el Reglamento (UE) 2019/2115, adoptado el 27 de noviembre de 2019 y aplicable desde el 1 de enero de 2021, que introdujo ajustes específicos para facilitar tanto el acceso inicial a la cotización en un SME Growth Market como la migración posterior de una empresa desde este mercado hacia un mercado regulado de mayor exigencia, a medida que la propia empresa crece.

\paragraph{Última reforma: Listing Package} 
La última reforma de este marco ha sido el Listing Act, que entró en vigor el 4 de diciembre de 2024, y buena parte de su contenido incide de manera directa y específica sobre el régimen del SME Growth Market (también sobre folletos, que veremos más adelante).

El Listing Act extiende a todos los emisores, y no solo a los que ya cotizaban en un SME Growth Market, el formato simplificado de listas de iniciados. Además, introduce la posibilidad de que una empresa que solicite por primera vez su admisión a negociación adopte una estructura de acciones con voto múltiple, permitiendo a los fundadores conservar el control de la sociedad tras la salida a cotización sin necesidad de diluir su capacidad decisoria, bajo salvaguardas de proporción máxima de voto y de protección de los accionistas sin voto reforzado (mediante una directiva específica sobre estructuras de voto múltiple).\footnote{%
    También revierte parcialmente el régimen de separación obligatoria entre el pago de servicios de análisis financiero y el pago de servicios de ejecución de órdenes que había introducido la propia MiFID II en 2018, permitiendo de nuevo, desde el 6 de junio de 2026, el pago conjunto de ambos servicios bajo determinadas condiciones de transparencia frente al cliente final, motivada precisamente porque aquella separación había reducido la cobertura de análisis financiero disponible sobre empresas de pequeña y mediana capitalización, agravando indirectamente el propio problema de liquidez que el SME Growth Market pretende resolver.
}

El Listing Act exige a la Comisión Europea adoptar, en los dieciocho meses siguientes a su entrada en vigor, una serie de actos delegados que desarrollen técnicamente estas modificaciones, para lo cual solicitó el asesoramiento técnico de la Autoridad Europea de Valores y Mercados. Ésta hizo público su informe final de asesoramiento técnico sobre las cuestiones relativas al abuso de mercado y al régimen del SME Growth Market el 7 de mayo de 2025, y la Comisión tiene previsto adoptar los correspondientes actos delegados antes de julio de 2026. Los Estados miembros debían, por su parte, haber transpuesto las modificaciones de la Directiva MiFID II antes del 5 de junio de 2026 y las relativas a las estructuras de voto múltiple antes del 5 de diciembre de 2026.\footnote{%
    Este calendario de implementación por fases —actos delegados de nivel dos culminando en 2026, transposición nacional de las directivas también fijada en 2026— sitúa al régimen del SME Growth Market en una fase de aplicación todavía incompleta a fecha de esta actualización (julio 2026): el marco legal ya existe y es vinculante en su contenido de primer nivel desde diciembre de 2024, pero buena parte de su desarrollo técnico definitivo seguirá materializándose a lo largo de 2026, por lo que conviene presentar este régimen en el examen como derecho en fase de implementación activa, no como un marco ya plenamente asentado en su versión final.
}


\subsection{Instrumentos que inciden sobre los agentes}
Frente a la categoría anterior, en la que el legislador europeo interviene sobre el instrumento negociable o sobre la infraestructura de intercambio, esta segunda categoría agrupa aquellos instrumentos cuyo objeto regulado es directamente un sujeto que participa en el mercado. 

Dentro de esta categoría conviene distinguir tres tipos de agente que reciben un tratamiento normativo distinto, precisamente porque desempeñan funciones económicas distintas dentro del ciclo de la financiación: el emisor, que demanda capital del mercado para financiar su actividad y que, por ello, está sujeto fundamentalmente a normas de divulgación de información en el momento concreto en que acude al mercado; el intermediario, que presta un servicio de conexión entre quien demanda y quien oferta capital, y que por ello está sujeto a un régimen de supervisión prudencial continuada, no limitada a una operación puntual sino extendida a toda su actividad en el tiempo; y el inversor u oferente de capital, cuya protección constituye el objeto directo de la norma, a través de productos diseñados específicamente para facilitar y proteger su participación en el mercado. 

\subsubsection{Emisores: Folletos}
La mayoría de las empresas que desean captar capital mediante una oferta pública de valores, o que tienen valores que van a ser admitidos a cotización en un mercado regulado, están obligadas a poner a disposición de los inversores un folleto, entendido como un documento legal que debe describir necesariamente tres elementos: 
\begin{la}
    \item \textbf{Información empresarial}. Información relativa a la empresa. En particular, la actividad principal de la empresa, su situación financiera y su estructura accionarial;
    \item \textbf{Información de la emisión}. Información relativa a la emisión: características de los valores concretos que se emiten o que se admiten a cotización; y,
    \item \textbf{Información adicional (\textit{due dilligence})}. Bajo el paraguas de \textit{due dilligence}, la información que un inversor razonable necesitaría conocer antes de decidir si invierte en esos valores.
\end{la}

La finalidad de esta exigencia es corregir directamente el fallo de mercado clásico de la información asimétrica entre el emisor, que conoce su propia situación financiera y sus perspectivas de negocio con un grado de detalle que el inversor externo no puede replicar por sí mismo, y el inversor, que debe adoptar su decisión de inversión con la información que el emisor decida divulgarle.

\paragraph{Estructura actual: ¿Cómo se regula este fallo de mercado?} 
El marco normativo vigente es el Reglamento (UE) 2017/1129, que sustituyó a la anterior Directiva de Folletos de 2003, precisamente para sustituir la armonización de mínimos por una armonización plena mediante un reglamento. 

Las normas de la Unión sobre folletos garantizan así la existencia de estándares de divulgación adecuados y equivalentes en todos los Estados miembros, de modo que los inversores se beneficien del mismo nivel mínimo de información con independencia de dónde esté establecido el emisor, y establecen además una protección mínima del inversor geográficamente uniforme en toda la Unión. 

Sobre esta base de contenido idéntico opera el elemento más característico del régimen: el \textbf{pasaporte único para los emisores}, en virtud del cual un folleto aprobado por la autoridad competente de un Estado miembro es automáticamente válido para realizar la misma oferta o admisión a cotización en cualquier otro Estado miembro.

\paragraph{Última reforma: Listing Act} 
Tal como ya adelantamos al cerrar el epígrafe de titulización (aunque éste está de nuevo en un proceso de reforma) y de nuevo al tratar el SME Growth Market, el Reglamento de Folletos de 2017 ha sido objeto de una modificación sustancial mediante el Listing Act, en vigor desde el 4 de diciembre de 2024, cuyo contenido específico en esta materia conviene detallar aquí con precisión.

La reforma introduce, en primer lugar, exenciones adicionales a la obligación de publicar un folleto y amplía el alcance de exenciones ya existentes, elevando además los umbrales cuantitativos por debajo de los cuales una oferta de valores queda excluida de esta obligación, y eliminando el umbral mínimo de un millón de euros que anteriormente marcaba, en sentido inverso, el límite por debajo del cual el propio Reglamento de Folletos no resultaba aplicable en absoluto. En segundo lugar, crea un nuevo folleto simplificado de continuación, el Folleto de Seguimiento de la Unión, destinado específicamente a emisiones secundarias realizadas por empresas que ya cuentan con valores admitidos a cotización en un mercado regulado o en un SME Growth Market, bajo la lógica de que una empresa ya sujeta a obligaciones continuadas de información no necesita repetir el mismo nivel de detalle que exigiría una primera salida a los mercados. En tercer lugar, sustituye el anterior Folleto de Crecimiento de la Unión por un nuevo Folleto de Emisión de Crecimiento de la Unión, que simplifica todavía más el documento exigible a las pequeñas y medianas empresas, en línea directa con la lógica de proporcionalidad que ya vimos al desarrollar el SME Growth Market. Y en cuarto lugar, estandariza y simplifica de forma general el contenido y el formato de cualquier folleto, con el objetivo declarado de facilitar su comparación entre distintas emisiones y distintos emisores.

La generalidad de estas modificaciones al Reglamento de Folletos resulta aplicable desde el 4 de diciembre de 2024, si bien determinadas disposiciones de mayor calado técnico, que exigían el desarrollo previo de actos delegados o de ejecución de segundo nivel, no han comenzado a aplicarse hasta un momento posterior dentro de 2026.\footnote{%
    La Comisión Europea abrió además, en paralelo a esta reforma, una consulta pública sobre la posible armonización de las normas de responsabilidad civil derivada del contenido de un folleto de emisión, cuestión que hasta la fecha permanece regulada de forma dispar entre los distintos Estados miembros pese a la armonización ya alcanzada sobre el contenido mismo del documento: el Reglamento de Folletos armoniza qué debe divulgarse, pero no armoniza, todavía, qué ocurre jurídicamente frente a un emisor cuando la información divulgada resulta ser incorrecta o engañosa, cuestión que sigue remitiéndose al derecho nacional de responsabilidad civil de cada Estado miembro.
}

\subsubsection{Intermediarios: Revisión del marco prudencial de las firmas de inversión}
Las empresas de servicios de inversión no demandan capital para financiar su propia actividad productiva, como hace el emisor que coloca un folleto, ni ofrecen su ahorro como haría un inversor particular, sino que prestan un servicio de conexión entre quienes demandan y quienes ofrecen capital, ejecutando órdenes, colocando instrumentos financieros o gestionando carteras por cuenta de terceros. 

Sus actividades se rigen, en cuanto a la conducta frente al cliente y la organización de sus servicios, por la Directiva de Mercados de Instrumentos Financieros (MiFID II) que ya introdujimos, pero además están sujetas a normas prudenciales específicas cuyo objeto no es proteger al inversor frente a una operación puntual mal informada, como ocurre con el folleto, sino otro de carácter prudencial que persigue un doble objetivo: reducir el \textbf{riesgo de quiebra} de la empresa de inversión y, con ello, reducir tanto el riesgo de que sus clientes sufran un perjuicio económico indebido por la insolvencia de su intermediario como el riesgo de que esa quiebra provoque perturbaciones en los mercados en los que la empresa opera, en una lógica de prevención de \textbf{externalidad negativa}.

\paragraph{Estructura actual: ¿Cómo se regulan estos intermediarios financieros?} 
Hasta el 25 de junio de 2021, las normas prudenciales aplicables a las empresas de inversión formaban parte, sin distinción sustancial respecto de las entidades de crédito, del propio Reglamento de Requisitos de Capital y de su Directiva hermana que ya trabajamos extensamente al desarrollar el Código Normativo Único. Sin embargo, sesde el 26 de junio de 2021 la mayoría de las empresas de inversión pasaron a regirse por un marco prudencial propio y específicamente diseñado para ellas: el Reglamento (UE) 2019/2033, relativo a los requisitos prudenciales de las empresas de servicios de inversión, conocido por sus siglas en inglés como IFR, y la Directiva (UE) 2019/2034, relativa a la supervisión prudencial de las empresas de servicios de inversión, conocida como IFD.\footnote{%
    La razón de ser de este nuevo régimen separado es que el marco CRR/CRD, pensado originalmente para entidades de crédito que captan depósitos del público y conceden crédito, resultaba mal calibrado para el perfil de riesgo real de la mayoría de las empresas de inversión, que no captan depósitos ni conceden crédito en sentido propio, y que por ello merecían requisitos de capital más sencillos y más coherentes con la naturaleza, el tamaño y la complejidad efectiva de su actividad.
}

El nuevo marco distingue tres categorías de empresas de inversión según su tamaño y su modelo de negocio, y esta clasificación determina qué régimen prudencial les resulta aplicable:
\begin{lnum}
    \item \textbf{Clase 3}. Las denominadas empresas de \textit{clase 3} son aquellas consideradas suficientemente pequeñas y no interconectadas con el resto del sistema financiero como para no plantear riesgo sistémico;
    \item \textbf{Clase 2}. Las empresas de \textit{clase 2} son aquellas que no cumplen los criterios para clasificarse como clase 3 pero tampoco alcanzan el tamaño de las categorías superiores; y,
    \item \textbf{Clase 1}. Las llamadas empresas de \textit{clase 1}, que negocian por cuenta propia o suscriben y colocan instrumentos financieros sobre la base de un compromiso firme y cuyos activos consolidados alcanzan los treinta mil millones de euros, y que precisamente por ese volumen de actividad deben solicitar autorización como entidades de crédito propiamente dichas. Este tercer nivel de empresas de inversión, las de mayor tamaño relativo, no se rigen por el nuevo marco sino que permanecen sujetas al régimen general CRR/CRD de las entidades de crédito.\footnote{%
        Junto a ellas, las llamadas empresas de \textit{clase 1 minus}, que desarrollan el mismo tipo de actividad pero con un volumen de activos consolidados de quince mil millones de euros, o de cinco mil millones cuando así lo determine su autoridad nacional competente conforme a criterios específicos de riesgo. Éstas siguen estando autorizadas como empresas de inversión y no como entidades de crédito, aunque quedan igualmente sujetas al régimen prudencial general en lugar de al régimen específico del IFR.
    }
\end{lnum}

La gran mayoría de las empresas de inversión de la Unión pertenecen a las dos primeras categorías sujetas plenamente al nuevo régimen IFR/IFD. Para estas, el IFR establece un requisito de capital compuesto por tres elementos: (i) un \textbf{requisito de gastos generales fijos}, equivalente a una cuarta parte de los gastos generales fijos anuales de la propia empresa; (ii) un \textbf{requisito de capital mínimo} permanente, que según el tipo de actividad concreta que desarrolle la empresa se sitúa en 75.000, 150.000 o 750.000 euros; y, (iii) un \textbf{requisito de capital global basado en el llamado \textit{factor K}}, que resulta de sumar una serie de indicadores de riesgo agrupados en tres bloques: riesgo para el cliente, riesgo para el mercado y riesgo para la propia empresa. 

\paragraph{Revisión periódica (Comisión): En curso} 
El IFR contempla un mandato de revisión periódica a cargo de la Comisión, que puede desembocar, si así lo considera oportuno, en una propuesta legislativa de modificación del régimen. 

Para preparar esa revisión, la Comisión remitió el 1 de febrero de 2023 una solicitud formal de asesoramiento técnico a la Autoridad Bancaria Europea y a la Autoridad Europea de Valores y Mercados, que publicaron conjuntamente, con cierto retraso respecto al calendario inicialmente previsto, su informe final de asesoramiento el 15 de octubre de 2025. 

La conclusión general de ambas autoridades es que el marco actual cumple razonablemente bien sus objetivos originales, pero identifica al mismo tiempo una serie de cuestiones técnicas concretas que justificarían modificaciones específicas del Reglamento, de la Directiva o de sus normas técnicas de desarrollo, entre ellas una recalibración de la metodología empleada para fijar los umbrales cuantitativos de quince mil y cinco mil millones de euros que separan a las empresas de \textit{clase 1 minus} del resto, tomando como referencia metodológica el propio umbral de treinta mil millones ya empleado para las empresas de \textit{clase 1}.\footnote{%
    Este proceso de revisión discurre en paralelo, y no debe confundirse, con la reforma de la Directiva de Requisitos de Capital que ya trabajamos extensamente al desarrollar el primer pilar de la Unión Bancaria: la CRD VI, aplicable desde el 11 de enero de 2026, modifica el artículo 8a de la propia Directiva de Requisitos de Capital, precepto que regula precisamente el procedimiento por el cual una empresa de inversión de \textit{clase 1} debe solicitar su autorización como entidad de crédito al superar el umbral de los treinta mil millones de euros; se trata, por tanto, de dos procesos de revisión distintos que convergen, sin embargo, en el mismo punto de fricción concreto: dónde situar exactamente la frontera cuantitativa entre una empresa de inversión y una entidad de crédito a todos los efectos.
} 

A fecha de esta actualización, la Comisión Europea no ha presentado todavía una propuesta legislativa formal de modificación del IFR o del IFD a partir de este asesoramiento técnico, aunque distintos observadores del sector consideran probable que dicha propuesta se presente durante el propio año 2026, en un estadio del proceso, por tanto, todavía anterior al que ya hemos visto alcanzado por la reforma de la titulización, y más próximo al estadio de mero diagnóstico técnico en el que se encuentra todavía la revisión del marco de bonos garantizados.

\subsubsection{Inversores: el producto paneuropeo de planes de pensiones (PEPP)}
Las pensiones personales, o pensiones privadas, son productos de ahorro a largo plazo a los que los particulares contribuyen de forma estrictamente voluntaria, y que cumplen una función complementaria (pensiones públicas o, en su caso, pensiones profesionales). 

Dentro de este contexto, el producto paneuropeo de planes de pensiones no regula ni al emisor ni al intermediario, sino directamente la posición jurídica y las opciones disponibles para el inversor u oferente de capital, es decir, para el propio ahorrador particular, con el objetivo de poner en contacto a ese ahorrador con oportunidades de inversión a largo plazo que de otro modo no tendría a su disposición de forma armonizada en toda la Unión.

\paragraph{Estructura actual: ¿Cómo se regula PEPP?} 
El producto paneuropeo de pensiones personales está regulado por el Reglamento (UE) 2019/1238, que sienta las bases jurídicas de un mercado paneuropeo de pensiones personales mediante la normalización de las características fundamentales que cualquier producto de este tipo debe cumplir con independencia del proveedor o del Estado miembro de comercialización: los requisitos de transparencia frente al ahorrador, las normas de inversión aplicables al capital acumulado, el derecho del titular a cambiar de proveedor o de modalidad de inversión a lo largo de la vida del producto, y el tipo de opciones de inversión que deben ofrecerse obligatoriamente. 

Se trata de un plan de pensiones voluntario que puede ser ofrecido en toda la Unión por una amplia gama de proveedores financieros, entre ellos compañías de seguros, gestoras de activos, bancos, determinadas empresas de servicios de inversión y ciertos fondos de pensiones de empleo, todos ellos bajo el mismo conjunto único de normas y con vocación de portabilidad transfronteriza: el titular puede continuar aportando al mismo producto aunque cambie de país de residencia dentro de la Unión, una característica que ningún producto nacional de pensiones ofrece por sí solo.

La finalidad declarada de este Reglamento es dar a los ahorradores más opciones y ofrecerles productos de pensiones personales más competitivos a la hora de ahorrar para la jubilación, todo ello sin renunciar a una sólida protección del consumidor: el producto incorpora, en su modalidad básica, un límite de costes fijado en el $1\%$ del capital acumulado por año, así como una garantía de protección del capital invertido frente a pérdidas, y desde marzo de 2025 los proveedores están además obligados a ofrecer subcuentas nacionales para, al menos, dos Estados miembros distintos, precisamente para reforzar en la práctica la portabilidad transfronteriza que constituye la razón de ser del producto.

\paragraph{Estado actual: Fracaso comercial y proceso de reforma.}
A diferencia de los demás instrumentos, el PEPP debe presentarse como un fracaso comercial ampliamente documentado por la propia industria, por los supervisores europeos y por asociaciones de defensa del inversor. 

Desde su lanzamiento efectivo en 2022, solo uno o dos proveedores han llegado a ofrecer efectivamente el producto en el mercado, una cifra que contrasta radicalmente con la ambición inicial del Reglamento de crear un mercado paneuropeo con una amplia gama de proveedores compitiendo entre sí. La Autoridad Europea de Seguros y Pensiones de Jubilación publicó en septiembre de 2024 un documento de trabajo que analiza en detalle las razones de este fracaso, identificando causas tanto del lado de la oferta como del lado de la demanda: 
\begin{la}
    \item \textbf{Oferta: Comercialización inviable}. Del lado de la oferta, el límite de costes del $1\%$ y la exigencia de asesoramiento obligatorio para la modalidad básica del producto han encarecido su distribución hasta el punto de resultar comercialmente inviable para la mayoría de los proveedores potenciales; y,
    \item \textbf{Demanda: Fiscalidad}. Del lado de la demanda, la falta de un tratamiento fiscal homogéneo entre Estados miembros (la fiscalidad de las pensiones sigue siendo competencia exclusiva de cada Estado miembro) hace que, en muchos países, el producto paneuropeo no disfrute de los mismos incentivos fiscales que sí reciben los productos de pensiones puramente nacionales, neutralizando en la práctica buena parte de su atractivo para el ahorrador medio.
\end{la}

Ante este diagnóstico, la Comisión Europea anunció en marzo de 2025, dentro de la comunicación que presenta la nueva estrategia de la Unión del Ahorro y las Inversiones, su intención de presentar una propuesta de revisión del Reglamento durante el último trimestre de 2025.

Esta propuesta, que se hizo pública en noviembre de 2025 como parte de un paquete más amplio de medidas sobre pensiones complementarias, se encuentra en tramitación ante el Parlamento Europeo y el Consejo. Su contenido más relevante consiste en: (i) eliminar el límite de costes rígido del $1\%$, sustituyéndolo previsiblemente por un mecanismo de evaluación de costes más flexible que pondere el coste frente al beneficio y la rentabilidad ofrecidos; (ii) simplificar la distribución del producto, permitiendo la comercialización de modalidades \textit{seguras por diseño} sin necesidad de asesoramiento obligatorio, precisamente para reducir el coste de distribución que hasta ahora ha desincentivado a los proveedores; y, (iii) explorar fórmulas de incorporación automática del producto en el ámbito laboral, en línea con una de las dos grandes propuestas alternativas que se han barajado en el debate previo a esta reforma, la de expandir el propio modelo del producto paneuropeo hacia un producto de ahorro a largo plazo de adscripción automática a escala de la Unión, frente a la propuesta alternativa, de corte más descentralizado, de dotar de una etiqueta europea de calidad a los productos de pensiones nacionales que ya funcionan razonablemente bien en determinados Estados miembros, en lugar de insistir en un producto paneuropeo único que hasta ahora no ha logrado el respaldo del mercado.\footnote{%
    Esta disyuntiva entre profundizar el modelo paneuropeo único (enfoque defendido en el Informe LETTA) y reconocer, mediante una etiqueta de calidad, la diversidad de productos nacionales ya existentes (enfoque defendido en el Informe NOYER sobre el ahorro y la inversión europeos) resulta especialmente relevante porque ilustra, en miniatura, la misma tensión de fondo que recorre toda la UMC: armonizar mediante un instrumento único de arriba abajo o reconocer, con un mínimo de estándares comunes, la heterogeneidad de soluciones nacionales ya existentes.
} 


\subsection{European Single Access Point (ESAP)}
El Punto de Acceso Único Europeo constituye una categoría propia. Su objetivo es más modesto y, al mismo tiempo, transversal a todo lo anterior: ofrecer un acceso gratuito, fácil de usar, centralizado y digital a la información financiera y relacionada con la sostenibilidad que las empresas europeas ya están obligadas a hacer pública en virtud de otras normas, facilitando así el proceso de toma de decisiones a una amplia gama de inversores, incluidos expresamente los inversores minoristas. 

\subsubsection{Estructura actual: ¿Qué es y cómo funciona el ESAP?} 
Los colegisladores de la Unión alcanzaron un acuerdo sobre la creación del ESAP en mayo de 2023, que se formalizó como Reglamento (UE) 2023/2859. Se trata de un reglamento de aplicación directa que encomienda a la Autoridad Europea de Valores y Mercados la operación técnica de la plataforma. Conviene subrayar, porque es la clave de su naturaleza jurídica, que el ESAP no impone a las empresas europeas ningún requisito adicional de divulgación de información: no crea nueva obligación sustantiva de transparencia alguna, sino que se limita a facilitar el acceso a la información que esas empresas ya están legalmente obligadas a publicar en aplicación de las directivas y reglamentos pertinentes. Es, en ese sentido preciso, un instrumento puramente instrumental: no añade contenido, sino que reduce el coste de búsqueda y acceso a un contenido que ya existía de forma dispersa entre registros nacionales, páginas web corporativas y bases de datos sectoriales distintas entre sí. Por tanto, las atribuciones a la Autoridad Europea de Valores y Mercados no suponen ninguna nueva competencia normativa de esta autoridad sobre el contenido sustantivo de la información, cuya obligatoriedad sigue derivando exclusivamente de las normas sectoriales ya existentes que el ESAP se limita a agregar y a hacer accesible desde un único punto de entrada.

El ESAP no recopila la información directamente de las empresas obligadas a divulgarla, sino que opera mediante un sistema de dos niveles que conviene entender con precisión. 

\paragraph{Primer nivel: Organismo de Recopilación}
En un primer nivel, la información es recogida por lo que el Reglamento denomina un \textit{organismo de recopilación}, que puede ser la propia autoridad nacional competente del sector correspondiente, otro organismo o registro nacional ya existente, o una de las propias autoridades europeas de supervisión, entre ellas la Autoridad Bancaria Europea o la Autoridad Europea de Seguros y Pensiones de Jubilación además de la propia Autoridad Europea de Valores y Mercados. 

\paragraph{Segundo nivel: Divulgación}
En un segundo nivel, esa información recopilada se transmite y se hace pública a través de la plataforma ESAP propiamente dicha, operada centralizadamente por la Autoridad Europea de Valores y Mercados. 

Este diseño en dos niveles evita que las empresas obligadas a divulgar información deban realizar una doble notificación, una a su autoridad o registro nacional habitual y otra directamente al ESAP, preservando así el principio de que el instrumento no añade carga administrativa nueva sino que reorganiza el acceso a la carga ya existente.


\subsubsection{Estado actual: Calendario de implementación (en curso)} 
El Reglamento del ESAP exige a la Autoridad Europea de Valores y Mercados poner la plataforma a disposición del público a partir del verano de 2027, con una introducción progresiva del alcance de la información recogida, estructurada en tres fases sucesivas en función de la prioridad de las normas sectoriales afectadas. 

La primera fase, que comprende la información derivada de la Directiva de Transparencia, del Reglamento de Folletos que ya trabajamos en este epígrafe y del Reglamento sobre Ventas en Corto, ha comenzado formalmente el 10 de julio de 2026, aunque en esta fase inicial la recopilación de información por parte de los organismos de recopilación y de la propia Autoridad Europea de Valores y Mercados no implica todavía su publicación pública en la plataforma, que seguirá reservada hasta que esta esté operativa a mediados de 2027. 

La segunda fase, prevista para el 10 de enero de 2028, incorporará la información que reportan las entidades de crédito conforme a la Directiva de Requisitos de Capital que ya conoces del Código Normativo Único, así como la información de sostenibilidad publicada en aplicación del Reglamento de Divulgación de Finanzas Sostenibles, del Reglamento de Taxonomía y de la Directiva sobre Informes de Sostenibilidad Corporativa, e incluirá además a entidades de dinero electrónico sujetas a la Directiva de Servicios de Pago y a la Directiva de Dinero Electrónico. 

Una tercera fase, sujeta a una cláusula de revisión que la Comisión Europea debe completar, en estrecha cooperación con la propia Autoridad Europea de Valores y Mercados, no más tarde del 10 de enero de 2029 para evaluar la implementación, el funcionamiento y la efectividad del proyecto hasta ese momento, está prevista, si se confirma tras esa revisión, para enero de 2030.\footnote{%
    Este calendario dilatado, que se extiende desde la entrada en vigor del Reglamento en 2023 hasta una tercera fase que no se completaría antes de 2030, convierte al ESAP en el instrumento de todo este epígrafe con el periodo de implementación más largo, superando incluso a la reforma de titulización o al Listing Act en cuanto a la distancia temporal entre la aprobación de la norma y su plena operatividad.
}

Por tanto, el ESAP, pese a estar ya en fase de implementación activa desde julio de 2026, es todavía, de largo, el proyecto menos maduro de cuantos hemos desarrollado.



\clearpage
\section{Conclusión} 


\subsection{Recapitulación}


\subsection{Extensiones y relación con otras partes del Temario}



\subsection{Opinión}

\subsubsection{Idea final (Salida o cierra)}

\clearpage
\pagenumbering{gobble}
\MyBibliography



\MyBibItem{DAWID y GATTI}{Chapter 2 - Agent-Based Macroeconomics}{2018}{https://doi.org/10.1016/bs.hescom.2018.02.006}


% ANEXOS
\clearpage
\anexo{\textbf{Preguntas Test}}
%\input{\TemaCode/questions.tex} 



\clearpage
\anexo{Historia de la Unión Bancaria: Una introducción detallada}\label{anx:historia_ub}

\begin{specialblock}{Fuente original}
    Todo este texto ha sido extraído de: \href{https://www.bruegel.org/sites/default/files/2024-11/Nicolas%20book_online.pdf}{\textit{Europe’s banking union at ten: unfinished yet transformative} VÉRON, Bruegel (2024)}

    La unión bancaria es el proyecto de la Unión Europea para integrar la política del sector bancario en la zona del euro¹. Tras décadas de debates y decisiones trascendentales adoptadas en 2012, se convirtió en una realidad en 2014, quince años después del propio euro, cuando el Banco Central Europeo (BCE) asumió la autoridad como supervisor bancario. Por tanto, 2024 marca no solo el vigésimo quinto aniversario del euro, sino también el décimo aniversario de la supervisión bancaria europea, el principal logro inicial de la unión bancaria. Este volumen revisa la génesis de la unión bancaria, describe y evalúa su desarrollo hasta la fecha y analiza los elementos que aún faltan y las perspectivas para su eventual culminación.

    Existe un sólido argumento, corroborado por algunos de los principales participantes en los avances decisivos de 2012, de que la decisión de otorgar al BCE un mandato de supervisión bancaria fue determinante para hacer posible su actuación decisiva durante el verano de 2012, que puso fin a la fase más perturbadora de la crisis de la zona del euro. Solo esto aportó beneficios incalculables a los europeos, así como al sistema financiero mundial. El marco de supervisión microprudencial² resultante, centrado en el BCE, está esencialmente completo, proporcionando una protección mucho mayor de la zona del euro frente a la inestabilidad financiera que la existente durante la primera década y media de la unión monetaria. Hasta el momento, la supervisión bancaria europea parece haber sido eficaz, tanto en comparación con el régimen anterior de supervisión nacional en los países de la zona del euro como con sus homólogos internacionales, especialmente Estados Unidos, donde quedaron de manifiesto graves deficiencias en la supervisión bancaria durante la crisis de la banca regional de marzo de 2023.
    
    En el ámbito de la gestión y resolución de crisis, por el contrario, el progreso hacia un marco coherente ha sido vacilante y sigue siendo insuficiente. Las razones de ello incluyen los desafíos interrelacionados derivados de la concentración de las exposiciones soberanas, la reticencia a establecer un sistema europeo de garantía de depósitos y las importantes diferencias entre las preferencias nacionales en cuanto a los compromisos entre rescatar con dinero público a bancos en dificultades o imponer el reparto de las cargas financieras a los acreedores del sector privado.
    
    El marco supervisor de la unión bancaria, más allá de su efecto a corto plazo como facilitador de las decisiones críticas del BCE en el verano de 2012, hace que la zona del euro —y, por implicación, la UE— sea mucho más fuerte y resiliente de lo que era anteriormente. Sin embargo, el carácter incompleto del marco de gestión y resolución de crisis, que implica obstáculos para la plena integración del mercado bancario, conlleva costes significativos tanto en términos del potencial de crecimiento perdido de la UE como de fragilidad en determinados escenarios de crisis, aunque no genere una situación de emergencia inmediata. La probabilidad de avances concretos hacia la culminación de la unión bancaria parece lamentablemente baja a corto plazo, a menos que una crisis genere un renovado impulso político.
    
    No obstante, ninguno de los obstáculos para completar la unión bancaria es intrínsecamente insuperable, incluso en aquellos países de la UE donde parecen ser más poderosos. Inevitablemente, dada la amplitud y complejidad de la materia, no abordo de manera extensa ciertos temas importantes relacionados con la unión bancaria, aunque mi análisis y mis recomendaciones se formulan teniendo debidamente presentes dichos temas. Entre ellos se incluyen, aunque no se limitan a, los marcos de política macroprudencial, de provisión de liquidez de emergencia y de conducta financiera (incluida la protección de los clientes bancarios frente a la comercialización indebida de productos y la política de lucha contra el blanqueo de capitales, ambas con implicaciones prudenciales); los bancos malos; las sanciones financieras; la posibilidad de que los países de la UE que no pertenecen a la zona del euro puedan incorporarse voluntariamente a la unión bancaria;\footnote{%
        Esta cuestión fue abordada casi simultáneamente en diciembre de 2019 por informes de alto nivel en Dinamarca y Suecia. Sin embargo, desde entonces no se han producido avances significativos. Véanse Government of Sweden (2019) y el comunicado de prensa del Ministerio danés de Industria, Empresa y Asuntos Financieros de 19 de diciembre de 2019, «Report from the Working Group on possible Danish Participation in the Banking Union», https://www.eng.em.dk/news/2019/dec/report-from-the-working-group-on-possible-danish-participation-in-the-banking-union.
    } y la interacción entre la unión bancaria y la secuencia de acontecimientos que condujo a la salida del Reino Unido de la UE.\footnote{%
        Se ha sugerido la existencia de un vínculo causal entre la unión bancaria y el Brexit tanto ex ante (por ejemplo, Barker y Parker, 2012) como ex post (por ejemplo, Rogers, 2017).
    }
    
\end{specialblock}

La relación entre los bancos y los soberanos de Europa se remonta a mucho tiempo atrás. 

Se sustenta en múltiples legados históricos e institucionales que los observadores externos a menudo encuentran difíciles de comprender, y de los que muchos europeos tampoco son conscientes. La mayoría de los bancos europeos remontan sus orígenes al siglo XIX y principios del siglo XX, el período de surgimiento del nacionalismo en su forma moderna. Muchos fueron creados en un contexto de proyectos de desarrollo nacionales y/o coloniales, y ocasionalmente participaron en la financiación de guerras intraeuropeas. Un cierto número fue patrocinado, directa o indirectamente, por gobiernos locales o nacionales.\footnote{%
    En cambio, en Estados Unidos, la creación de bancos ha sido en gran medida un fenómeno ascendente y emprendedor (\textit{bottom-up}), aunque estuvo sometida a requisitos regulatorios generales en una etapa comparativamente temprana. Una coalición estable durante la mayor parte de los siglos XIX y XX limitó la capacidad de los bancos estadounidenses para expandirse por todo el país y ejercer influencia política a nivel federal (por ejemplo, CALOMIRIS y HABER, 2014). En las primeras décadas del país existía un nexo financiero entre los bancos y el gobierno a nivel de los distintos estados de Estados Unidos, pero ese vínculo directo desapareció en gran medida a lo largo del siglo XIX (GELPERN y VÉRON, 2018).
} 

A cambio de este patrocinio, los bancos europeos con frecuencia consideraban natural o inevitable facilitar la financiación del Estado, dirigir el crédito hacia empresas, sectores o proyectos favorecidos por el gobierno y respaldar las políticas y estrategias gubernamentales tanto en el ámbito nacional como en el extranjero. Las formas en que estas relaciones se configuraron y desarrollaron variaron ampliamente entre los países europeos. En muchos casos, ello implicó menores beneficios para los ahorradores y otros grupos de interés del sistema bancario, un patrón que los economistas denominan \textit{represión financiera}. A la inversa, \textit{nacionalismo bancario} es una expresión paraguas imprecisa, pero útil, para resumir la inclinación de los gobiernos a proteger y promover campeones bancarios nacionales en el contexto europeo cargado de historia. Aunque la relación entre los bancos europeos y los gobiernos durante el último siglo o los dos últimos siglos ha sido compleja y rara vez de plena alineación, en términos generales ha sido más estrecha que una relación mantenida a distancia.

$$\textbf{1. PRIMEROS DEBATES: INTEGRACIÓN DE LA POLÍTICA BANCARIA EUROPEA}$$
En las fases iniciales de la integración europea, los negociadores y legisladores no mostraron ningún interés en alterar estos legados nacionales. El Tratado de Roma (1957) incluía disposiciones específicas que exigían la unanimidad para las \textit{medidas relativas a la protección del ahorro, en particular la concesión de crédito y el ejercicio de la profesión bancaria} (art. $55$). Esto hizo más difícil armonizar la legislación bancaria que en la mayoría de los demás sectores de servicios, para los cuales la votación por mayoría cualificada era la norma (Maes, 2007, p. 27). La primera iniciativa de la Comisión Europea para la integración de la política bancaria, que comenzó en 1965, no llegó a contemplar una supervisión prudencial centralizada, y aun así se había agotado a principios de la década de 1970 (Mourlon-Druol, 2016). Una directiva de 1973 sobre la libertad de establecimiento en el sector bancario cambió poco, ya que no intentó abordar los divergentes requisitos regulatorios nacionales ni los obstáculos transfronterizos vinculados a los controles nacionales de capital.

Tras estas primeras experiencias, la Comisión Europea adoptó un enfoque decididamente incremental respecto a la integración del sector bancario, siendo su primer paso la denominada \textit{primera directiva bancaria} (1977). De manera significativa, los representantes de la industria bancaria tenían, en general, una visión menos favorable de la integración de las políticas a escala europea que sus homólogos de otros sectores, como el manufacturero. No emprendieron ninguna consolidación transfronteriza a gran escala, optando en su lugar, desde finales de la década de 1950 en adelante, por alianzas flexibles o clubes de bancos de distintos países, destinados a generar sinergias sin poner fundamentalmente en cuestión el principio de que cada banco conservaría su identidad nacional (DRACH, 2021). En un documento de posición de 1981, la Federación Bancaria Europea afirmó que \textit{las diferencias nacionales debían preservarse} (DRACH, 2020, p. 780). 

A finales de la década de 1980 y principios de la de 1990, la liberalización de los movimientos de capital y el programa del MIU de la Comisión Europea coincidieron con actitudes más favorables hacia la integración bancaria europea. Incluso entonces, sin embargo, la industria bancaria mostró poco entusiasmo por las iniciativas de armonización y prefirió defender opciones más flexibles de reconocimiento mutuo que preservaran las diferencias regulatorias nacionales y la protección de los campeones bancarios nacionales (DRACH, 2020).

Pasando de la regulación a la supervisión prudencial, se crearon diversos organismos paneuropeos para fomentar la coordinación transfronteriza. Siguiendo un patrón recurrente, cada una de estas iniciativas suscitó inicialmente esperanzas entre los funcionarios directamente implicados de que pudieran desembocar en un futuro marco supervisor integrado. Pero tales esperanzas, que solo se expresaban en privado, normalmente desaparecían tras las primeras reuniones, y \textbf{el principio de que la autoridad supervisora correspondía exclusivamente al nivel nacional permaneció firmemente defendido}. En primer lugar estuvo el Comité de Gobernadores, creado en 1964 por iniciativa de los ministros de Finanzas, en el que los gobernadores de los bancos centrales nacionales de la Comunidad Europea se reunían mensualmente en terreno neutral en Basilea. Este comité desempeñó posteriormente un papel decisivo en la creación del BCE, pero inicialmente no tuvo un impacto directo sobre la supervisión, aunque solo fuera porque esa función correspondía a los bancos centrales únicamente en tres de los seis países. En segundo lugar, y de forma separada, en 1972 las autoridades prudenciales de los Estados miembros tomaron una iniciativa propia para establecer un grupo de contacto que también se reunía regularmente, generalmente conocido por su nombre en francés, \textit{Groupe de Contact}. En tercer lugar, en 1978, la Comisión Europea promovió la creación de un Comité Consultivo Bancario (\textit{Banking Advisory Committee}, BAC), que reunía a altos funcionarios del banco central, la autoridad supervisora (cuando era distinta del primero) y el ministerio de Finanzas de cada país, celebrando normalmente una reunión por trimestre. En cuarto lugar, entre 1988 y 1990, frustrado por la reticencia de la Comisión Europea a otorgar mayores competencias al BAC, el Comité de Gobernadores estableció en su seno un Subcomité de Supervisión Bancaria (\textit{Banking Supervisory Sub-Committee}, BSSC), que se reunía antes que el BAC y contaba con una secretaría permanente proporcionada por el Banco de Pagos Internacionales en Basilea (MOURLON-DRUOL, 2025).

Las expectativas de una futura integración supervisora se manifestaron desde muy temprano, aunque únicamente de manera confidencial, con el entendimiento de que el asunto sería demasiado controvertido para plantearlo públicamente. A la inversa, desde muy pronto también se expresó explícitamente, a puerta cerrada, la oposición de principio de al menos algunos Estados miembros. Así, en diciembre de 1973, uno de los miembros del Groupe de Contact escribió a sus colegas: \textit{Aunque todavía pueda parecer bastante lejano, podría concebirse nuestro “Club” informal como el núcleo de una futura autoridad europea de supervisión bancaria} (GOODHART, 2011). Pero en 1974, el Banco de Francia dejó claro que el Groupe no debía recibir poderes delegados de supervisión bancaria (DRACH, 2019). Del mismo modo, en 1978, un funcionario de la Comisión Europea escribió en correspondencia interna que uno de los objetivos de la creación del BAC era \textit{proporcionar un precursor de una institución supervisora de la CEE}. Pero más tarde, en 1978, otro funcionario de la Comisión señaló, en el período previo a la reunión inaugural del BAC, que las \textit{principales personalidades del Ministerio [francés de Finanzas] han mantenido la mayor parte del tiempo una visión bastante escéptica de la coordinación bancaria de la CEE, al menos cuando se trata de cualquier decisión que implique una cierta transferencia de competencias supervisoras} (MOURLON-DRUOL, 2025). Siguiendo el mismo patrón, en mayo de 1990, el secretario del BSSC escribió al miembro del subcomité del Banco de Inglaterra que el BSSC era \textit{considerado como el precursor de la dirección encargada de la formulación de políticas de una División de Supervisión Bancaria del ECBS}, utilizándose entonces este último acrónimo para designar lo que se convertiría en el BCE (MOURLON-DRUOL, 2025). Pero la creación del BCE no contempló inicialmente una función sólida de supervisión bancaria dentro de la institución.

Durante estas primeras décadas, las conexiones entre un mercado único de servicios bancarios, la integración monetaria y la integración de la política bancaria fueron con frecuencia claramente identificadas desde un punto de vista analítico. Por ejemplo, el Comisario Europeo de Presupuesto y Control Financiero e Instituciones Financieras, Christopher Tugendhat, señaló en 1978 que \textit{sería más fácil aplicar una política monetaria común si tuviéramos sistemas bancarios supervisados con arreglo a normas comunes equivalentes y, eventualmente, de estructuras similares. Además, en la medida en que el Sistema Monetario Europeo [que se establecería en 1979] dé lugar a una mayor liberalización de los capitales, debemos asegurarnos de poder supervisar a los actores de este futuro mercado común de capitales, es decir, las entidades de crédito europeas, sobre la base de normas equivalentes}. En 1982, PADOA-SCHIOPPA, entonces director general de Asuntos Económicos y Financieros de la Comisión Europea, señaló con referencia al sistema financiero europeo: \textit{En la medida en que, por tanto, exista una necesidad de control y supervisión, surge la cuestión de determinar el nivel apropiado en el que deberían ejercerse y, a este respecto, existe un sólido argumento a favor de ejercerlos a nivel comunitario} (MOURLON-DRUOL, 2016). En 1991, en un momento en que las negociaciones sobre la unión monetaria ya estaban en marcha, pero cuyo resultado aún no se conocía, funcionarios estadounidenses, reflexionando sobre las enseñanzas de la crisis estadounidense de las asociaciones de ahorro y préstamo de la década anterior, identificaron lúcidamente el riesgo de una carrera supervisora hacia el mínimo común denominador derivada de la interacción entre la integración del MIU y el mantenimiento de la supervisión en el nivel nacional.\footnote{%
    Su contribución a un simposio celebrado en Berlín merece ser citada extensamente: \textit{Obviamente, las decisiones sobre dónde localizar las actividades bancarias dependen de muchos factores. Sin embargo, las diferencias en la supervisión podrían ser uno de los factores que afectaran a los bancos de la C[omunidad] E[uropea]. Los bancos podrían optar por establecer sus sedes centrales (o filiales con capitalización separada) en cualquiera de los 12 países en función de las percepciones sobre el rigor regulatorio (o la falta de este). Aunque nos resulta difícil juzgarlo desde el exterior, es posible que intereses regionales también pudieran dar lugar a lo que en Estados Unidos se ha denominado una “competencia en laxitud” entre los supervisores bancarios. Con el fin de estimular el desarrollo económico dentro de sus fronteras nacionales o de atraer empleos de sedes centrales en los servicios financieros, algunos países podrían sentirse inclinados a utilizar concesiones en el rigor o el coste de la supervisión como incentivo. [...] Esto podría tener un efecto grave sobre la seguridad y la solidez de los bancos en la CE si la convergencia regulatoria diera lugar a una supervisión cada vez más laxa} (SWAIM y WESSELS, 1991)
}

$$\textbf{1.1. Unión Monetaria: La cláusula habilitante del Tratado de Maastricht (1992)}$$
Aunque la unión monetaria europea había sido objeto de debate político al menos desde principios de la década de 1970, su verdadero punto de partida fue el denominado Informe del Comité Delors, presentado en abril de 1989. 

Sin embargo, ese texto fue claramente tímido en materia de supervisión bancaria. La recomendación del informe de que el Sistema Europeo de Bancos Centrales (SEBC) \textit{participaría en la coordinación de las políticas de supervisión bancaria de las autoridades supervisoras} (DELORS, 1989) implicaba que la supervisión seguiría ejerciéndose en el nivel nacional.\footnote{%
    Esto parece no haber obedecido a una falta de consideración de las cuestiones supervisoras en las deliberaciones del Comité, sino más bien a una falta de consenso sobre un enfoque más contundente. Los relatores del Comité, GUNTER BAER (funcionario alemán del Banco de Pagos Internacionales) y PADOA-SCHIOPPA, habían hecho referencia a las cuestiones supervisoras en uno de los dos documentos que redactaron antes de la primera reunión del Comité, sugiriendo que una moneda única debía ir asociada a una supervisión bancaria centralizada (VIANELLI, 2022). Durante ese trabajo preparatorio, el presidente del Bundesbank, PÖHL, defendió que se otorgara al futuro BCE un papel activo, aunque no decisorio, en la supervisión bancaria, similar al que desempeñaba el Bundesbank en Alemania: \textit{El BCE [...] debería estar estrechamente implicado en las actividades cotidianas de supervisión bancaria} (VAN DEN BERG, 2005). En enero de 1989, el gobernador del banco central de los Países Bajos, DUISENBERG, propuso que el futuro SEBC \textit{será responsable de la formulación de la política de supervisión bancaria a nivel comunitario y de la coordinación de las políticas de supervisión bancaria de las autoridades supervisoras nacionales}. Este texto se incluyó en una versión posterior del borrador del informe del Comité, pero fue atenuado hasta la formulación citada anteriormente en rondas posteriores (VAN DEN BERG, 2005). Un miembro del Comité, LAMFALUSSY, figuró entre quienes intentaron poner de relieve el desafío supervisor (MAES, 2017). Otro miembro del Comité Delors, el economista THYGESEN, mencionó la preocupación de algunos participantes, aparentemente los gobernadores de los bancos centrales de Alemania y posiblemente también de Dinamarca, de que otorgar al futuro BCE un mandato supervisor \textit{conduciría inevitablemente a una onerosa supervisión política y constituiría una amenaza para [su] autonomía} (ENDERLEIN y RUBIO, 2014). Jacques DE LAROSIÈRE, gobernador del Banco de Francia en aquel momento y miembro del Comité Delors, recordaría posteriormente: \textit{No dedicamos mucho tiempo a estas cuestiones} [de estabilidad financiera y supervisión bancaria],\textit{ también porque la supervisión bancaria podría haber sido una cuestión divisiva, ya que existían diferencias significativas en las responsabilidades en este ámbito entre los distintos bancos centrales} (MAES y PÉTERS, 2021)
} Sin embargo, el desafío supervisor fue retomado posteriormente por el Comité de Gobernadores y su Subcomité de Supervisión Bancaria (Banking Supervisory Sub-Committee, BSSC). Sobre la base del Informe Delors, el Consejo, en abril de 1990, encomendó efectivamente al Comité de Gobernadores la preparación de un proyecto de Estatuto para el futuro BCE. En noviembre de 1990, el Comité de Gobernadores, basándose en los trabajos del BSSC, acordó proponer que al BCE pudiera otorgársele explícitamente una función de supervisión bancaria: \textit{El BCE podrá formular, interpretar y aplicar políticas relativas a la supervisión prudencial de las entidades de crédito y otras instituciones financieras respecto de las cuales sea designado como autoridad supervisora competente}.

Sin embargo, a diferencia de sus bancos centrales, los gobiernos nacionales más influyentes no estaban dispuestos a renunciar a lo que consideraban un instrumento esencial para utilizar el sector bancario al servicio de sus objetivos de política; en otras palabras, la represión financiera. Mirando retrospectivamente al período fundacional del euro, Erkki Liikanen, entonces embajador de Finlandia ante las Comunidades Europeas, observó que \textit{la única parte que faltaba era que se debería haber contado con la unión bancaria, la supervisión bancaria a nivel europeo, la resolución bancaria y el seguro de depósitos. Eso se debatió en aquel momento. Pero los Estados miembros no lo querían; afirmaban: nosotros podemos gestionar nuestros bancos, sabemos que están en buena situación}. En una clara referencia al control directo ejercido por el Ministerio alemán de Finanzas sobre su autoridad supervisora bancaria, el funcionario de la Comisión Europea Alexander Italianer, que posteriormente se convertiría en Secretario General de la Comisión, señaló lacónicamente que una responsabilidad a nivel europeo en materia de estabilidad financiera \textit{también significaría que algunos ministerios de Finanzas verían transferidas al BCE sus competencias prudenciales. Esto explica presumiblemente en parte las dificultades para alcanzar un acuerdo sobre el papel del BCE en el ámbito prudencial} (ITALIANER, 1993).

En una fase temprana del proceso, en agosto de 1989, el Tesoro francés expresó su posición de que cualquier función supervisora para el futuro BCE \textit{debía mantenerse al margen de la negociación}. En mayo de 1991, durante la primera ronda de discusión formal del proyecto de Estatuto elaborado por el Comité de Gobernadores, el Tesoro francés volvió a liderar la oposición a dicha función, con el apoyo de sus homólogos británico y alemán y pese al asesoramiento discrepante del Banco de Francia; los Estados miembros más pequeños que favorecían la supervisión por parte del BCE no formaron una coalición consistente y, en cualquier caso, no podían igualar el peso político combinado de los tres países más grandes.\footnote{%
    Durante una importante reunión celebrada el 6 de noviembre de 1991, en la que se examinó la propuesta de la Presidencia neerlandesa relativa a la cláusula habilitante, Francia, Alemania, Luxemburgo y el Reino Unido votaron a favor de eliminar completamente el texto propuesto, mientras que Portugal, Italia, España y Dinamarca votaron por mantenerlo o reforzarlo (VAN DEN BERG, 2005). La Comisión Europea se mostró mayoritariamente contraria a atribuir una función supervisora al BCE, aparentemente por razones burocráticas. En una nota de la Comisión de junio de 1990 sobre esta cuestión, se argumentaba que \textit{no existe razón alguna para suponer que una supervisión centralizada por parte del Eurofed} [sic] \textit{vaya a ser más eficiente o menos propensa al fracaso que la supervisión ejercida por las autoridades bancarias nacionales} (MOURLON-DRUOL, 2025).
} 

No obstante, el equipo neerlandés que asumió la presidencia rotatoria del Consejo durante el último semestre de la negociación (la segunda mitad de 1991) logró forjar un acuerdo sobre una \textbf{cláusula habilitante} que preservaba la posibilidad de una futura función supervisora para el BCE, aunque sometida a una exigente condición de unanimidad. Los ministros aprobaron dicha cláusula en sus reuniones de los días 1 a 3 de diciembre de 1991. Como consecuencia, la supervisión prudencial no tuvo que debatirse en la reunión final del Consejo Europeo que ultimó el texto del tratado los días 9 y 10 de diciembre de 1991 (MOURLON-DRUOL, 2025), celebrada en la ciudad neerlandesa de Maastricht, de la que el tratado tomó su nombre. La cláusula habilitante quedó así consagrada en el artículo $105(6)$ del Tratado de Maastricht. Tras su posterior revisión, dicha cláusula pasó a denominarse desde 2009 artículo $127(6)$ del TFUE.\footnote{%
    Establecía lo siguiente: \textit{El Consejo podrá, por unanimidad, a propuesta de la Comisión y previa consulta al BCE y previo dictamen conforme del Parlamento Europeo, encomendar al BCE tareas específicas respecto de políticas relacionadas con la supervisión prudencial de las entidades de crédito y otras instituciones financieras, con excepción de las empresas de seguros}. Como veremos, a mediados de 2012 se consideró que esta compleja redacción constituía una base jurídica suficiente para establecer, en esencia, las mismas disposiciones que el Comité de Gobernadores había propuesto en su proyecto de Estatuto del BCE de finales de 1990, en lo relativo a la autoridad supervisora efectiva, aunque no respecto de la formulación de políticas.
}

En el período inmediatamente posterior a la adopción del Tratado de Maastricht, la ausencia de un marco integrado de política bancaria que acompañara a la unificación monetaria no pasó desapercibida para los observadores. Escribiendo en 1992, el veterano historiador financiero KINDLEBERGER citó \textit{la regulación de los bancos como la primera de «una serie de cuestiones que debían resolverse en el camino hacia la integración monetaria y financiera europea} (KINDLEBERGER, 1993). La jurista Rosa Lastra observó: \textit{la ausencia de una función de supervisión bancaria claramente definida para el SEBC} [Sistema Europeo de Bancos Centrales], \textit{existiendo ya una “licencia bancaria única”, podría tener efectos perturbadores en caso de producirse una crisis sistémica} (LASTRA, 1992). El economista financiero Alberto GIOVANNINI escribió: \textit{cabe sospechar que una unión monetaria incompleta, no acompañada de reformas de instituciones como los sistemas de pagos y la supervisión bancaria, podría dar lugar a costes de eficiencia que fácilmente compensarían los beneficios estimados de la introducción de una moneda única} (GIOVANNINI, 1993). EICHENGREEN predijo: \textit{Las autoridades} [supervisoras] \textit{nacionales se verán presionadas para conceder ventajas regulatorias a los bancos nacionales} […] \textit{Pero muchos de los costes de una desregulación competitiva, en forma de inestabilidad financiera, serán soportados por la} [Comunidad] \textit{Europea en su conjunto} (EICHENGREEN, 1993). Sin embargo, estas lúcidas advertencias quedaron olvidadas durante los años posteriores.

$$\textbf{1.2. De Maastricht a la crisis financiera de 2007}$$
Los vínculos entre la integración del mercado bancario, la unificación monetaria y la política del sector bancario fueron debatidos extensamente en las décadas de 1960, 1970 y 1980 y a principios de la década de 1990. En cambio, estos vínculos apenas figuraron en los numerosos debates sobre el diseño de la UEM europea y sus perspectivas entre el período inmediatamente posterior al Tratado de Maastricht y el estallido de la gran crisis financiera en 2007. 

Un examen de las opiniones, en gran medida escépticas, de los economistas estadounidenses sobre el euro entre 1989 y 2002 mencionó las preocupaciones relativas al sector bancario una sola vez (JONUNG y DREA, 2009). Se trataba de una polémica de 1999 que pronosticaba el inminente colapso del euro y que, de paso, identificaba los futuros rescates del sector bancario como uno de los diversos factores impulsores de futuros déficits públicos insostenibles en Francia e Italia, basándose en los casos entonces recientes del Crédit Lyonnais de Francia en 1994-1995 y del Banco di Napoli de Italia en 1995 (CALOMIRIS, 1999). En otras palabras, durante ese período hubo muchas predicciones de desastre para el euro, pero el círculo vicioso entre bancos y soberanos que casi quebró la unión monetaria en 2011-2012 no figuraba entre los puntos débiles identificados.\footnote{%
    El escenario de CALOMIRIS, además de errar en el pronóstico, identificó un componente de la dinámica entre bancos y soberanos, pero no el círculo completo, puesto que su análisis no contemplaba que una mayor tensión fiscal socavaría, a su vez, la solidez del sector bancario. Una de las muchas otras predicciones de colapso de la zona del euro publicadas en aquel momento se basaba en un escenario en el que al BCE se le concedería prematuramente autoridad supervisora en virtud de la cláusula habilitante del Tratado de Maastricht en 1999 y se encontraría incapaz de ejercer esta autoridad con tanta eficacia como lo habían hecho los supervisores nacionales, contribuyendo así a la acumulación de riesgo financiero en el sistema bancario de la zona del euro (LASCELLES, 1996). Este escenario es exactamente lo contrario de lo que realmente ocurrió en la década de 2000.
}

Al acercarse el final de su mandato como presidente del BCE en octubre de 2019, DRAGHI reflexionó que \textit{de forma no muy distinta de la crisis ecológica, la crisis de la zona del euro ha puesto al descubierto múltiples bucles de retroalimentación que anteriormente no se comprendían bien, por ejemplo, entre soberanos, bancos y empresas} (DRAGHI, 2019). En lo que respecta a la evolución de las políticas, hubo algunos esfuerzos encaminados a una mayor armonización del marco regulatorio bancario, pero adoptaron la forma de directivas, en contraposición a reglamentos, como la Directiva bancaria de 2000 y la Directiva de requisitos de capital de 2006, lo que significaba que estos textos todavía debían ser transpuestos al ordenamiento legislativo de cada Estado miembro, con las consiguientes diferencias entre los países de la UE.\footnote{%
    Respectivamente, la Directiva 2000/12/CE, de 20 de marzo de 2000, y la Directiva 2006/49/CE, de 14 de junio de 2006.
}


Durante ese período también se produjeron algunos cambios incrementales en la arquitectura supervisora de la Unión, pero sin alterar fundamentalmente el principio de supervisión nacional. Cuando el BCE comenzó sus operaciones en 1998 en Fráncfort, estableció un Comité de Supervisión Bancaria interno, compuesto por supervisores bancarios nacionales y banqueros centrales de los países de la zona del euro, con el objetivo de facilitar la cooperación transfronteriza y el intercambio de información supervisora pertinente. Del mismo modo que el BCE había sucedido (indirectamente) al Comité de Gobernadores, su Comité de Supervisión Bancaria era el heredero del BSSC del Comité. Posteriormente, el BCE intentó influir en las decisiones relativas a las arquitecturas supervisoras nacionales de los distintos países de la UE, argumentando que debía otorgarse a los bancos centrales nacionales autoridad de supervisión prudencial sobre los bancos para lograr sinergias con sus funciones esenciales de vigilancia de la infraestructura financiera, seguimiento macroprudencial y operaciones de liquidez en caso de crisis. El BCE añadió que tales disposiciones reforzarían la independencia y la profesionalidad de la supervisión bancaria y la gestión del riesgo sistémico, y restó importancia a los contraargumentos basados en consideraciones sobre posibles conflictos de intereses, riesgo moral y acumulación excesiva de poder por parte de los bancos centrales (BCE, 2001). Esta posición fue impulsada, entre otros, por Tommaso Padoa-Schioppa, que se incorporó al BCE como miembro del Comité Ejecutivo desde la creación de la institución en junio de 1998 (BINI SMAGHI, 2011; DE RYNCK, 2014). PADOA-SCHIOPPA fue también uno de los primeros defensores de la armonización plena de la legislación de la UE sobre los requisitos prudenciales bancarios, para la cual acuñó la expresión \textbf{código normativo único} (Padoa-Schioppa, 2004).

Sin embargo, el BCE fracasó claramente en su labor de presión. Esto quedó ilustrado de la manera más evidente por la creación en 2002 de BaFin (Bundesanstalt für Finanzdienstleistungsaufsicht) como autoridad integrada de supervisión financiera de Alemania, que, al igual que su predecesora, quedó sometida a la estrecha supervisión del Ministerio de Finanzas, pero separada del Bundesbank.\footnote{%
    Cabe señalar que, en las discusiones que finalmente condujeron a la creación de BaFin, el Bundesbank argumentó que debía otorgársele autoridad de supervisión bancaria de conformidad con las recomendaciones del BCE de aquel momento, pero fue desautorizado en el proceso político de toma de decisiones (Schüler, 2004, página 12).
}

En 2004, la UE creó un Comité de Supervisores Bancarios Europeos (CEBS) autónomo, que era independiente del BCE y parecía consolidar el enfoque alternativo según el cual los supervisores se mantenían separados de los bancos centrales, como también había ocurrido desde hacía tiempo en Bélgica, Suiza y Escandinavia y, más recientemente, en el Reino Unido con la creación de la Autoridad de Servicios Financieros a finales de 2001. El Groupe de Contact quedó integrado en el CEBS, al que se dotó de una secretaría permanente en Londres. En la práctica, el CEBS volvió en gran medida redundante al Comité de Supervisión Bancaria del BCE —con el que compartía varios miembros, concretamente aquellos bancos centrales que eran supervisores bancarios—, al menos en términos de supervisión microprudencial. El CEBS organizaba consultas periódicas entre los supervisores bancarios nacionales y trabajos conjuntos sobre cuestiones de interés común, pero sin la autoridad ni los medios para imponer la coherencia supervisora. Mientras tanto, el Comité de Supervisión Bancaria del BCE promovió intentos de simulación de crisis y pruebas de resistencia, pero estas iniciativas encontraron una oposición considerable y, en última instancia, fueron incapaces de generar un análisis compartido realista de las vulnerabilidades del sistema.\footnote{%
    Se celebraron dos «ejercicios de pruebas de resistencia del Eurosistema», en abril de 2005 y mayo de 2006, respectivamente, y fueron seguidos por una conferencia a mediados de julio de 2007, poco antes del comienzo de la gran crisis financiera (véase https://www.ecb.europa.eu/pub/conferences/html/sfi_conf.it.html). Se describen de manera somera en las páginas 204-206 de las actas publicadas de la conferencia (BCE, 2008).
}

Durante todo ese período, se consideraba comúnmente que cualquier integración de políticas en el ámbito de los servicios financieros tendría lugar primero en el ámbito de los mercados mayoristas y la regulación de valores, y no en el de la banca, puesto que esta última mantenía vínculos tan estrechos con los gobiernos y las políticas nacionales. Esta priorización quedó reflejada, por ejemplo, en el plan de acción de servicios financieros de la Comisión Europea de 1999, un programa para la actividad legislativa de los años siguientes en ese ámbito (Comisión Europea, 1999), y en el énfasis puesto en la supervisión de los mercados de valores en un destacado informe de un comité presidido por el banquero central Alexandre Lamfalussy (Lamfalussy, 2001).\footnote{%
    El propio CEBS se inspiró en el Comité de Reguladores Europeos de Valores, creado en 2001 conforme a la recomendación del informe Lamfalussy. Entretanto, Lamfalussy había dirigido el Instituto Monetario Europeo, que había sucedido al Comité de Gobernadores en enero de 1994 y dio paso al BCE en junio de 1998.
}

Se siguieron formulando diagnósticos sobre las deficiencias de la arquitectura de política bancaria de la zona del euro, pero las advertencias no llegaron a convertirse en un consenso de política europea. En un análisis en profundidad realizado en 1998, el Fondo Monetario Internacional (FMI) criticó la ausencia de una función explícita de prestamista de última instancia para el BCE y pidió una mayor integración de las políticas: «Con el tiempo, la introducción del euro […] puede requerir la centralización de la vigilancia financiera, la gestión del riesgo sistémico y la resolución de crisis. […] el BCE estará en el centro de los mercados financieros europeos sin las herramientas necesarias para evaluar de manera independiente la solvencia de las contrapartes ni las herramientas para prestar apoyo directo a instituciones solventes pero ilíquidas. No es probable que esto sea sostenible, y el BCE puede verse pronto obligado a asumir un papel rector y coordinador en la gestión de crisis y la supervisión bancaria» (FMI, 1998, páginas 106 y 110).

En un discurso destacado pronunciado poco después, Padoa-Schioppa rechazó las preocupaciones del FMI sobre la gestión de crisis basándose en que el BCE estaría preparado para responder a las necesidades de liquidez de emergencia, una afirmación que quedó ampliamente validada cuando la crisis estalló una década después. No obstante, coincidió en que era «absolutamente necesario» avanzar hacia una mayor integración supervisora y «permitir que surja una especie de supervisor colectivo de la zona del euro que pueda actuar con tanta eficacia como si existiera un único supervisor». Posiblemente por consideraciones diplomáticas, expresó simultáneamente su escepticismo respecto de que dicha integración tuviera que implicar la activación de la cláusula habilitante del Tratado de Maastricht: «Aunque el Tratado contiene una disposición que permite la asignación de tareas supervisoras al BCE, personalmente no confío en el supuesto de que esta cláusula vaya a ser activada» (Padoa-Schioppa, 1999).

Sin embargo, ni el FMI ni Padoa-Schioppa ni otros analistas de la época llegaron a contemplar un escenario en el que los fallos de supervisión fueran tales que las preocupaciones sobre la solvencia bancaria amenazaran la solvencia crediticia soberana de uno o varios Estados miembros. Del mismo modo, el economista Paul De Grauwe, en un análisis por lo demás singularmente lúcido de los riesgos de inestabilidad financiera sistémica en la unión monetaria que estaba a punto de constituirse, no contempló explícitamente un escenario de tensión crediticia soberana, y mucho menos de ruptura de la zona del euro.\footnote{
Paul De Grauwe, «The euro and financial crises», Financial Times, 20 de febrero de 1998.
}

Los inversores tampoco incorporaron en los precios la posibilidad de futuros problemas de crédito soberano, quizá en parte cegados por los incentivos efectivos a la disciplina fiscal durante la década de 1990 entre los países que querían incorporarse al euro desde su inicio.\footnote{%
    Entre los países que inicialmente decidieron adoptar el euro, únicamente Grecia tuvo que afrontar un retraso, incorporándose a la unión monetaria con dos años de demora, el 1 de enero de 2001. Sin embargo, las preocupaciones sobre la durabilidad de esa disciplina fiscal se hicieron evidentes ya en 2003, cuando Francia y Alemania neutralizaron conjuntamente los esfuerzos de la Comisión Europea por aplicar el procedimiento de déficit excesivo establecido por el Tratado de Maastricht.
}

Los diferenciales de deuda soberana de todos los futuros países de la zona del euro —es decir, la diferencia entre el tipo de interés determinado por el mercado de su deuda a largo plazo, por ejemplo con vencimiento a diez años, y el del emisor de referencia, a saber, Alemania— descendieron a niveles extremadamente bajos durante la segunda mitad de la década de 1990 y permanecieron en ellos hasta 2008, un año después del inicio de la gran crisis financiera. La mayoría de los participantes en el mercado y muchos responsables de política habían adoptado la creencia de que había surgido un mercado único de servicios financieros, en el cual las condiciones financieras se habían igualado entre todos los países de la zona del euro. La posibilidad de fragmentación de ese sistema financiero siguiendo líneas nacionales, mediante el aumento de los diferenciales por país hasta niveles significativos, parecía inverosímil para la mayoría, pese a las disposiciones consagradas en los tratados según las cuales los países de la zona del euro no podían mutualizar su deuda existente ni siquiera en casos de turbulencia financiera.

Mientras tanto, las autoridades nacionales europeas permitieron que los bancos bajo su supervisión ampliaran sus balances a un ritmo acelerado desde mediados de la década de 1990, si no antes (Bayoumi, 2017). Como resumió el historiador financiero Harold James (2012b, página 393): «Los primeros diez años de existencia del euro estuvieron ensombrecidos por dos sagas de larga duración, ambas objeto de gran atención pública y que parecían definir la lucha en torno a la moneda: las disputas sobre el carácter de la política monetaria y quién la elaboraba, y las disputas sobre la competitividad. Los debates desviaron la atención de las cuestiones principales, la insostenibilidad del enfoque de Europa —y del mundo— respecto de la banca y los efectos de la explosión de la financiarización sobre el crédito público».

Solo la experiencia real y dolorosa de la crisis revelaría el nexo de riesgo entre los sistemas bancarios y las finanzas soberanas como la mayor vulnerabilidad práctica de la zona del euro.\footnote{%
    Esta opinión se ha convertido en estándar desde la crisis. Como un ejemplo entre muchos, Charles Grant (2015) citó «cinco graves defectos de diseño» en el diseño inicial de la zona del euro, el segundo de los cuales era «que los planes para la unión monetaria carecían de disposiciones para una “unión bancaria”, que ahora se reconoce como un componente esencial».
}

$$\textbf{2. LA CRISIS DE LA ZONA EURO}$$
El 30 de julio de 2007, IKB Deutsche Industriebank, un banco alemán de tamaño medio, anunció su rescate por parte de la institución financiera pública KfW, convirtiéndose así en la primera víctima bancaria de lo que pronto se transformaría en la gran crisis financiera. La crisis se intensificó en los meses posteriores, incluida la corrida bancaria de Northern Rock en el Reino Unido a mediados de septiembre de 2007, y alcanzó toda su fuerza en septiembre-octubre de 2008 con una sucesión de acontecimientos traumáticos que se produjeron en rápida secuencia, siendo el más intensamente recordado la quiebra de Lehman Brothers el 15 de septiembre de 2008.

A partir de 2009-2010, los acontecimientos se transformaron en una crisis de la zona del euro. Otras jurisdicciones en las que se habían producido turbulencias durante los dos primeros años de la crisis —Estados Unidos, el Reino Unido, Suiza y Dinamarca— retornaron a la estabilidad financiera. La crisis de la zona del euro alcanzó su punto culminante en el período comprendido entre el verano de 2011 y el verano de 2012, con secuelas dramáticas pero comparativamente limitadas en Chipre en 2013 y Grecia en 2015, y había concluido en lo esencial a mediados de 2017.

Aquí solo se mencionan algunos puntos destacados por haber allanado el camino hacia la unión bancaria, seleccionados de una secuencia compleja de acontecimientos que varios autores han descrito y analizado de manera más exhaustiva (por ejemplo, Bastasin, 2015; Bayoumi, 2017; Tooze, 2018; Rehn, 2020).


$$\textbf{2.1. Fracaso supervisor y círculo vicioso entre bancos y soberanos}$$

En la mayoría de las narrativas mediáticas y percepciones públicas, el componente fiscal estuvo en el centro de la crisis de la zona del euro, por lo que a menudo se la denomina «crisis de deuda soberana del euro». La crisis también reveló debilidades estructurales en varios países, especialmente en Grecia. Sin embargo, en gran medida, la historia subyacente fue una de fragilidad del sector bancario, al menos tanto como de incontinencia presupuestaria o rigideces estructurales. De los denominados países en crisis dentro de la zona del euro, el elemento fiscal fue inequívocamente dominante únicamente en Grecia. En Irlanda, España, Chipre y Eslovenia, el sector bancario fue el verdadero epicentro de las turbulencias, siendo Portugal un caso intermedio\textsuperscript{31}. Aún más importante, las vulnerabilidades del sector bancario en Francia, Alemania e Italia influyeron en las posiciones de muchos de estos países de mayor tamaño y, como consecuencia, también en las decisiones de la UE.

Considerada desde el punto de vista de la UE, por tanto, una característica central de la crisis fue el fracaso colosal y casi universal de la regulación y la supervisión bancarias. En principio, estos dos tipos de fracaso —regulatorio y supervisor— no están necesariamente correlacionados. Es posible imaginar que un entorno regulatorio laxo sea compensado por una supervisión rigurosa, porque los supervisores suelen estar facultados para imponer requisitos más exigentes que los mínimos establecidos en la regulación bancaria, con el fin de garantizar que los bancos sigan siendo seguros y sólidos, una capacidad a menudo denominada examen supervisor o discrecionalidad supervisora.\footnote{
La discrecionalidad supervisora está incorporada en los denominados acuerdos de Basilea, que definen el marco internacional para la supervisión prudencial de los bancos. Estos acuerdos toman su nombre del Comité de Supervisión Bancaria de Basilea, que los establece y que, a su vez, tiene su sede en el Banco de Pagos Internacionales de Basilea, Suiza. Distinguen entre requisitos mínimos uniformes del «pilar 1» derivados de la regulación aplicable y requisitos específicos para cada banco del «pilar 2» derivados del examen supervisor. Un «pilar 3» adicional gira en torno a requisitos de divulgación destinados a generar presión por parte de los inversores en las acciones —cuando cotizan— y los bonos de los bancos, conocida en el lenguaje de Basilea como disciplina de mercado.
}

A la inversa, pueden producirse fallos supervisores incluso en presencia de una regulación adecuada si esta no es aplicada correctamente por las autoridades supervisoras. Dado que una supervisión sólida puede compensar parcialmente una regulación laxa, mientras que una regulación sólida no puede compensar una supervisión deficiente, el énfasis aquí se pone generalmente más en la supervisión que en la regulación. En la práctica, los fallos regulatorios y supervisores suelen estar correlacionados, especialmente cuando ambos derivan, al menos en parte, de un entorno en el que las autoridades públicas son propensas a conceder a los bancos un trato más generoso del que exigiría el interés público.

En el período previo a la crisis, todos los bancos europeos eran considerados por sus supervisores nacionales conformes con los requisitos de capital aplicables. Pero estos requisitos eran demasiado bajos y la aplicación del cumplimiento era demasiado laxa, con el resultado de que la mayoría de los bancos estaban gravemente infracapitalizados según cualquier criterio económico razonable. De los 13 países de la zona del euro al comienzo de la crisis, únicamente Finlandia no experimentó ningún caso de gran notoriedad pública de incumplimiento flagrante de la supervisión prudencial nacional que desembocara en un rescate gestionado por el Gobierno a costa de los contribuyentes durante la década comprendida entre 2007 y 2017.\footnote{%
    No hubo ningún caso independiente de fracaso supervisor grave en Luxemburgo, pero ese país estuvo implicado, junto con Bélgica, Francia y los Países Bajos, en los colapsos de gran notoriedad de Dexia y Fortis, como se señala, por ejemplo, en la Evaluación de la Estabilidad del Sistema Financiero del FMI de 2011 (FMI, 2011). Chipre, Malta, Eslovaquia y los países bálticos, que se incorporaron a la zona del euro después del comienzo de la crisis, quedan excluidos de este cómputo, aunque la observación de un fracaso supervisor generalizado también resulta aplicable, al menos, a Chipre y Letonia entre ellos.
}

Este pobre historial se hizo manifiestamente visible al comienzo de la secuencia de la crisis; un destacado informe de la UE publicado en febrero de 2009 —y analizado más adelante— señaló con severidad que «las pruebas muestran claramente que la función de prevención de crisis de los supervisores en la UE no se ha desempeñado adecuadamente y no es apta para su finalidad» (Larosière, 2009, p. 42). La experiencia europea también fue marcadamente distinta de la de Estados Unidos, donde los problemas más dramáticos se produjeron en «entidades no bancarias» sometidas a regímenes de supervisión más débiles: compañías de financiación al consumo como Household International, intermediarios bursátiles como Bear Stearns y Lehman Brothers, instituciones de ahorro como IndyMac y Washington Mutual, aseguradoras como AIG y empresas patrocinadas por el Gobierno como Fannie Mae y Freddie Mac. La supervisión bancaria estadounidense distó mucho de ser impecable, pero no fracasó de manera ni remotamente tan generalizada como en Europa, aunque la espiral descendente de la crisis acabó amenazando también la viabilidad de muchos bancos estadounidenses.\footnote{%
    Citigroup ha sido citado ampliamente como un caso de fallo supervisor estadounidense. El presidente de la Reserva Federal, Ben Bernanke, declaró en 2009 que «de quizá las 13, 13 instituciones financieras más importantes de Estados Unidos, 12 corrían riesgo de quebrar en un período de una o dos semanas» en el punto culminante de la crisis en septiembre-octubre de 2008 (Financial Crisis Inquiry Commission, 2011, p. 354).
}

Las autoridades estadounidenses también estaban, en general, mejor preparadas que sus homólogas europeas para intervenir en casos de quiebra bancaria.\footnote{%
    En tres ocasiones, las autoridades estadounidenses activaron el mecanismo jurídico conocido como excepción por riesgo sistémico, que permitía un apoyo selectivo a bancos frágiles que, de otro modo, podrían quebrar: para Wachovia el 29 de septiembre de 2008, Citigroup el 23 de noviembre de 2008 y Bank of America el 16 de enero de 2009. Wachovia fue adquirido posteriormente por Wells Fargo en una operación que eliminó la necesidad de apoyo público. Citigroup y Bank of America recibieron ambos una inyección de acciones preferentes y se beneficiaron de una garantía pública de activos (FDIC, 2017, capítulo 3).
}

En suma, la zona del euro experimentó un fracaso sorprendentemente general de su régimen de supervisión bancaria, basado casi exclusivamente en la supervisión nacional. La explicación más directa es que la combinación de la integración de los mercados y monetaria, por una parte, y la conservación de una supervisión prudencial casi exclusivamente nacional, por otra, distorsionó los incentivos de los supervisores hacia el nacionalismo bancario, en contraposición a su mandato aparente de garantizar la seguridad y la solidez.\footnote{%
    En algunos casos, el «nacionalismo» operó en el nivel subnacional, por ejemplo, en las comunidades autónomas de España, las provincias y municipios de Italia o los \textit{Länder} de Alemania. Véanse, por ejemplo, Otero-Iglesias et al. (2016) sobre España y Hellwig (2018) sobre Alemania.
}

Una ilustración vívida de esa deriva fue proporcionada por la adquisición en 2007 de ABN AMRO, el grupo bancario internacional con sede en Ámsterdam, por un consorcio que incluía al Royal Bank of Scotland (RBS) británico, al Fortis del Benelux y al Santander español, que poco después revendió de manera rentable Banca Antonveneta, un banco italiano que ABN AMRO acababa de adquirir en 2006, a Banca Monte dei Paschi di Siena. Particularmente en el Reino Unido, Bélgica e Italia, hubo un examen muy insuficiente de las respectivas operaciones, en parte porque la expansión de los campeones locales se consideraba intrínsecamente deseable. El legado de la adquisición de ABN AMRO desempeñó un papel significativo en la posterior caída de los tres adquirentes: RBS y Fortis en 2008, y Monte dei Paschi gradualmente en años posteriores. A la inversa, las autoridades nacionales actuaron en muchos casos para proteger a los bancos nacionales frente a su adquisición por homólogos extranjeros, utilizando su autoridad legal para examinar los cambios en el control accionarial mucho más allá de lo que habría estado justificado por consideraciones prudenciales.\footnote{%
    Véase, por ejemplo, Xavier Vives, «European Banks Future on the Urge to Merge», Wall Street Journal, 13 de mayo de 2005, https://blog.iese.edu/xvives/files/2011/09/208.pdf.
}

Junto con la supervisión ineficaz, el otro legado —relacionado con ella— con el que Europa entró en la gran crisis financiera fue una elevada disposición de los gobiernos a utilizar sus propios recursos —el dinero de los contribuyentes— para compensar las pérdidas sufridas por los participantes en el mercado tras las quiebras bancarias o, en términos coloquiales, rescatarlos. Simplificando en exceso una historia bastante compleja, este enfoque de «bolsillos profundos» para la resolución de crisis bancarias se convirtió en una práctica estándar en Europa durante el período de entreguerras, siendo el Creditanstalt de Austria, el Danat-Bank, el Dresdner Bank y el Commerzbank de Alemania, y los tres bancos más grandes de Italia —Banca Commerciale Italiana, Credito Italiano y Banco di Roma— algunos de los casos con mayor repercusión internacional. A partir de entonces, los casos en los que los acreedores de bancos en quiebra perdieron dinero, como Bankhaus Herstatt en 1974 y Bank of Credit and Commerce International en 1991, fueron excepciones más que la regla, y el hecho de que los titulares de derechos perdieran dinero se consideraba normalmente un fracaso de las políticas.\footnote{%
    La crisis de Herstatt dio lugar a varias iniciativas para crear elementos de una red de seguridad y de reducción del riesgo, incluida la creación de un seguro de depósitos (\textit{Einlagesicherungsfonds}) y de una institución especializada para proporcionar liquidez (\textit{Liquiditäts-Konsortialbank}) por parte de los bancos alemanes del sector privado, seguida de la creación del sistema de protección institucional de las cajas de ahorros alemanas en 1975; la fundación del Comité de Supervisión Bancaria de Basilea (Goodhart, 2011); y, tras mucho retraso, la creación del sistema CLS, que comenzó a operar en 2002 para reducir el riesgo en las operaciones de divisas.
}

En Europa, por tanto, al igual que en otras partes del mundo, incluidos Canadá y Japón, se entendía que una buena práctica de gestión de crisis se reducía a una supervisión eficaz para que ningún banco quebrara —prevención de crisis— o, si el colapso era inevitable, a una intervención temprana para minimizar el coste del rescate para los contribuyentes y liquidar los activos problemáticos de la manera más eficiente posible. En esta línea, la crisis bancaria sueca de 1991-1992, en la que el Gobierno reembolsó a todos los acreedores de los bancos quebrados, fue considerada generalmente un éxito de política. A la inversa, en casos de rescates a gran escala, como el de Crédit Lyonnais en 1993, la indignación pública se dirigió hacia la mala gestión del banco y el fracaso de la supervisión pública, y no hacia el principio del uso correctivo del dinero de los contribuyentes, que se consideraba inevitable una vez que las pérdidas se habían materializado.

Así, en las primeras fases de la crisis iniciada en 2007, se consideraba natural que un banco en quiebra fuera rescatado por el Gobierno de su país de origen. En varios casos, incluidos IKB y RBS, incluso los accionistas fueron compensados en una medida significativa con fondos públicos; todos los acreedores, incluidos los subordinados, fueron reembolsados íntegramente en todos estos primeros casos. Un caso extremo fue Irlanda, donde el Gobierno garantizó formalmente a finales de septiembre de 2008 todos los depósitos bancarios y casi toda la deuda bancaria, partiendo de la premisa de que todos los bancos eran solventes, lo que posteriormente resultó no ser cierto.\footnote{%
    Dos años después, Irlanda también proporcionaría el primer caso durante la crisis de la zona del euro en el que los acreedores bancarios subordinados sufrieron pérdidas, después de que la garantía insostenible hubiera precipitado a todo el país hacia un programa de asistencia.
}

Al igual que la tolerancia supervisora, el enfoque de «bolsillos profundos» para la resolución de crisis bancarias estaba directamente relacionado con el nacionalismo bancario. Se consideraba que los países que optaran por una disciplina de mercado más estricta y por imponer pérdidas a los acreedores de bancos en quiebra verían a «sus» bancos lastrados por mayores costes de financiación y posiblemente adquiridos por rivales extranjeros mejor respaldados (Véron, 2013). Padoa-Schioppa, por entonces ministro de Finanzas de Italia, pidió una revisión de la solidez bancaria de todo el sector a nivel europeo, lo que implicaba clasificar a las entidades más débiles y reestructurar los casos sin esperanza. Pero esto fue rechazado de manera decisiva por todos los demás dirigentes políticos (Enria, 2020b).\footnote{%
    Véase también Tommaso Padoa-Schioppa, «Europe needs a single financial rule book», Financial Times, 10 de diciembre de 2007, https://www.ft.com/content/b3c5f9c0-a750-11dc-a25a-0000779fd2ac.
}

En su lugar, se acordó, mediante un amplio consenso y por particular insistencia de Alemania, que cada país se ocuparía de «sus» bancos, incluso en el punto culminante de la dislocación de los mercados financieros a principios de octubre de 2008 (Bastasin, 2015, capítulos 1 y 2). Todavía en octubre de 2011, la posición oficial de la UE era que «los gobiernos nacionales deberían proporcionar apoyo [financiero]» a cualquier banco que no fuera capaz de mantener una posición financiera sólida mediante mecanismos de mercado.\footnote{%
    Declaración de la Cumbre del Euro de 26 de octubre de 2011, p. 15, https://www.consilium.europa.eu/uedocs/cms_data/docs/pressdata/en/ec/125644.pdf. El control de las ayudas estatales de la UE mitigó algunos de los efectos más distorsionadores de dicha asistencia financiera, pero no pretendía evitarla por completo.
}

Sin embargo, a medida que se revelaba que un número creciente de bancos de toda Europa carecía de solidez, el enfoque nacional de «bolsillos profundos» consistente en rescatar a todos los bancos problemáticos se volvió cada vez menos asequible. Ya en marzo de 2008, como mostró posteriormente una investigación del FMI, tomó forma «una crisis bancaria distintivamente europea», en la que «los diferenciales soberanos tendían a aumentar con la creciente demanda de apoyo por parte de sectores financieros nacionales debilitados, especialmente en países con menores perspectivas de crecimiento y mayores cargas de deuda» (Mody y Sandri, 2011). En los países donde los sistemas bancarios eran grandes en relación con la economía en general debido a un elevado apalancamiento y/o a la expansión en el extranjero, ese pasivo contingente amenazaba la propia sostenibilidad de las finanzas públicas, hasta un punto que no se había anticipado durante el período formativo de la unión monetaria, como se expuso en el capítulo anterior.

A su vez, la incipiente tensión financiera soberana fue respondida de manera natural mediante «persuasión moral» por parte de las autoridades públicas sobre los bancos nacionales para que facilitaran la financiación soberana, una forma de represión financiera que posteriormente se vería aún más favorecida por las operaciones de financiación a largo plazo del BCE a finales de 2011.\footnote{%
    El presidente francés Sarkozy resumió memorablemente el mecanismo: «Esto significa que cada Estado puede recurrir a sus bancos, que tendrán liquidez a su disposición» (citado por Paul Taylor, «Exclusive - ECB limits bond buying, eurozone looks to banks», Reuters, 9 de diciembre de 2011, https://www.reuters.com/article/uk-eurozone-ecb/exclusive-ecb-limits-bond-buying-eurozone-looks-to-banks-idUKTRE7B80OA20111209). Sobre el consiguiente aumento de las exposiciones soberanas nacionales, véase James (2012a).
}

En el otoño de 2010, Irlanda tuvo que afrontar austeridad y la humillante experiencia de un programa de asistencia externa debido al coste exorbitante de rescatar a sus bancos.\footnote{%
    Esto no carecía totalmente de precedentes para un país de la UE, puesto que el colapso de Parex Bank en 2008 había obligado a Letonia a solicitar asistencia al FMI. Sin embargo, el contexto era marcadamente diferente y la causalidad más multifacética, dados los desafíos macroeconómicos de Letonia. El economista Anders Aslund (2010, p. 27) resumió: «Aunque Letonia también era un accidente a punto de ocurrir, podría haber resistido durante varios meses sin un programa del FMI de no haber sido por Parex Bank».
}

En el verano de 2011, estas dinámicas de contagio ya no se limitaban a países relativamente pequeños como Grecia, Irlanda y Portugal\textsuperscript{44}. España, Italia y Francia también experimentaron diversas formas de tensión. En España, las fragilidades del sector bancario y, especialmente, de las cajas de ahorros afloraron desde finales de 2008, pero su magnitud se hizo más dramáticamente visible a finales de 2011, cuando asumía el poder un nuevo Gobierno encabezado por Mariano Rajoy (Baudino et al., 2023). La aceleración posterior del proceso de revelación condujo a las dimisiones, con dos semanas de diferencia a finales de la primavera de 2012, del gobernador del Banco de España y del subgobernador responsable de supervisión (Véron, 2016, p. 29). La correspondiente carga contingente de rescatar a los acreedores de los bancos en dificultades hizo que los inversores estuvieran cada vez más nerviosos respecto a prestar al Estado español.

En Italia, el elevado nivel de deuda pública provocó un aumento paralelo de los diferenciales de los bonos, impulsado por las preocupaciones acerca de la solvencia crediticia soberana y, a su vez, un deterioro de las condiciones de financiación de los bancos del país. En Francia, durante unas semanas a finales de agosto y principios de septiembre de 2011, los inversores comenzaron a cuestionar si la históricamente sólida garantía del Estado sobre los grandes bancos del país podría convertirse en una fuente de vulnerabilidad para la propia firma financiera del Gobierno. En consecuencia, los bancos tuvieron dificultades para acceder a financiación en dólares estadounidenses y los diferenciales soberanos franceses también comenzaron a aumentar\textsuperscript{45}. Aunque el foco del mercado sobre Francia resultó ser efímero, el episodio tuvo un profundo impacto sobre los banqueros y los responsables de política en París, lo que desempeñaría un papel clave en su abandono colectivo del nacionalismo bancario a mediados de 2012, como se detalla más adelante. Las dinámicas desordenadas se ilustran en la Figura 1, que muestra los diferenciales de la deuda a diez años de los tres países respecto de Alemania.

\imagenfit{Fig1.png}{Francia, Italia, España, diferencial de los bonos soberanos a 10 años respecto de Alemania (puntos porcentuales)}{Bloomberg}

A principios de 2009, a partir de la observación de los acontecimientos, particularmente en Islandia e Irlanda, el personal del FMI elaboró lo que parece haber sido el primer diagnóstico, en el contexto de la zona del euro, del círculo vicioso entre sectores bancarios nacionales sobredimensionados y soberanos que insisten en garantizarlos (Mody, 2009). La narrativa del círculo vicioso entre bancos y soberanos captó el núcleo del mecanismo de contagio que alimentaba la inestabilidad financiera en la zona del euro y fue adoptada de manera cada vez más amplia por académicos y otros analistas durante 2010 y 2011. A principios de 2012, su reconocimiento como motor de la crisis de la zona del euro se había convertido en una cuestión de consenso abrumador en los círculos responsables de las políticas.

En la segunda mitad de 2011, la creciente conciencia del círculo vicioso entre bancos y soberanos dio lugar, lógicamente, a propuestas procedentes tanto del interior como del exterior de las instituciones responsables de las políticas para una integración de la política fiscal, una integración de la política bancaria o ambas. En el plano conceptual, el economista jefe del FMI, Olivier Blanchard, resumió el círculo vicioso de manera nítida y gráfica en una serie de presentaciones a responsables de política durante agosto y septiembre de 2011 (Véron, 2016, p. 7). En noviembre, el influyente Consejo Alemán de Expertos Económicos convirtió la descripción del círculo vicioso entre bancos y soberanos en un elemento destacado de su muy comentado informe anual, junto con un enérgico llamamiento a Alemania para que actuara con el fin de romperlo (GCEE, 2011; Schäfer, 2017, p. 94).

El punto culminante de la defensa retórica de la integración de la política fiscal durante esa secuencia fue un discurso de la canciller alemana Angela Merkel ante el Bundestag el 2 de diciembre de 2011, en el que pidió una «unión fiscal» (\textit{Fiskalunion})\textsuperscript{46}. Sin embargo, Merkel dejó claro que su visión no implicaría la mutualización de las deudas de los Estados miembros ni la emisión conjunta, sino más bien una camisa de fuerza para las políticas nacionales que garantizara la prudencia fiscal en todos los Estados miembros. Como tal, podía considerarse una manera de reducir el riesgo de crédito soberano a largo plazo, pero a corto plazo probablemente agravaría el círculo vicioso en lugar de mitigarlo o romperlo.

En cuanto a la integración de la política bancaria, se produjo un acontecimiento clave en agosto de 2011, cuando tanto la Autoridad Bancaria Europea (ABE) como el FMI (Lagarde, 2011) comenzaron a realizar llamamientos públicos en favor de la recapitalización directa de los bancos europeos frágiles mediante un fondo común de la zona del euro, muy probablemente el Mecanismo Europeo de Estabilidad (MEDE), que en aquel momento se encontraba en proceso de creación.\footnote{%
    Market News International, «EBA Calls for Direct EFSF Bank Lending, More Capital – Press», 30 de agosto de 2011, https://www.forexlive.com/news/!/eba-calls-for-direct-efsf-bank-lending-more-capital-press-20110830. La ABE fue establecida mediante legislación de la UE el 1 de enero de 2011 como una versión algo reforzada del anterior CEBS, que había sido creado en 2004 —véase el capítulo 2—. Estaba dirigida por Andrea Enria, que había trabajado con Padoa-Schioppa en el Banco de Italia, había sido secretario general del Comité de Supervisión Bancaria del BCE entre 1999 y 2004 y del CEBS entre 2004 y 2008, y que posteriormente sería responsable de la Supervisión Bancaria del BCE entre 2019 y 2023.
}

El enfoque de recapitalización directa mediante el MEDE fue apoyado de inmediato por el Gobierno estadounidense, que lo consideraba funcionalmente equivalente a sus propias «pruebas de resistencia», que en gran medida habían logrado restablecer la confianza en el sistema bancario estadounidense en la primavera de 2009. En una reunión de ministros europeos de Finanzas celebrada a mediados de septiembre de 2011 en Breslavia, Polonia, el secretario del Tesoro estadounidense, Timothy Geithner, defendió un enfoque contundente para sanear el sector bancario de la zona del euro. El 4 de noviembre, el presidente Barack Obama exhortó de manera similar a los dirigentes de Francia, Alemania, Italia y España en una sesión paralela de la cumbre del Grupo de los Veinte celebrada en Cannes (Bastasin, 2015, pp. 328 y 334). A principios de diciembre, el presidente del Consejo Europeo, Herman Van Rompuy, economista de formación que había trabajado en el Banco Nacional de Bélgica a mediados de la década de 1970, pidió a su vez, en un informe público dirigido a los dirigentes europeos, «introducir la posibilidad de que el MEDE recapitalice directamente las instituciones bancarias» (Van Rompuy, 2011).

La expresión «unión bancaria» fue introducida por primera vez por el autor en el debate público en diciembre de 2011, quince días después del discurso de Merkel ante el Bundestag y con un énfasis inicial en el paralelismo entre su retórica de unión fiscal y lo que debía hacerse en materia de política bancaria (Véron, 2011; De Rynck, 2015, p. 9)\textsuperscript{48}. Su adopción se generalizó tras un artículo publicado en el Financial Times a principios de abril de 2012, que la utilizó en relación con una presentación durante la cual el miembro del Comité Ejecutivo del BCE Jörg Asmussen defendió la integración a nivel europeo tanto de la supervisión prudencial como de los recursos financieros para la gestión de crisis. En aquel momento, todavía no estaba claro si la unión bancaria debía limitarse a la zona del euro o si también incluiría al Reino Unido y a todo el mercado único de la UE.\footnote{%
    Alex Barker, «Eurozone weighs union on bank regulation», Financial Times, 3 de abril de 2012, https://www.ft.com/content/f03ab0bc-7d84-11e1-81a5-00144feab49a. Todavía en junio de 2012, algunos observadores contemplaban la concesión de un mandato supervisor directo a la ABE, lo que implicaba un ámbito geográfico extendido a todo el mercado único (Barker y Parker, 2012).
}

La Comisión Europea comenzó a referirse a la unión bancaria en su discurso público a finales de mayo de 2012.\footnote{%
    Véanse James G. Neuger, «EU Weighs Direct Aid to Banks as Antidote to Crisis», Bloomberg, 31 de mayo de 2012, https://www.bloomberg.com/news/articles/2012-05-30/eu-weighs-direct-aid-for-banks-common-bonds-as-crisis-antidote; y el memorando de la Comisión Europea de 6 de junio de 2012, «The banking union», https://ec.europa.eu/commission/presscorner/detail/en/MEMO_12_413.
}

Incluso después de eso, los países de la UE tardaron más de un año en acordar colectivamente el uso de la expresión y la ambición que señalaba.\footnote{
Reflexionando sobre la cuestión, Van Rompuy (2014a) señaló: «Recuerdo vívidamente cómo, apenas dos años antes, fuimos cuidadosos de no emplear el término “unión bancaria”, por la preocupación de que fuera políticamente demasiado sensible y la gente levantara barricadas».
}

Entretanto, había comenzado a formarse en la comunidad de responsables de las políticas de la UE un consenso sobre lo que la unión bancaria debía significar en la práctica. Tanto el BCE como el FMI desempeñaron un papel importante en la catalización de ese consenso.

Este último había estado sistemáticamente por delante de los acontecimientos en años anteriores al defender la integración de la política bancaria europea (por ejemplo, Decressin et al., 2007; Fonteyne et al., 2010). En noviembre de 2011, el vicepresidente del BCE, Vítor Constâncio, declaró en un discurso público: «Para la zona del euro lo diré claramente: necesitamos, para las instituciones bancarias transfronterizas, una Autoridad Europea de Resolución, que incluya o esté combinada con un Fondo de Resolución, así como un Supervisor Europeo» (Constâncio, 2011). La entonces directora gerente del FMI, Christine Lagarde, sugirió una hoja de ruta específica en enero de 2012: «Para romper el bucle de retroalimentación entre soberanos y bancos, necesitamos una mayor distribución del riesgo a través de las fronteras en el sistema bancario. A corto plazo, un mecanismo para toda la zona del euro que tenga capacidad para adquirir participaciones directas en bancos [es decir, la recapitalización directa mediante el MEDE] ayudará a romper este vínculo. Mirando más adelante, la unión monetaria necesita estar respaldada por una integración financiera en forma de supervisión unificada, una autoridad única de resolución bancaria con un respaldo común y un fondo único de seguro de depósitos» (Lagarde, 2012).

El presidente del BCE, Mario Draghi, se hizo eco a su vez de esta visión en abril, aunque con una formulación menos precisa: «Garantizar una UEM [Unión Económica y Monetaria] que funcione adecuadamente implica reforzar la supervisión y la resolución bancarias a nivel europeo» (Draghi, 2012). En la reunión informal del Consejo Europeo de 23 de mayo de 2012, los dirigentes nacionales pidieron a Van Rompuy que elaborara en unas pocas semanas un informe sobre cómo detener el rápido deterioro de las condiciones financieras en la zona del euro. En esa ocasión, el presidente francés François Hollande, recién elegido a principios de ese mismo mes y alentado por el primer ministro italiano Mario Monti, se pronunció ante los periodistas a favor de la unión bancaria, algo que su predecesor Nicolas Sarkozy nunca había hecho (Bastasin, 2015, p. 374).

$$\textbf{2.2. Una semana decisiva a finales de junio de 2012}$$

El impulso político hacia la unión bancaria, que se había acumulado gradualmente durante la segunda mitad de 2011 y la primera mitad de 2012, llegó súbitamente a un punto de decisión durante una secuencia de tres reuniones clave a finales de junio de ese año, coincidiendo con lo que uno de los principales protagonistas ha descrito como el momento más peligroso de toda la crisis de la zona del euro (Rehn, 2020, p. 172)\textsuperscript{52}. El viernes 22 de junio, en Roma, Mario Monti recibió a Hollande, Merkel y Van Rompuy, como preparación para la reunión del Consejo Europeo prevista para la semana siguiente; Rajoy, de España, se unió a ellos por iniciativa propia. Durante este encuentro, los dirigentes no entraron en aspectos técnicos, pero acordaron que debían acelerarse los esfuerzos para encontrar una solución a los problemas bancarios que alimentaban el círculo vicioso entre bancos y soberanos, y que, con ese fin, sus respectivos ministros de Finanzas debían reunirse de manera urgente y secreta.

La reunión de los ministros se programó para el martes siguiente por la noche en el aeropuerto Charles de Gaulle (CDG), cerca de París, una elección de lugar inusual que combinaba las ventajas de la accesibilidad y la discreción. El lunes 25 de junio, España solicitó formalmente un programa de asistencia, incluida ayuda del Mecanismo Europeo de Estabilidad, que estaba a punto de ser establecido, para abordar las crecientes preocupaciones sobre su sector bancario.

El martes 26 de junio, en Bruselas, Van Rompuy publicó el informe que los dirigentes le habían pedido que preparara el 23 de mayo. Puesto que lo había preparado en coordinación con sus homólogos de la Comisión Europea —José Manuel Barroso—, el Eurogrupo —Jean-Claude Juncker—\textsuperscript{53} y el BCE —Mario Draghi—, el texto pasó a conocerse como el «Informe de los Cuatro Presidentes» (Van Rompuy, 2012)\textsuperscript{54}. En él, Van Rompuy presentó un plan para la reforma de la zona del euro, proporcionando un importante punto de referencia para los debates posteriores, mediante una combinación de cuatro esfuerzos interrelacionados: unión bancaria —«un marco financiero integrado [que] eleva la responsabilidad de la supervisión al nivel europeo y establece mecanismos comunes para resolver bancos y garantizar los depósitos de los clientes»—; unión fiscal —«un marco presupuestario integrado [que implica] pasos proporcionales hacia la emisión común de deuda [y] diferentes formas de solidaridad fiscal»—; unión económica —«un marco integrado de política económica [para] promover el crecimiento sostenible, el empleo y la competitividad»—; y unión política —«garantizar la necesaria legitimidad democrática y la rendición de cuentas de la toma de decisiones dentro de la UEM, sobre la base del ejercicio conjunto de la soberanía para las políticas comunes y la solidaridad»—. La colocación de la unión bancaria en primera posición sugería que sería el principal foco de actuación en el futuro inmediato.\footnote{%
    En parte como respuesta a la creciente tensión del sector bancario en España y a los llamamientos de varios Estados miembros para complementar la supervisión bancaria nacional con un mecanismo de vigilancia más fiable, parece que el BCE realizó en mayo-junio de 2012 un análisis preliminar interno sobre la viabilidad de asumir un mandato de supervisión bancaria sobre la base del artículo 127(6) del TFUE (De Rynck, 2015, p. 11). Este análisis incluía un esbozo aproximado de cómo las tareas correspondientes podrían separarse operativamente de la ejecución de la política monetaria, como posteriormente quedó consagrado en el Reglamento del MUS.
}

Horas después de la publicación del informe de Van Rompuy, los ministros de Finanzas de los cuatro países y algunos funcionarios adicionales se reunieron según lo previsto en una anodina sala de conferencias del hotel Sheraton del aeropuerto CDG. Eran los ministros Wolfgang Schäuble por Alemania, Pierre Moscovici por Francia, Vittorio Grilli por Italia y Luis de Guindos por España; los altos funcionarios ministeriales Thomas Steffen —Alemania— y Ramon Fernandez —Francia—; el jefe adjunto de gabinete de Hollande encargado de asuntos económicos, Emmanuel Macron; junto con el comisario europeo Olli Rehn, el jefe de gabinete de Van Rompuy, Frans van Daele, el presidente del Grupo de Trabajo del Euro, Thomas Wieser, el director general de Asuntos Económicos de la Comisión, Marco Buti, el asesor de Rehn, Taneli Lahti, y el director del Departamento Europeo del FMI, Reza Moghadam. Durante aquella reunión, cuya propia existencia permaneció secreta hasta que un artículo de prensa la reveló dieciocho meses después, se llevó a cabo el esencial intercambio político de concesiones que posteriormente definiría la unión bancaria, no sin nuevos giros ni ironías.\footnote{%
    Peter Spiegel y Alex Barker, «Banking union falls short of EU goal», Financial Times, 19 de diciembre de 2013, https://www.ft.com/content/f1c23942-68cd-11e3-bb3e-00144feabdc0.
}

Hollande, aparentemente impulsado por Monti, había expresado su apoyo a la unión bancaria en la reunión de Bruselas del 23 de mayo. Esto era nuevo. El Gobierno francés había defendido sistemáticamente la recapitalización directa de los bancos por parte del MEDE, haciéndose eco de las recomendaciones de la ABE y el FMI de agosto de 2011. En cambio, la puesta en común a nivel europeo de la supervisión prudencial nunca había recibido anteriormente el apoyo francés. A principios de la década de 1990, como se describió en el capítulo anterior, el Tesoro francés había figurado entre quienes se oponían a la inclusión en el Tratado de Maastricht de una función supervisora para el BCE. A principios de 2012, la supervisión bancaria seguía siendo considerada por la mayoría de los altos funcionarios —y banqueros— franceses como un instrumento de soberanía nacional que debía mantenerse seguro en París, aunque el propio Tesoro había comenzado a considerar opiniones alternativas, como se detalla más adelante. Varias semanas después de la reunión de mayo, los medios seguían situando a los supervisores bancarios franceses, junto con sus homólogos neerlandeses y alemanes, entre las «autoridades nacionales muy atrincheradas que no tienen intención alguna de renunciar al poder» (Barker y Parker, 2012).

En ese contexto, cuando Schäuble realizó una maniobra inicial en CDG consistente en aceptar el principio de la recapitalización directa por parte del MEDE de los bancos españoles —el foco inmediato de preocupación en aquel momento— con la condición de que existiera una auténtica supervisión bancaria europea, fue a la vez un avance decisivo y una postura de línea dura. Fue un avance decisivo porque Alemania había rechazado hasta entonces de plano cualquier recapitalización directa por parte del MEDE. Al mismo tiempo, Schäuble adoptaba una postura de línea dura porque tenía buenas razones para creer que su condición no sería aceptable para sus homólogos franceses y que la supervisión prudencial a nivel europeo sería considerada una exigencia alemana a la que era necesario oponerse.\footnote{
Schäfer (2017) señaló que Hollande, en sus declaraciones posteriores al Consejo Europeo informal de mayo de 2012, se había pronunciado firmemente a favor de la integración supervisora (p. 117): «Yo mismo dije que quiero que los mecanismos de supervisión financiera, las garantías de depósitos y la resolución de crisis estén integrados. […] Cuanto más se coordine y centralice, mejor será la respuesta en materia de supervisión, resolución de crisis y, sobre todo, garantías de depósitos». Pero, basándose en sus entrevistas, Schäfer también indicó la opinión existente entre los funcionarios franceses en el momento de la reunión de CDG de que la propuesta de Schäuble de un supervisor conjunto «fue planteada como una condición que nunca aceptaríamos» (p. 121).
}

De hecho, las opiniones en Alemania sobre la integración europea de la supervisión bancaria distaban mucho de ser unánimes. Había partidarios de larga data, incluidos los bancos comerciales encabezados por Deutsche Bank, que la habían defendido durante muchos años. En cambio, los influyentes bancos públicos y cooperativos eran hostiles a cualquier puesta en común de la política bancaria a nivel europeo. Temían que inevitablemente condujera a un trato menos favorable de sus acuerdos idiosincrásicos. El Bundesbank se había mostrado durante mucho tiempo hostil a que el BCE desempeñara una función en la supervisión de bancos individuales, sobre la base de que ello perjudicaría la integridad de la política monetaria,\footnote{
Según un testigo cercano, «Alemania se había opuesto firmemente a ella [la supervisión bancaria común] incluso cuando la crisis financiera se encontraba en su punto culminante, en 2009» (van Middelaar, 2019, p. 56). En una entrevista, Lamfalussy recordó que, durante las primeras fases de la gran crisis financiera, el BCE no tenía conocimiento de la verdadera situación del sector bancario de la zona del euro «porque los alemanes se oponían a ello. Pensaban que, al asumir esa nueva tarea [de vigilancia del sector bancario, aunque solo fuera en el nivel macroprudencial], el BCE corría el riesgo de pervertir su mandato básico, que es vigilar la moneda» (Lamfalussy, 2013, p. 173; traducción del autor).
} y continuó argumentando que sería peligroso embarcarse en una unión bancaria sin una unión fiscal simultánea.\footnote{
James Wilson, «Bundesbank warns on EU banking union», Financial Times, 12 de junio de 2012, https://www.ft.com/content/79c17794-b467-11e1-bb68-00144feabdc0.
}

Sin embargo, inesperadamente, Macron, hablando en nombre de Hollande, aceptó de inmediato el quid pro quo y orientó la discusión hacia cómo aplicar en la práctica la idea de supervisión integrada de Schäuble. Por razones de conveniencia, defendió el uso de la cláusula habilitante del Tratado de Maastricht como base jurídica para la supervisión europea, lo que también había sido recomendado ese mismo día en el Informe de los Cuatro Presidentes de Van Rompuy.\footnote{
Van Rompuy (2012), p. 4. La activación de la cláusula habilitante del artículo 127(6) del TFUE no carecía técnicamente de precedentes, puesto que ya se había utilizado como base para un reglamento adoptado en 2010 —Reglamento (UE) 1096/2010— con el fin de permitir que el BCE apoyara al recién creado Comité Europeo de Riesgo Sistémico. Sin embargo, ese texto tenía un impacto y un carácter controvertido incomparablemente menores que el establecimiento de una supervisión bancaria europea integrada. Schäfer (2017, p. 121) corrobora la opinión de que las «principales discusiones [en CDG] tuvieron lugar entre Schäuble y Macron».
}

Mientras tanto, todos los participantes tomaron una decisión consciente de no incluir el seguro de depósitos en el orden del día, aunque había sido mencionado explícitamente en el informe de Van Rompuy.\footnote{
A lo largo de este texto, las expresiones «seguro de depósitos» y «garantía de depósitos» se utilizan como sinónimos. La legislación de la UE se refiere a los mecanismos nacionales como sistemas de garantía de depósitos, mientras que «seguro de depósitos» se utiliza generalmente cuando se consideran mecanismos similares a nivel europeo.
} Esta elección también estuvo motivada por razones de conveniencia, porque eran conscientes de la sensibilidad política del seguro de depósitos, particularmente en Alemania, y por la mayor urgencia percibida de la recapitalización directa de los bancos españoles por parte del MEDE.\footnote{
Irónicamente, un subconjunto de países de la UE, incluidos Francia y Alemania, había respaldado formalmente el principio de un seguro europeo de depósitos apenas unos días antes, en la declaración de los dirigentes del Grupo de los Veinte (G20) en Los Cabos, México, el 19 de junio de 2012: «Los miembros de la zona del euro del G20 adoptarán todas las medidas necesarias para salvaguardar la integridad y la estabilidad de la zona, mejorar el funcionamiento de los mercados financieros y romper el bucle de retroalimentación entre soberanos y bancos. […] Con ese fin, apoyamos la intención de considerar medidas concretas hacia una arquitectura financiera más integrada, que abarque la supervisión bancaria, la resolución y recapitalización, y el seguro de depósitos». Sin embargo, posteriormente las autoridades alemanas nunca parecieron sentirse vinculadas por ese compromiso.
}

La aceptación francesa súbita y sin reservas de la supervisión bancaria europea debe entenderse en el contexto de un entorno que cambiaba rápidamente en Francia. El episodio de turbulencias de mercado de agosto-septiembre de 2011 había hecho añicos certezas mantenidas durante mucho tiempo en París sobre la optimalidad del nexo francés entre bancos y Estado, en el que el Estado rescataría a cualquier banco en quiebra —como había hecho con Crédit Lyonnais— y los bancos aceptarían la orientación estatal de sus préstamos en cuestiones de interés nacional. Dentro del Tesoro francés, a principios de 2012 hubo intensos debates entre quienes defendían la simbiosis tradicional con la comunidad bancaria nacional y quienes consideraban que defender el crédito del Estado francés podría exigir establecer una nueva distancia entre este y los bancos. El director del Tesoro en aquel momento, Ramon Fernandez, estaba convencido de la necesidad de la recapitalización directa por parte del MEDE y dispuesto a aceptar cambios en la postura tradicional francesa sobre soberanía supervisora para conseguirla.

El ciclo electoral de apenas unas semanas antes, con las elecciones presidenciales ganadas por Hollande a principios de mayo de 2012 y las elecciones parlamentarias ganadas por su partido de centroizquierda en junio, había llevado al poder a un nuevo equipo político cuyos principales miembros tenían comparativamente poco bagaje en materia de política bancaria, al haber estado ausentes de los múltiples episodios de gestión de crisis de los cinco años anteriores.\footnote{%
    De Rynck (2017, p. 129) se refirió a la aceptación francesa de la supervisión bancaria europea en 2017 como un «giro de 180 grados respecto de su preferencia de 2010 por mantener la supervisión en el nivel nacional». Van Middelaar (2019, p. 55) sugirió que la transición de Sarkozy a Hollande permitió a Van Rompuy fomentar una mentalidad de «ningún tabú» al incluir los desafíos estructurales de la arquitectura financiera de la zona del euro en el orden del día de la cumbre del 23 de mayo. Puede señalarse, sin embargo, que la política del sector bancario no fue un tema de división partidista durante las campañas electorales de la primavera.
}

El avance decisivo en el aeropuerto CDG preparó el terreno para el siguiente y más público gran paso: la reunión del Consejo Europeo y la cumbre de la zona del euro de los días 28 y 29 de junio en Bruselas. Aunque gran parte de la cobertura de esa reunión se ha centrado en una maniobra combativa de Monti para imponer su visión de apoyo europeo a la deuda soberana italiana mediante compras de bonos por parte del MEDE, ese episodio resultó ser considerablemente menos trascendente que las dinámicas de toma de decisiones sobre la unión bancaria que tuvieron lugar en paralelo —y que Monti también apoyaba—.

Al mismo tiempo que los jefes de Estado y de Gobierno se reunían en el edificio Justus Lipsius de Bruselas durante la noche del jueves 28 de junio, sus asesores —los «sherpas»— y altos funcionarios de los ministerios de Finanzas se reunieron en otra sala del mismo edificio, donde trabajaron en varios borradores sucesivos sobre la combinación de la supervisión bancaria europea y la recapitalización bancaria directa por parte del MEDE que los participantes en la reunión de CDG habían acordado dos días antes. La cuestión crítica era la articulación secuencial entre los dos componentes del quid pro quo. Se reconocía que la creación de un sistema viable de supervisión bancaria europea, basado en el artículo 127(6) del TFUE y, por consiguiente, centrado en el BCE, constituía un esfuerzo a medio plazo, que requería al menos varios meses de proceso legislativo y quizá un año para la posterior preparación operativa. En las intensas condiciones de crisis del momento, eso parecía un plazo muy largo, con múltiples incertidumbres procedimentales a lo largo del camino y ninguna garantía de eficacia final.

En cambio, la recapitalización directa por parte del MEDE de los bancos necesitados estaría llena de dificultades, pero, si se ejecutaba bien, podría lograrse comparativamente con gran rapidez, y la señal de puesta en común del riesgo era considerada por la mayoría de los defensores de la medida como una cuestión de urgencia inmediata.\footnote{%
    La estructura precisa de las operaciones de recapitalización directa nunca se hizo explícita públicamente, aunque es probable que se detallara en borradores compartidos entre los negociadores. Un funcionario del MEDE la describió mucho más tarde como «realizar una gran inversión de capital en la parte viable de un banco sistémico en quiebra» (MEDE, 2019, p. 294).
}

Un borrador inicial, preparado por la secretaría del Grupo de Trabajo del Euro dirigida por Wieser junto con la Comisión Europea, intentó cuadrar el círculo haciendo referencia a que la recapitalización directa por parte del MEDE tendría lugar después de que se hubiera establecido el mecanismo supervisor centrado en el BCE, pero añadía la frase: «Se establecerán soluciones provisionales para la recapitalización directa de las entidades [de crédito] por parte del MEDE durante el tiempo en que se esté creando este mecanismo». La reunión de los sherpas concluyó alrededor de la medianoche con una versión revisada del proyecto de declaración que incluía una versión ligeramente modificada de esa frase, preservando el concepto de «soluciones provisionales».

Pero a ello se opusieron, durante las primeras horas del viernes 29 de junio, tres de los dirigentes políticos reunidos: Jyrki Katainen, de Finlandia; Mark Rutte, de los Países Bajos; y Merkel. Draghi no defendió las «soluciones provisionales», sino que sugirió la posibilidad de una aplicación retroactiva mediante la cual el MEDE, en un futuro a medio plazo tras el inicio de la supervisión bancaria europea, asumiría los instrumentos de capital que inicialmente hubieran sido proporcionados por los países individuales, asumiendo así finalmente el riesgo correspondiente de manera común. Para Draghi, el beneficio crítico del paquete parece haber sido la firme afirmación de la confianza de los dirigentes en el BCE, expresada mediante la concesión de autoridad supervisora sobre los bancos, que presumiblemente consideró más importante que la puesta en común inmediata del riesgo a través del MEDE.\footnote{%
    Una consideración adicional para el BCE era que su extensión de liquidez a los bancos, enormemente ampliada y a más largo plazo, mediante las operaciones de financiación a largo plazo que Draghi había anunciado a finales de 2011, lo exponía a un riesgo insostenible a menos que obtuviera una mayor comprensión de las condiciones financieras de las entidades receptoras (De Rynck, 2015, p. 11).
}

Merkel también apoyó firmemente el concepto de supervisor único durante la reunión e insistió en añadir un plazo de finales de 2012 para su elaboración legislativa (van Middelaar, 2019, p. 57). La declaración final, adoptada en medio del cansancio general a las 4:35 de la madrugada del 29 de junio (Rehn, 2020, p. 188) y publicada inmediatamente después, mantuvo la referencia a la recapitalización directa por parte del MEDE, pero sin mención alguna a «soluciones provisionales». Tras las múltiples rondas de redacción, el texto clave era enrevesado: «La Comisión presentará en breve propuestas sobre la base del artículo 127(6) para un mecanismo único de supervisión. Pedimos al Consejo que examine estas propuestas con carácter urgente antes de finales de 2012. Cuando se establezca un mecanismo único de supervisión efectivo, que implique al BCE, para los bancos de la zona del euro, el MEDE podría, tras una decisión ordinaria, tener la posibilidad de recapitalizar directamente a los bancos».

Sin embargo, esta tortuosa formulación iba precedida de una declaración de propósito inusualmente directa y contundente: «Afirmamos que es imperativo romper el círculo vicioso entre bancos y soberanos»\textsuperscript{66}.

En suma, la cumbre de finales de junio proporcionó una base sólida para la futura supervisión bancaria del BCE, aunque con numerosos riesgos e incertidumbres sobre su aplicación final. En materia de gestión de crisis, respaldó el principio de recapitalización directa por parte del MEDE, pero aplazó su aplicación indefinidamente. No asumió compromiso alguno sobre un seguro común de depósitos. La principal desviación respecto del equilibrio de políticas que los participantes en la reunión de CDG habían esbozado se refería a la recapitalización directa por parte del MEDE. Resulta difícil separar las causas de esa diferencia de la compleja relación entre Merkel y Schäuble, que pudo haber implicado enfoques distintos sobre la política y la gestión de crisis, así como una rivalidad política de larga data. También resulta difícil separar las causas de la dinámica de la coalición alemana de aquel momento, que incluía a un socio —el Partido Democrático Libre— aún más reticente que Schäuble respecto de la distribución del riesgo en toda la zona del euro.

A primera vista, la preferencia de Merkel parece haber sido el rechazo completo de la recapitalización directa por parte del MEDE, en consonancia con la postura alemana anterior, mientras que Schäuble se sentía vinculado por su compromiso en CDG.\footnote{%
    La versión menos obstructiva de esa postura, formulada también por el Bundesbank en junio de 2012, consistía en vincular la unión bancaria a la condición de una unión fiscal simultánea. Barker y Parker (2012) escribieron así el 18 de junio de 2012: «Antes de exponer a los contribuyentes alemanes a pasivos extranjeros —como el seguro de depósitos o participaciones directas en bancos—, Angela Merkel, canciller alemana, quiere controles federales sobre los bancos nacionales y una unión fiscal. En otras palabras, algún control compartido de la tributación y el gasto nacionales».
}

La formulación de la declaración del 29 de junio, sin la referencia a soluciones provisionales, se situaba en algún punto intermedio entre estas dos posiciones, dejando a ninguno de los dos alemanes plenamente satisfecho. Como se detallará más adelante, la visión restrictiva de Merkel sobre el uso del MEDE terminó prevaleciendo a través de los acontecimientos posteriores, pero ello fue parcialmente compensado por la creación de un mecanismo único de resolución que incorporaba cierto grado de distribución del riesgo, que a su vez podría ser finalmente «respaldado» por el MEDE. La otra parte del acuerdo, relativa a la puesta en común de la autoridad de supervisión bancaria, resultó ser extremadamente resistente.

$$\textbf{2.3. La unión bancaria y la resolución de la crisis}$$

Las semanas que siguieron a la declaración del 29 de junio fueron caóticas, pero decisivas. En algún momento de principios de julio, Merkel parece haber tomado la firme decisión de que una salida de Grecia de la zona del euro era demasiado arriesgada como para intentarla. Entretanto, en las negociaciones sobre la aplicación de las decisiones adoptadas a finales de junio, los funcionarios alemanes dieron un paso más atrás en el compromiso de utilizar el MEDE para recapitalizaciones bancarias directas en España, incluso en un futuro relativamente lejano. Parte de la motivación de ese comportamiento inusual pudo haber estado relacionada con la preocupación de que pudiera debilitar la posición que defendían simultáneamente ante el Tribunal Constitucional alemán de que el MEDE era compatible con la Ley Fundamental de Alemania.\footnote{
En la práctica de la UE, las decisiones de las cumbres de dirigentes, como la de los días 28 y 29 de junio de 2012, suelen respetarse como vinculantes para todas las partes.
}

El tribunal celebró una tensa audiencia el 10 de julio en Karlsruhe antes de aprobar finalmente la estructura del MEDE en septiembre, aunque con determinadas salvedades (MEDE, 2019, capítulo 26). A medida que se filtraban noticias sobre la marcha atrás alemana, los mercados se mostraban cada vez más nerviosos.\textsuperscript{69} Finalmente, el 26 de julio, ante una audiencia de participantes en los mercados financieros en Londres, Draghi pronunció las memorables palabras: «Dentro de nuestro mandato, el BCE está dispuesto a hacer todo lo que sea necesario para preservar el euro; y créanme, será suficiente».\textsuperscript{70} Como ilustra la Figura 1, ello marcó un importante punto de inflexión, tras el cual la crisis de los mercados se desescaló rápidamente, incluso antes de que el BCE diera seguimiento formalmente a las palabras de Draghi anunciando su programa de Operaciones Monetarias de Compraventa (Outright Monetary Transactions, OMT) el 6 de septiembre.

En varios discursos posteriores, Draghi insinuó la existencia de un vínculo causal entre la muestra de apoyo de los dirigentes políticos al BCE el 29 de junio y su histórica declaración de intenciones en Londres (por ejemplo, Draghi, 2013). Sin embargo, nunca hizo plenamente explícita esa relación, ya que ello habría supuesto menoscabar el principio sacrosanto de independencia de la política monetaria del BCE respecto de cualquier consideración política. Van Rompuy podía permitirse ser más directo, al menos una vez transcurrido cierto tiempo. Dos años más tarde, dedicó un discurso a esta cuestión en el BCE de Fráncfort, quince días después del inicio efectivo de su mandato de supervisión prudencial, el 4 de noviembre de 2014, y quince días antes del final de su propio mandato como presidente del Consejo Europeo.

En su intervención, Van Rompuy destacó la importancia de la cumbre de finales de junio, cuyas dinámicas esenciales describió en los siguientes términos: «La unión bancaria constituye el mayor salto adelante desde la creación del euro. […] La cumbre de junio de 2012 fue quizá el Consejo Europeo más importante de mis cinco años en el cargo. […] Los dirigentes percibieron que el momento exigía un avance cualitativo. También comprendieron que debía producirse en el ámbito de la unión bancaria, el asunto más urgente de todos. […] Algunos querían comenzar por la supervisión bancaria, para evitar nuevos problemas; otros preferían actuar primero en favor de una mayor solidaridad en el ámbito bancario, para superar las dificultades heredadas del pasado. […] Vinculamos dos decisiones políticas: la creación de un mecanismo único de supervisión para todos los bancos de la zona del euro y la posibilidad de que los bancos en dificultades recibieran capital directamente de un fondo común de rescate» (es decir, el MEDE).

Van Rompuy fue más lejos de lo que Draghi jamás haría en público: «Nunca olvidaré que, un par de horas más tarde, aquel viernes, Mario Draghi entró en mi despacho […] Un hombre sometido a una enorme presión; por primera vez en los ocho meses durante los cuales le había visto trabajar, ahora parecía aliviado. “Herman”, me dijo, “¿te das cuenta de lo que habéis hecho todos anoche? Esto es el cambio de juego que necesitábamos”. El compromiso de los dirigentes políticos con la supervisión bancaria europea creó la oportunidad que su propia institución necesitaba para intensificar su papel en la crisis, con palabras —ya famosas palabras [en Londres el 26 de julio]— y con acciones, las OMT (Outright Monetary Transactions), que ambas llegaron aquel verano. Fue un punto de inflexión». Van Rompuy añadió de manera significativa: «Para un banco central, ser independiente no significa estar desconectado. Por eso los presidentes del BCE asisten a las reuniones del Consejo Europeo, de las Cumbres del Euro, del Eurogrupo y del Parlamento Europeo» (Van Rompuy, 2014b).\footnote{%
    La afirmación de Van Rompuy de que el BCE debía ser «independiente, pero no desconectado» pudo haber constituido un discreto homenaje al fallecido Padoa-Schioppa, quien subrayaba con frecuencia que la independencia del BCE «no debía significar soledad institucional» (por ejemplo, Bini Smaghi, 2011).
}

Esta interpretación ha sido corroborada por otros participantes clave en aquellos acontecimientos. El antiguo miembro del Comité Ejecutivo del BCE Peter Praet comentó en una entrevista de 2019: «El pánico de los mercados en 2012 solo pudo ser detenido por Mario Draghi. Naturalmente, el trasfondo del éxito de su frase “haré todo lo que sea necesario” fue la reunión del Consejo Europeo de jefes de Estado y de Gobierno de junio, sobre la puesta en marcha de la unión bancaria y de los mecanismos de gestión de crisis. Ese fue, por tanto, el contexto político».\footnote{%
    Rebecca Christie, «Convincing the markets: Peter Praet revisits the ECB’s unconventional monetary policy response to the euro crisis», \textit{Finance & Development}, septiembre de 2019, https://www.imf.org/en/Publications/fandd/issues/2019/09/peter-praet-on-ECB-and-euro-crisis-trenches.
}

Varios participantes franceses han descrito igualmente una relación directa entre la cumbre de la zona del euro y la actuación del BCE.\footnote{
En una entrevista a finales de 2022, el antiguo director del Tesoro francés, Ramon Fernandez, recordaba: «Durante todos esos años existió una especie de \textit{quid pro quo} entre los gobiernos y el BCE. Los gobiernos necesitaban ver al BCE desempeñar su papel, y el BCE no podía actuar sin una actuación muy firme de los gobiernos. […] El BCE es totalmente independiente, pero podía actuar cuando los gobiernos actuaban. Y, en los grandes momentos, eso fue exactamente lo que ocurrió. Por ejemplo, en 2012, cuando los gobiernos avanzan hacia la unión bancaria, el BCE / Mario Draghi puede pronunciar su discurso del “haré todo lo que sea necesario” a finales de julio». Véase Tim Gwynn Jones, «In the Room» podcast, diciembre de 2022, https://uk-podcasts.co.uk/podcast/in-the-room-1/ramon-fernandez. Hollande expuso de forma igualmente sucinta esta interpretación en sus memorias: «Aquella noche [28-29 de junio], nuestro compromiso del eje franco-alemán salvó la moneda única, derrotó a los especuladores de los mercados y, sobre todo, dio al presidente del BCE margen para una política más acomodaticia. Un mes después, en pleno verano, el 26 de julio de 2012, Mario Draghi declaró que “haría todo lo que fuera necesario para preservar el euro”. Angela Merkel y yo decidimos publicar un comunicado conjunto al día siguiente. Ese día marcó el principio del fin de la crisis de la zona del euro. La cumbre europea del 28 de junio de 2012 fue histórica. ¿Se ha reconocido suficientemente?». (Hollande, 2018, capítulo 6; traducción del autor).
}

Otros participantes siguen negando que existiera tal relación. Naturalmente, resulta imposible determinar con completa certeza si el BCE habría intervenido del modo en que lo hizo de no haberse adoptado previamente la decisión de establecer una supervisión bancaria europea. A comienzos de septiembre, muchos participantes en los mercados ya habían concluido que la fase existencial de la crisis de la zona del euro había terminado. Irónicamente, ello redujo a su vez parte de la presión sobre los responsables políticos para cumplir los compromisos asumidos y facilitó el posterior vaciamiento de la decisión del 29 de junio sobre la recapitalización bancaria directa mediante el MEDE. Mientras tanto, la decisión de los dirigentes relativa a la supervisión bancaria del BCE se aplicó fielmente.

La manera en que los dirigentes de la zona del euro lograron adoptar una decisión tan radical ha sido objeto de análisis tentativos por parte de varios autores (por ejemplo, Epstein y Rhodes, 2014; De Rynck, 2015; Glöckler, Lindner y Salines, 2017; Schäfer, 2017; Nielsen y Smeets, 2018). Del relato aquí presentado emergen cinco elementos fundamentales.

En primer lugar, el sistema existente basado en la supervisión bancaria nacional había fracasado de forma tan generalizada que ya no podía defenderse.

En segundo lugar, el imperativo de evitar la ruptura de la zona del euro, considerada razonablemente por todos los implicados como una opción inaceptable, concentró la atención de los responsables políticos y les obligó a establecer prioridades claras. Resulta llamativo que, en el momento de la decisión, las posiciones tradicionales de política de los respectivos sectores bancarios nacionales parecieran influir mucho menos en los cálculos de los principales responsables políticos que en circunstancias normales.

En tercer lugar, el papel central del círculo vicioso entre bancos y soberanos como motor del contagio y de la desestabilización fue aceptado como un hecho, proporcionando una sólida base para alcanzar un consenso sobre las decisiones de política.

En cuarto lugar, se produjo una alineación fortuita de circunstancias, en la que coincidieron oportunidad y urgencia como consecuencia, entre otros factores, de los ciclos electorales en Francia y Grecia —con una sensación de frágil respiro en este último país, tras dos elecciones parlamentarias sucesivas celebradas el 6 de mayo y el 17 de junio— y de la percepción de una crisis que se agravaba rápidamente en el sector bancario español.

En quinto lugar, varias personas demostraron cualidades de liderazgo extraordinarias: funcionarios de la UE que impulsaron las decisiones y dirigentes nacionales que aceptaron compromisos políticamente costosos. Tanto Francia como Alemania realizaron concesiones significativas, seguidas por otros Estados miembros. Francia renunció al principio de soberanía nacional sobre la supervisión bancaria. Alemania terminó retrocediendo respecto de la recapitalización directa mediante el MEDE, pero aceptó de manera decisiva que todos sus bancos —incluidos los reacios bancos públicos y cooperativos, sobre los que se tratará con mayor detalle en el capítulo siguiente— quedaran cubiertos por el nuevo mecanismo de supervisión.

Como señalaron Nielsen y Smeets (2018, p. 1251), «no había nada inevitable en la unión bancaria». Fue necesario un notable ejercicio de liderazgo político para comprender, en medio de una negociación particularmente intensa, que el acuerdo alcanzado en las primeras horas del 29 de junio de 2012 producía un saldo positivo para todos los países implicados.





\end{document}